% concatenated from completecohom.tex
\documentclass[12pt,a4paper]{article}

%% Input encoding 'utf8'
\usepackage[utf8]{inputenc}
%% Babel provides support for British English
\usepackage[english]{babel}

%% The AMS-LaTeX extensions for mathematical typesetting
\usepackage{amsmath,amssymb,amsfonts,mathrsfs}

%This will allow us to typeset theorems like examples, proofs and
%similar in a swift manner
\usepackage[amsmath,thmmarks]{ntheorem}

% Theorem-like environments

\numberwithin{equation}{section}

\newtheorem{thm}{Theorem}[section]
\newtheorem{cor}[thm]{Corollary}
\newtheorem{lem}[thm]{Lemma}
\newtheorem{prop}[thm]{Proposition}
\newtheorem{rem}[thm]{Remark}
\newtheorem{defn}[thm]{Definition}
\newtheorem{axm}[thm]{Axiom}
\newtheorem{exl}[thm]{Example}
\newtheorem{ques}[thm]{Question}
\newtheorem{notn}[thm]{Notation}
\newtheorem{conv}[thm]{Convention}

% Proof environment with a small square as a "qed" symbol
\theoremstyle{nonumberplain}
\theorembodyfont{\normalfont}
\theoremsymbol{\ensuremath{\square}}
\newtheorem{proof}{Proof}

% This is required for a reference
\usepackage{hyperref}

\title{On a Completion of Cohomological Functors Generalising Tate Cohomology I}
\date{1 April 2026}
\author{Max Gheorghiu}

% This formatting tool provides a nice margin
\usepackage[margin=2.5cm]{geometry}

% This is needed for all commuting diagrams
\usepackage{tikz-cd}
% This is needed to write double brackets for free profinite modules
\usepackage{stmaryrd}
% For attaching pictures
\usepackage{graphicx}

% This provides the symbol for an infinite free product of groups
\newcommand{\bigast}{\mathop{\scalebox{2}{\raisebox{-0.1ex}{$\ast$}}}}%
% This provides the symbol for a proper normal subgroup
\newcommand{\triangleleftneq}{%
  \mathrel{\ooalign{$\lneq$\cr\raise.22ex\hbox{$\lhd$}\cr}}}

%end the preamble and start the document
\begin{document}

%avoiding indentation at the beginning of each paragraph
\setlength{\parindent}{0cm}

%create a title page from the preamble info
\maketitle

\begin{abstract}
Tate cohomology has been generalised by several authors using different constructions that have applications in group theory, ring theory and homotopical algebra. Therefore, there is a need for a uniform account that explains why their underlying approaches all lead to the same conclusions. The key notion in such a uniform theory is a specific completion of cohomological functors that is constructed under mild assumptions. This completion takes Tate cohomology to settings where it has never been introduced such as in condensed mathematics. Through the latter, one can define Tate cohomology for any $T1$ topological group.
\end{abstract}

\tableofcontents

\section{Introduction}\label{sec:intro}

Originally, Tate cohomology was developed only for finite groups by J.\! Tate in the paper~\cite{tat52} from 1952 motivated by class field theory. As it unites group homology and group cohomology of finite groups in a convenient manner, it attracted the attention of group theorists~\cite[p.~128]{bro82}. It was generalised to groups of finite virtual cohomological dimension in F.\! T.\! Farrell's paper~\cite{far77} from 1977 and then to all groups in D.\! J.\! Benson and J.\! F.\! Carlson's paper~\cite{ben92} from 1992, F.\! Goichot's paper~\cite{goi92} from 1992 and G.\! Mislin's paper~\cite{mis94} from 1994. The different approaches of the above papers have been used in generalised constructions by several authors with applications to group theory, ring theory and homotopical algebra. There is a need for a uniform account that explains why the approaches underlying these constructions all lead to the same conclusions: a theory that does not only work for discrete groups or modules over a ring, but also in greater generality. The main goal of this paper to provide a uniform treatment of the entire theory, putting in place detailed proofs that establish the equivalence of different definitions of the different authors. \\

This is the first in a series of two papers where we establish the fundamental properties of our resulting uniform theory in the second paper~\cite{ghe24}. For instance, our theory detects whether a group (object) has finite cohomological dimension. It satisfies a form of dimension shifting, meaning that one can express a Tate cohomology group of any degree isomorphically in terms of any other degree. It satisfies an Eckmann–Shapiro Lemma, meaning that Tate cohomology of finite index (closed) subgroups of a discrete or profinite group can be described by Tate cohomology of the entire group. Cup products and Yoneda products are developed in Tate cohomology under certain conditions and their key properties are established. We refer to the paper for the details. \\

Since the theory developed in this present paper is based on the above mentioned approaches, let us briefly summarise them. In the way they are originally phrased, the approaches pertain to all groups and more generally, to modules over a ring. D.\! J.\! Benson and J.\! F.\! Carlson in~\cite{ben92} obtain their Tate cohomology groups $\widehat{H}^{\bullet}(G, A)$ by colimits of (stable) Hom-functors using a stabilisation process, making their approach akin to stable module theory. F.\! Goichot in~\cite{goi92} defines Tate cohomology groups $\widehat{H}^{\bullet}(G, A)$ using a hypercohomology construction which he attributes to P.\! Vogel. Thus, we shall refer to this approach as P.\! Vogel's approach from now on. Two approaches can be found in G.\! Mislin's paper~\cite{mis94}. In the first, Tate cohomology groups $\widehat{H}^{\bullet}(G, A)$ are obtained by colimits of ordinary cohomology groups through connecting homomorphisms. In the second, Tate cohomology groups $\widehat{H}^{\bullet}(G, A)$ are obtained as colimit of left satellite functors of ordinary cohomology groups. Left satellite functors were treated as an alternative foundation of homological algebra in the third chapter of H.\! Cartan and S.\! Eilenberg's book on homological algebra~\cite{car56} from 1956. We refer to Section~\ref{sec:constructions} for the details of each of these approaches. \\

Our theory generalises a few constructions from the literature that we outline in the following. It generalises F.\! T.\! Farrell's approach from~\cite{far77} as we demonstrate in~\cite[Lemma~1.11]{ghe24}. F.\! T.\! Farrell formulated it originally for groups of finite virtual cohomological dimension where P.\! Symonds extended it to profinite groups of finite virtual cohomological dimension in~\cite[p.~34]{sym07}. Our theory applies to all groups and to all profinite groups. L.\! L.\! Avramov and O.\! Veliche generalised the above mentioned approaches to Ext-functors of modules over a commutative Noetherian ring $R$ and used them to characterise when $R$ is regular, complete intersection or Gorenstein. Our theory encompasses their ring theoretic work. At the end of the introduction we detail how our theory relates to generalised constructions of Tate cohomology in the realm of homotopical algebra. The generality of our theory prompts us to introduce the following terminology. Whenever we apply our theory it to groups or to topological groups, we follow the convention from~\cite{kro95} and term the resulting generalisation complete cohomology. Analogously, we term the generalised Ext-functors resulting from our theory completed Ext-functors. \\

We provide a detailed outline of our theory in the next three paragraphs. Connecting homomorphisms can be seen as a starting point of its development. They are one of the most important structures for cohomology. Particularly in group cohomology, they constitute fundamental computational tools. Thus far, authors have only demonstrated that the complete cohomology groups $\widehat{H}^{\bullet}(G, A)$ resulting from the above approaches are isomorphic. Even if authors consider connecting homomorphisms $\widehat{\delta}^{\bullet}: \widehat{H}^{\bullet}(G, -) \rightarrow \widehat{H}^{\bullet+1}(G, -)$, they do not devise any explicit formulae for them which created issues in the following two instances. G.\! Mislin in~\cite{mis94} provides a connecting homomorphism for D.\! J.\! Benson and J.\! F.\! Carlson's approach that is incorrect (Remark~\ref{rem:mislinsmistake}). S.\! Guo and L.\! Liang in~\cite{guo23} provide a connecting homomorphism for P.\! Vogel's approach that might render it incompatible to the other approaches (Question~\ref{ques:comparewithguoliang}). To remedy these issues, we provide the first explicit formulae in the literature for the connecting homomorphisms in terms of each approach (Theorem~\ref{thm:satelliteconnmorphisms}, Definition~\ref{defn:resolconnmorph}, Theorem~\ref{thm:naiveconnmorph}, Definition~\ref{defn:hyperconnmorph}). As it has not been done in the literature either, we develop for any morphism $f: A \rightarrow B$ the first formulae for induced morphisms $\widehat{H}^n(G, f): \widehat{H}^n(G, A) \rightarrow \widehat{H}^n(G, B)$ in terms of each construction (Theorem~\ref{thm:inducedsatellitemorphs}, Definition~\ref{defn:resolindmorph}, Definition~\ref{defn:inducednaivemorph}, Definition~\ref{defn:extchainmaps}). Our novel formulae for connecting homomorphisms and induced morphisms provide the cornerstone for our uniform treatment of the entire theory in which we prove that all approaches are equivalent. \\

The key notion in our uniform treatment of the theory is a completion of cohomological functors. Given the prominence of connecting homomorphisms, one can view cohomological functors as a generalisation of group cohomology only preserving the connecting homomorphism. More specifically, if $\mathcal{C}$, $\mathcal{D}$ are abelian categories, then a family of additive functors $(T^n: \mathcal{C} \rightarrow \mathcal{D})_{n \in \mathbb{Z}}$ is a cohomological functor if there are connecting homomorphisms $(\delta^n: T^n \rightarrow T^{n+1})_{n \in \mathbb{Z}}$ satisfying two straightforward axioms~\cite[p.~201--202]{kro95}. In particular, if $G$ is a discrete group, $R$ a discrete ring and $A$ a  discrete $R$-module, then  setting $H_R^n(G, -) = 0$ and $\mathrm{Ext}_R^n(A,-) = 0$ for $n < 0$ renders group cohomology and Ext-functors into cohomological functors~\cite[p.~201]{kro95}, \cite[p.~295]{mis94}. G.\! Mislin shows in~\cite{mis94} that the generalisations of Tate cohomology by his two approaches can be obtained by a completion of cohomological functors. Following the convention from~\cite[p.~197/202]{kro95}, a Mislin completion of a cohomological functor $T^{\bullet}: \mathcal{C} \rightarrow \mathcal{D}$ is another cohomological functor $\widehat{T}^{\bullet}: \mathcal{C} \rightarrow \mathcal{D}$ together with a morphism $\Phi^{\bullet}: T^{\bullet} \rightarrow \widehat{T}^{\bullet}$ such that $\widehat{T}^n(P) = 0$ for any projective $P \in \mathrm{obj}(\mathcal{C})$ and $n \in \mathbb{Z}$. It satisfies the universal property that any morphism $T^{\bullet} \rightarrow V^{\bullet}$ to a cohomological functor $V^{\bullet}: \mathcal{C} \rightarrow \mathcal{D}$ also vanishing on projectives factors uniquely through $\Phi^{\bullet}$~\cite[Definition~2.1]{mis94}. By its universal property, any Mislin completion is unique up to isomorphism~\cite[p.~202]{kro95}. This completion of cohomological functors tracing back to G.\! Mislin's paper~\cite{mis94} bears the advantage that it relates group cohomology with complete cohomology in a unique manner. \\

Let us detail the extent and the generality of the uniform theory resulting from our use of Mislin completions. The premise is to formulate the approaches by D.\! J.\! Benson, J.\! F.\! Carlson, P.\! Vogel and G.\! Mislin in terms of Mislin completions. As G.\! Mislin's concern was group cohomology, he only constructed in~\cite{mis94} Mislin completions of functors $T: \mathcal{C} \rightarrow \mathcal{D}$ in case $\mathcal{C}$ and $\mathcal{D}$ are categories of modules over a ring. We generalise his two approaches under the mild assumptions that $\mathcal{C}$ has enough projective objects and that in $\mathcal{D}$ all countable direct limits exist and are exact (Theorem~\ref{thm:satellitemislin} and Theorem~\ref{thm:resolsatellites}). Direct limits are a specific kind of colimits~\cite[p.~14]{rib10} and are vital to these constructions. We establish the unprecedented fact that there are countably many distinct constructions of Mislin completions (Theorem~\ref{thm:countmanyconstrs}). In particular, that implies that there are countably many distinct generalisations of Tate cohomology. If $\mathbf{Ab}$ denotes the category of abelian groups, then we generalise P. Vogel’s approach from~\cite{goi92} as well as D. J. Benson and J. F. Carlson’s approach~\cite{ben92} so that they give rise to the Mislin completion of Ext-functors $\widehat{\mathrm{Ext}}_{\mathcal{C}}^{\bullet}(A, -): \mathcal{C} \rightarrow \mathbf{Ab}$ (Theorem~\ref{thm:resolbenson} and Theorem~\ref{thm:hyperresol}). As indicated above, we develop the first explicit formulae for connecting homomorphisms and induced morphisms of Mislin completions in terms of each of the above mentioned constructions. In summary, Mislin completions together with the formulae for their connecting homomorphisms and induced morphisms form our uniform theory that generalises D.\! J.\! Benson, J.\! F.\! Carlson, P.\! Vogel and G.\! Mislin's approaches and that thus generalises Tate cohomology. \\

As a remarkable feature, our generalisation of Tate cohomology is applicable to condensed mathematics and thus to most topological groups of interest. Condensed mathematics a powerful novel theory developed by D.\! Clausen and P.\! Scholze in 2018~\cite{sch19}. In a nutshell, it provides a unified approach for studying topological groups, rings and modules~\cite[p.~6]{sch19}. In~\cite{sch20}, P.\! Scholze writes that he wants ``to make the strong claim that in the foundations of mathematics, one should replace topological spaces with condensed sets''. Here, a condensed set can be thought as a sheaf of sets. Let us substantiate P.\! Scholze's claim. In practice, most topological spaces (resp.\! groups, rings, etc) are $T1$, meaning that all their points are closed. These can be functorially turned into condensed sets (resp.\! groups, rings, etc), meaning that one can work only with the corresponding condensed object. We refer to Section~\ref{sec:condmaths} for the details. Condensed mathematics provides an example of a Mislin completion that cannot be treated through any previous framework from the literature. If $\mathcal{R}$ is a condensed ring, $\mathrm{Cond}(\mathrm{Mod}(\mathcal{R}))$ the category of condensed $\mathcal{R}$-modules and $\mathrm{Cond}(\mathbf{Ab})$ the category of condensed abelian groups, then the  Mislin completion of the enriched Ext-functor
\[ \underline{\widehat{\mathrm{Ext}}}_{\mathcal{R}}^{\bullet}(A, -): \mathrm{Cond}(Mod(\mathcal{R})) \rightarrow \mathrm{Cond}(\mathbf{Ab}) \]
can be constructed only by means of our uniform theory (Remark~\ref{rem:differencetoguoliang}). One can define complete cohomology for all condensed groups and thus, for all $T1$ topological groups, meaning for most topological groups as the following list demonstrates (Theorem~\ref{thm:condgroupcohom}).
\begin{itemize}
\item Complete cohomology can be defined for all Lie groups and in particular, for $p$-adic analytic groups.
\item Complete cohomology can be defined for any Galois group or equivalently, for any profinite group. More specifically, a profinite group is an inverse limit of finite discrete groups, meaning that it can be assembled from finite groups in a particular manner. Profinite groups are of interest because they possess properties of infinite groups such as finite generation, but also properties of finite groups such as Sylow subgroups. Moreover, they have recently featured in $3$-manifold topology via the study of profinite rigidity of their fundamental groups. Condensed mathematics is particularly well suited for investigating the cohomology of profinite groups (Lemma~\ref{lem:solidandctscohom}).
\item Complete cohomology can be defined for totally disconnected locally compact (tdlc) groups that are becoming increasingly fashionable. Tdlc groups can be defined as locally profinite groups, thus they generalise both discrete and profinite groups. One can interpret their profinite open subgroups as taking the role of finite subgroups in discrete groups. Tdlc groups play a crucial role in understanding the structure of the more general locally compact groups. Examples of tdlc groups also include $p$-adic analytic groups and automorphisms of locally finite trees with the compact-open topology.
\item However, complete cohomology cannot be defined for algebraic groups via condensed mathematics. Namely, any topological group can be turned into a condensed group because any cartesian product of two topological spaces results in a product of their corresponding condensed sets. In contrast, a product of two algebraic varieties with the Zariski topology does not result in a product of the corresponding condensed sets.
\end{itemize}

Calculations of complete cohomology have been performed beyond finite groups, but are notoriously difficult. This is already indicated by the fact that the complete cohomology group $\widehat{H}^0(G, \mathbb{Z})$ for a group $G$ vanishes if and only if $G$ has finite cohomological dimension~\cite[Lemma~4.2.4]{kro95}. In~\cite{dem08}, F.\! Dembegioti calculates the zeroeth complete cohomology group of a class of polycyclic groups, but also explains why one cannot determine a general formula calculating the zeroeth complete cohomology groups for all polycyclic groups. On the one hand, the cohomology of the free abelian group of countable rank $\bigoplus_{n \in \mathbb{N}} \mathbb{Z}$ is isomorphic to its complete cohomology by~\cite[p.~297]{mis94}. On the other hand, all cohomology groups of the free product $\bigast_{n \in \mathbb{N}} \mathbb{Z}^n$ are distinct from its complete cohomology groups~\cite[p.~432]{jo09}. In~\cite[Corollary~2.4 and Theorem~A]{dem05}, F.\! Dembegioti calculates the zeroeth complete cohomology group of $\bigast_{n \in \mathbb{N}} \mathbb{Z}^n$. Moving away from (discrete) groups, profinite groups are topological groups that are inverse limits of finite discrete groups, meaning that they can be assembled from finite groups in a particular manner. In his paper~\cite[p.~34]{sym07}, P.\! Symonds generalises F.\! T.\! Farrell's approach to profinite groups of finite virtual cohomological dimension. Using this, he computes in~\cite{sym04} the complete cohomology of the Morava stabiliser group $S_{p-1}$ with coefficients in the moduli space $E_{p-1}$ for odd primes $p$. To our knowledge, this is the only calculation of complete cohomology for a non-discrete topological group. \\

Because our theory enables one to define complete cohomology for most topological groups, the legitimate question arises to what calculations it can be employed. As a starting point, we suggest to compute complete cohomology groups of profinite groups because there is already an extensive theory on their cohomology as can be seen in~\cite{rib10} and \cite{wil98}. Inspired by the above mentioned work of F.\! Dembegioti, we propose the following two questions.

\begin{ques}
Let $\prod_{\mathbb{N}} \mathbb{Z}_p$ denote the free abelian pro-$p$ group over $\mathbb{N}$ as in~\cite[Example~3.3.8(c)]{rib10}. Is its cohomology $H_{\mathbb{Z}_p}^{\bullet}(\prod_{\mathbb{N}} \mathbb{Z}_p, -)$ isomorphic to its complete cohomology $\widehat{H}_{\mathbb{Z}_p}^{\bullet}(\prod_{\mathbb{N}} \mathbb{Z}_p, -)$ as a cohomological functor?
\end{ques}

\begin{ques}
Set $G_n = \mathbb{Z}_p^n$ for $n \in \mathbb{N}$ and $G_{\infty} = \lbrace 1 \rbrace$. Define $G = \bigsqcup_{n \in \mathbb{N} \cup \lbrace \infty \rbrace} G_n$ as the free pro-$p$ product over the one-point compactification of the integers $\mathbb{N} \cup \lbrace \infty \rbrace$ using the definitions of~\cite[p.~137--140]{rib17}. Can one compute the zeroeth complete cohomology $\widehat{H}_{\mathbb{Z}_p}^0(G, A)$ for a $\mathbb{Z}_p\llbracket G \rrbracket$-module $A$?
\end{ques}

It is unlikely that one can compute the above complete cohomology groups using previous generalisations of Tate cohomology. Namely, in the case of discrete groups, the complete cohomology groups of the free abelian group $\bigoplus_{n \in \mathbb{N}} \mathbb{Z}$ and of the free product $\bigast_{n \in \mathbb{N}} \mathbb{Z}^n$ cannot be calculated using F.\! T.\! Farrell's approach~\cite[Example~5.3]{cor97}, \cite[p.~119--120]{dem05}. In particular, P.\! Symonds generalisation of F.\! T.\! Farrell's approach from~\cite[p.~34]{sym07} is unlikely to apply to the above questions. \\

We detail how our work relates to generalised constructions of Tate cohomology in the realm of homotopical algebra. Recently, S. Guo and L. Liang have generalised in~\cite{guo23} the above constructions by D.\! J.\ Benson and J.\ F.\! Carlson~\cite{ben92}, by G.\! Mislin~\cite{mis94} and by P.\! Vogel~\cite{goi92} for relative Ext-functors over any abelian category $\mathcal{C}$ with a special precovering subcategory $\mathcal{W}$. Although they prove that their completed Ext-functors $\widehat{\mathrm{Ext}}_{\mathcal{WC}}^n(A, -)$ form a cohomological functor, we do not know whether they also form a Mislin completion of the Ext-functors $\mathrm{Ext}_{\mathcal{WC}}^{\bullet}(A, -)$ (Question~\ref{ques:comparewithguoliang}). Using a version of D.\! J.\ Benson and J.\ F.\! Carlson's approach, A.\! Beligiannis and I.\! Reiten define in~\cite[Section~IX.2]{bel07} a completed relative Ext-functor $\widehat{\mathrm{Ext}}_{(\mathcal{X}, \mathcal{Y})}^{\bullet}(-, -)$ for an abelian category $\mathcal{C}$ containing a particular pair of subcategories $(\mathcal{X}, \mathcal{Y})$. They prove that $\mathrm{Ext}_{\mathcal{X} \cap \mathcal{Y}}^{\bullet}(A, -) \rightarrow \widehat{\mathrm{Ext}}_{(\mathcal{X}, \mathcal{Y})}^{\bullet}(A, -)$ forms a Mislin completion for any $A \in \mathrm{obj}(\mathcal{C})$. J.\! Hu et al.\! use in~\cite{hu21} a version of P.\! Vogel's approach to establish completed Ext-functors $\widehat{\mathrm{Ext}}_{\mathcal{C}}^{\bullet}(-, -)$ for an extriangulated category $\mathcal{C}$ where extriangulated categories are a generalisation of exact categories and triangulated categories. Among their applications, they provide a criterion for the validity of the Wakamatsu Tilting Conjecture. A generalisation of Tate cohomology in the language of spectra can be found in~\cite{kle02}. For an overview of more generalisations of Tate cohomology and their applications, the reader is referred to M.\! Paganin's survey paper~\cite{pag16}. Our contribution stands out because it provides a uniform treatment of the entire theory and establishes fundamental properties such as dimension shifting, an Eckmann–Shapiro Lemma and cohomology products. \\

Let us point out two variants of the constructions considered in this paper that are covered in the literature, but that we shall not pursue. One could perform the constructions of D.\! J.\! Benson, J.\! F.\! Carlson, G.\! Mislin and P.\! Vogel for \textit{homological functors} instead of cohomological functors. O.\! Celikbas et al.\! generalise these constructions for homological functors of modules over rings, but they turn out to be not equivalent. Hence, one cannot expect to encounter the same uniform theory as we develop it for cohomological functors. As was first observed by B.\! E.\! A.\! Nucinkis in~\cite{nuc98}, one could perform the above constructions dually by assuming that there are enough injectives instead of enough projectives. More specifically, a ``Nucinkis completion'' of a cohomological functor $T^{\bullet}: \mathcal{C} \rightarrow \mathcal{D}$ is a morphism $T^{\bullet} \rightarrow \widetilde{T}^{\bullet}$ such that $\widetilde{T}^{\bullet}$ vanishes on injectives and is universal with respect to this property. S.\! Guo and L.\! Liang perform all their constructions of completed relative Ext-functors in a dual manner in~\cite{guo23}. Together with X.\! Yang they provide in~\cite[Main Theorem]{guol23} a criterion for when $\widehat{\mathrm{Ext}}_{\mathcal{C}}^{\bullet}(-, -)$ is isomorphic to $\widetilde{\mathrm{Ext}}_{\mathcal{C}}^{\bullet}(-, -)$. Such criteria are also established by A.\! Beligiannis and I.\! Reiten in~\cite[Theorem~IX.4.4]{bel07} as well as by S.\! Hu et al.\! in~\cite[Theorem~4.4]{huz21}. This ``Nucinkis completion'' would not affect the cohomology of discrete and profinite groups because they already vanish on injectives~\cite[p.~61]{bro82},~\cite[p.~9]{ser94}. As this completion does not generalise Tate cohomology, we shall not consider it further. \\

We showcase in Section~\ref{sec:constructions} the relevant constructions of Mislin completions that we are generalising in this paper. The bulk of the paper consists of Section~\ref{sec:satellites}, Section~\ref{sec:resols}, Section~\ref{sec:naiveconstruction} and Section~\ref{sec:hypercohom} in which we prove that each of these approaches gives rise to a Mislin completion. Accordingly, we develop explicit formulae for the induced morphisms and connecting homomorphisms for each approach in the corresponding section. More specifically, we prove in Section~\ref{sec:satellites} that one approach by G.\! Mislin forms a Mislin completion. In each subsequent section we demonstrate this for one other approach each by constructing isomorphisms of cohomological functors to an approach that already forms Mislin completions. In this manner, we cover another approach also due to G.\! Mislin in Section~\ref{sec:resols}, the approach due to D.\! J.\! Benson and J.\! F.\! Carlson in Section~\ref{sec:naiveconstruction} and the approach due to P.\! Vogel in Section~\ref{sec:hypercohom}. As a noteworthy feature, we apply our results to the powerful novel theory of condensed mathematics in Section~\ref{sec:condmaths}.

\subsection{Notation and terminology}

We adopt the convention that the natural numbers $\mathbb{N}$ include $1$, but do not include $0$. We write $\mathbb{N}_0$ for $\mathbb{N} \cup \lbrace 0 \rbrace$. Moreover, we adopt B.\! Poonen's convention from~\cite{poo18} that a ring is an abelian group together with a totally associative product, meaning a binary associative relation admitting an identity element. In particular, ring homomorphisms are understood to map identity elements to identity elements. We use the symbol $\varinjlim$ only to denote direct limits. If $\mathcal{D}$ is a category, then $\varinjlim_{\mathcal{D}}$ denotes a direct limit in this category and if $I$ is a directed set, then $\varinjlim_{i \in I}$ denotes a direct limit indexed over $I$. We write $\varinjlim_{k \in \mathbb{N}} (M_k, \mu_k)$ if $\mu_k: M_k \rightarrow M_{k+1}$ are the morphisms giving rise to the direct limit. Lastly, we use the same numbering to label diagrams and equations.

\section{Outline of constructions}\label{sec:constructions}

Let us axiomatically define cohomological functors as they are one of the key notions in this paper. For two abelian categories $\mathcal{C}$, $\mathcal{D}$ we call a family of additive functors $(T^n: \mathcal{C} \rightarrow \mathcal{D})_{n \in \mathbb{Z}}$ a cohomological functor if it satisfies the following two axioms~\cite[p.~201--202]{kro95}.

\begin{axm}\label{axm:delta}
For every $n \in \mathbb{Z}$ and short exact sequence $0 \rightarrow A \rightarrow B \rightarrow C \rightarrow 0$ in $\mathcal{C}$, there is natural connecting homomorphism $\delta^n: T^n(C) \rightarrow T^{n+1}(A)$.
\end{axm}

Being natural means in this context that for every commuting diagram in $\mathcal{C}$
\begin{center}
\begin{tikzcd}
0 \arrow[r] & A \arrow[r] \arrow[d, "f"] & B \arrow[r] \arrow[d] & C \arrow[r] \arrow[d, "g"] & 0 \\
0 \arrow[r] & A' \arrow[r] & B' \arrow[r] & C' \arrow[r] & 0
\end{tikzcd}
\end{center}
with exact rows there is a commuting diagram
\begin{center}
\begin{tikzcd}
T^n(A) \arrow[r, "\delta^n"] \arrow[d, "T^n(g)"] & T^{n+1}(A) \arrow[d, "T^{n+1}(f)"] \\
T^n(A') \arrow[r, "\delta^n"] & T^{n+1}(C')
\end{tikzcd}
\end{center}
in $\mathcal{D}$.

\begin{axm}\label{axm:les}
For every short exact sequence $0 \rightarrow A \xrightarrow{\iota} B \xrightarrow{\pi} C \rightarrow 0$ in $\mathcal{C}$ there is a long exact sequence
\[ \dots {} \xrightarrow{T^{n-1}(\pi)} T^{n-1}(C) \xrightarrow{\delta^{n-1}} T^n(A) \xrightarrow{T^n(\iota)} T^n(B) \xrightarrow{T^n(\pi)} T^n(C) \xrightarrow{\delta^n} T^{n+1}(A) \xrightarrow{T^{n+1}(\iota)} {} \dots \]
\end{axm}

Let us axiomatically define morphisms of cohomological functors as in~\cite[p.~202]{kro95}.

\begin{axm}\label{axm:morphisms}
Let $(T^{\bullet}, \delta^{\bullet})$ and $(U^{\bullet},\varepsilon^{\bullet})$ be cohomological functors from $\mathcal{C}$ to $\mathcal{D}$. Then a family of natural transformations $(\nu^n: T^n \rightarrow U^n)_{n \in \mathbb{Z}}$ is a morphism of cohomological functors if for every $n \in \mathbb{Z}$ and any short exact sequence $0 \rightarrow A \rightarrow B \rightarrow C \rightarrow 0$ in $\mathcal{C}$, the square
\begin{center}
\begin{tikzcd}
  T^n(C) \arrow[r, "\delta^n"] \arrow[d, "\nu^n"]
    &T^{n+1}(A) \arrow[d, "\nu^{n+1}"] \\
  U^n(C) \arrow[r, "\varepsilon^n"]
    &U^{n+1}(A)
\end{tikzcd}
\end{center}
commutes.
\end{axm}

We generalise G.\! Mislin's Definition~2.1 from~\cite{mis94} to the greatest extent.

\begin{defn}[Mislin completion]\label{defn:mislincompletion}
Let $(T^{\bullet}, \delta^{\bullet})$ be a cohomological functor from $\mathcal{C}$ to $\mathcal{D}$. Then its Mislin completion is a cohomological functor $(\widehat{T}^{\bullet}, \widehat{\delta}^{\bullet})$ from $\mathcal{C}$ to $\mathcal{D}$ together with a morphism $\nu^{\bullet}: T^{\bullet} \rightarrow \widehat{T}^{\bullet}$ satisfying the following universal property: \newline
\textnormal{1.} $(\widehat{T}^{\bullet}, \widehat{\delta}^{\bullet})$ vanishes on projectives, meaning that $\widehat{T}^n(P) = 0$ for every projective object $P \in \mathrm{obj}(\mathcal{C})$ and every $n \in \mathbb{Z}$. \newline
\textnormal{2.} If $(U^{\bullet}, \varepsilon^{\bullet})$ is any cohomological functor vanishing on projectives, then each morphism $T^{\bullet} \rightarrow U^{\bullet}$ factors uniquely as $T^{\bullet} \xrightarrow{\nu^{\bullet}} \widehat{T}^{\bullet} \rightarrow U^{\bullet}$.
\end{defn}

By virtue of their universal property, Mislin completions are unique up to isomorphism in the following sense. If $(U^{\bullet}, \varepsilon^{\bullet})$ is another Mislin completion of $(T^{\bullet}, \delta^{\bullet})$, then there is an isomorphism $\mu^{\bullet}: \widehat{T}^{\bullet} \rightarrow U^{\bullet}$, meaning that $\mu^n(M): \widehat{T}^n(M) \rightarrow U^n(M)$ is an isomorphism for any $n \in \mathbb{Z}$ and $M \in \mathrm{obj}(\mathcal{C})$~\cite[p.~202]{kro95}. This allows us to state G.\! Mislin's definition from~\cite[p.~297]{mis94} in the greatest generality.

\begin{defn}[Axiomatic, Mislin]\label{defn:axiomatic}
${} \quad {}$
\begin{itemize}
\item For any $A \in \mathrm{obj}(\mathcal{C})$ extend the (enriched or unenriched) Ext-functors to a cohomological functor by setting $\mathrm{Ext}_{\mathcal{C}}^n(A, -) = 0$ for $n < 0$. Define completed Ext-functors as the Mislin completion $(\widehat{\mathrm{Ext}}_{\mathcal{C}}^{\bullet}(A, -), \widehat{\delta}^{\bullet})$.
\item Analogously, if $G$ is a group object in $\mathcal{C}$ and $R$ a ring object, then extend group cohomology to a cohomological functor $(H_R^{\bullet}(G, -), \delta^{\bullet})$ by imposing $H_R^n(G, -) = 0$ for $n < 0$. Define complete cohomology as the Mislin completion $(\widehat{H}_R^{\bullet}(G, -), \widehat{\delta}^{\bullet})$.
\end{itemize}
\end{defn}

To ensure that Mislin completions exist, the domain category $\mathcal{C}$ of the cohomological functor needs to have enough projectives on the one hand. On the other hand, in the codomain category $\mathcal{D}$ all countable direct limits need to exist and be exact, which we define as follows.

\begin{defn}\label{defn:directlimit}
A partially ordered set $(I, \leq)$ is a directed set if for every $i, j \in I$ there is $k \in I$ such that $i,j \leq k$~\cite[p.~1]{rib10}. According to~\cite[p.~14]{rib10}, a diagram $\lbrace D_i \rbrace_{i \in I}$ in $\mathcal{D}$ indexed over a directed set is called a direct system in $\mathcal{D}$. More formally, $I$ can be turned into a category whose objects are its elements and there is a unique morphism $i \rightarrow j$ whenever $i \leq j \in I$. Then a direct system is a covariant functor $I \rightarrow \mathcal{D}, i \mapsto D_i$. A direct limit $\varinjlim_{i \in I} D_i$ in $\mathcal{D}$ is a colimit of a direct system $\lbrace D_i \rbrace_{i \in I}$. Direct limits in $\mathcal{D}$ are called exact if for every direct system of short exact sequence $\lbrace 0 \rightarrow A_i \rightarrow B_i \rightarrow C_i \rightarrow 0 \rbrace_{i \in I}$ also
\[ 0 \rightarrow \varinjlim_{i \in I} A_i \rightarrow \varinjlim_{i \in I} B_i \rightarrow \varinjlim_{i \in I} C_i \rightarrow 0 \]
is a short exact sequence~\cite[\href{https://stacks.math.columbia.edu/tag/079A}{Tag 079A}]{stacks-project}.
\end{defn}

In order to clarify where the above assumptions are needed, we present what we term the satellite functor construction, which is due to G.\! Mislin. First, we introduce left satellite functors. Using the above assumption that $\mathcal{C}$ has enough projective objects, there is for any $M \in \mathrm{obj}(\mathcal{C})$ a short exact sequence $0 \rightarrow K \rightarrow P \rightarrow M \rightarrow 0$ in $\mathcal{C}$ with $P$ projective. For a cohomological functor $(T^{\bullet}, \delta^{\bullet})$, we define the zeroeth left satellite functor of $T^n$ as $S^0T^n := T^n$, the first left satellite functor as
\[ S^{-1}T^n(M) := \mathrm{Ker}\big(T^n(K) \rightarrow T^n(P)\big) \]
and the $k^{\text{th}}$ left satellite functor as $S^{-k}T^n := S^{-1}(S^{-k+1}T^n)$ for $k \geq 2$~\cite[p.~36]{car56}. It is shown in~\cite[Section~III.1]{car56} that left satellite functors do not depend on the choice of short exact sequence. Since they have been defined as kernels, it follows from Axiom~\ref{axm:les} that $\delta^n: T^n \rightarrow T^{n+1}$ induces a morphism $\underline{\delta}^n: T^n(M) \rightarrow S^{-1}T^{n+1}(M)$ and therefore
\[ S^{-k}\underline{\delta}^{n+k}: S^{-k}T^{n+k}(M) \rightarrow S^{-k-1}T^{n+k+1}(M) \]
for any $k \in \mathbb{N}$~\cite[p.~207--208]{kro95}. We extend G.\! Mislin's construction from~\cite[p.~293]{mis94} by explicitly using our assumption that all countable  direct limits exist in the codomain category $\mathcal{D}$ of $T^{\bullet}$.

\begin{defn}[Satellite functor construction, Mislin]
The Mislin completion of a cohomological functor $(T^{\bullet}, \delta^{\bullet})$ can be defined as
\[ \widehat{T}^n(M) := \varinjlim_{k \in \mathbb{N}_0} (S^{-k}T^{n+k}(M), S^{-k}\underline{\delta}^{n+k}) \]
for any $M \in \mathrm{obj}(\mathcal{C})$ and $n \in \mathbb{Z}$. Accordingly,
\[ \widehat{H}_R^n(G, M) := \varinjlim_{k \in \mathbb{N}_0} (S^{-k}H_R^{n+k}(G, M), S^{-k}\underline{\delta}^{n+k}) \]
is a definition of complete cohomology.
\end{defn}

In order that the above forms a cohomological functor, we require the very assumption that all countable direct limits in $\mathcal{D}$ are exact. \\

Let us go over to what we term the resolution construction that occurs in~\cite[Lemma~B.3]{cel17} and can be retrieved from page~299 in G.\! Mislin's paper~\cite{mis94}. If $(M_n)_{n \in \mathbb{N}_0}$ is a projective resolution of $M \in \mathrm{obj}(\mathcal{C})$, let us define $\widetilde{M}_0 := M$ and $\widetilde{M}_k := \mathrm{Ker}(M_{k-1} \rightarrow \widetilde{M}_{k-1})$ for $k \in \mathbb{N}$. This is called the $k^{\text{th}}$ syzygy of $M_{\bullet}$ in the Gorenstein context~\cite[p.~89]{pag16}. The choice of our notation is meant to reflect that our syzygies do not necessarily arise from a specific choice of projective resolution as in \cite{kro95} and \cite{mis94}. Since for every $k \in \mathbb{N}_0$ we have the short exact sequence $0 \rightarrow \widetilde{M}_{k+1} \rightarrow M_k \rightarrow \widetilde{M}_k \rightarrow 0$, there is a connecting homomorphism $\delta^{n+k}: T^{n+k}(\widetilde{M}_k) \rightarrow T^{n+k+1}(\widetilde{M}_{k+1})$ for every $n \in \mathbb{Z}$. Then the following definition makes it more apparent why Mislin completions vanish on projective objects.

\begin{defn}[Resolution construction, Mislin]\label{defn:resolsinoutline}
The Mislin completion of a cohomological functor $(T^{\bullet}, \delta^{\bullet})$ can be defined as 
\[ \widehat{T}^n(M) := \varinjlim_{k \in \mathbb{N}_0} (T^{n+k}(\widetilde{M}_k), \delta^{n+k}) \]
for any $n \in \mathbb{Z}$ and $M \in \mathrm{obj}(\mathcal{C})$. Accordingly,
\[ \widehat{H}_R^n(G, M) := \varinjlim_{k \in \mathbb{N}_0} (H_R^{n+k}(G, \widetilde{M}_k), \delta^{n+k}) \]
is a definition of complete cohomology.
\end{defn}

The next two constructions only give rise to completed unenriched Ext-functors. Let $A_{\bullet}$, $B_{\bullet}$ are projective resolutions of $A, B \in \mathrm{obj}(\mathcal{C})$ and let $\widetilde{f}_{n+k}: \widetilde{A}_{n+k} \rightarrow \widetilde{B}_k$ be a morphism for $n \in \mathbb{Z}$ and $k \in \mathbb{N}_0$ such that $n+k \geq 0$. Then we can write the commuting diagram
\begin{center}
\begin{tikzcd}
    0 \arrow[r] & \widetilde{A}_{n+k+1} \arrow[r] \arrow[d, "\widetilde{f}_{k+1}"] & A_{n+k} \arrow[r] \arrow[d, "f_{k+1}"] & \widetilde{A}_{n+k} \arrow[r] \arrow[d, "\widetilde{f}_k"] & 0 \\
    0 \arrow[r] & \widetilde{B}_{k+1} \arrow[r] & B_{k+1} \arrow[r] & \widetilde{B}_k \arrow[r] & 0
\end{tikzcd}
\end{center}
whose terms arise as follows. Since the bottom row is exact and the term $A_{n+k}$ projective, there is a lift $f_k$ of $\widetilde{f}_k$ making the right-hand square commute. Because $\widetilde{B}_{k+1} \rightarrow B_{k+1}$ is a kernel, there is a morphism $\widetilde{f}_{k+1}$ making the left-hand side commute. If $\mathrm{Hom}_{\mathcal{C}}(-, -)$ denotes the (unenriched) Hom-functor in $\mathcal{C}$, then $\mathrm{Hom}_{\mathcal{C}}(\widetilde{A}_{n+k}, \widetilde{B}_k)$ is an abelian group by virtue of $\mathcal{C}$ being an abelian category. We define $\mathcal{P}_{\mathcal{C}}(\widetilde{A}_{n+k}, \widetilde{B}_k)$ to be the subgroup of $\mathrm{Hom}_{\mathcal{C}}(\widetilde{A}_{n+k}, \widetilde{B}_k)$ consisting of all morphisms factoring through a projective object and write the quotient as $[\widetilde{A}_{n+k}, \widetilde{B}_k]_{\mathcal{C}} := \mathrm{Hom}_{\mathcal{C}}(\widetilde{A}_{n+k}, \widetilde{B}_k)/ \mathcal{P}_{\mathcal{C}}(\widetilde{A}_{n+k}, \widetilde{B}_k)$~\cite[p.~203]{kro95}. As in the case of modules over a ring covered by~\cite[p.~204]{kro95}, one can prove that
\begin{align*}
t_{\widetilde{A}_{n+k}, \widetilde{B}_k}: [\widetilde{A}_{n+k}, \widetilde{B}_k]_{\mathcal{C}} &\rightarrow [\widetilde{A}_{n+k+1}, \widetilde{B}_{k+1}]_{\mathcal{C}}, \\
\widetilde{f}_k + \mathcal{P}_{\mathcal{C}}(\widetilde{A}_{n+k}, \widetilde{B}_k) &\mapsto \widetilde{f}_{k+1} + \mathcal{P}_{\mathcal{C}}(\widetilde{A}_{n+k+1}, \widetilde{B}_{k+1})
\end{align*}
is a well defined homomorphism. Using this, we generalise the following construction from D.\! J.\! Benson and J.\! F.\! Carlson's paper~\cite[p.~109]{ben92}.

\begin{defn}[Naïve construction, Benson \& Carlson]
For any $n \in \mathbb{Z}$, we can define the $n^{\text{th}}$ completed unenriched Ext-functor as
\[ \widehat{\mathrm{Ext}}_{\mathcal{C}}^n(A, B) := \varinjlim_{k \in \mathbb{N}_0, n+k \geq 0} ([\widetilde{A}_{n+k}, \widetilde{B}_k]_{\mathcal{C}}, t_{\widetilde{A}_{n+k}, \widetilde{B}_k}) \, . \]
In particular, if $R_{\bullet}$ is a projective $R[G]$-resolution of $R \in \mathrm{obj}(Mod_R(G))$, we can define complete unenriched cohomology as
\[ \widehat{H}_R^n(G, B) := \varinjlim_{k \in \mathbb{N}_0, n+k \geq 0} ([\widetilde{R}_{n+k}, \widetilde{B}_k]_{\mathcal{C}}, t_{\widetilde{R}_{n+k}, \widetilde{B}_k}) \, . \]
\end{defn}

Lastly, we present what we call the hypercohomology construction of complete cohomology. We define the chain complex $(A_n')_{n \in \mathbb{Z}}$ by $A_n' = A_n$ for $n \geq 0$ and $A_n' = 0$ for $n < 0$ and similarly $(B_n')_{n \in \mathbb{Z}}$~\cite[p.~209]{kro95}. Define the hypercohomology complex $(\mathrm{Hyp}_{\mathcal{C}}(A_{\bullet}', B_{\bullet}')_n, d^n)_{n \in \mathbb{Z}}$ by having $n$-cochains 
\[ \mathrm{Hyp}_{\mathcal{C}}(A_{\bullet}', B_{\bullet}')_n = \prod_{k \in \mathbb{Z}} \mathrm{Hom}_{\mathcal{C}}(A_{k+n}', B_k') \, . \]
To ease notation in the following, we view abelian groups as $\mathbb{Z}$-modules. If we denote by $a_n: A_n' \rightarrow A_{n-1}'$ and $b_n: B_n' \rightarrow B_{n-1}'$ the differentials induced from the respective projective resolution, we define for $n \in \mathbb{Z}$ the differential
\begin{align}
d^n: \mathrm{Hyp}_{\mathcal{C}}(A_{\bullet}', B_{\bullet}')_n &\rightarrow \mathrm{Hyp}_{\mathcal{C}}(A_{\bullet}', B_{\bullet}')_{n+1} \label{eq:hypercohombd} \\ (\varphi_{n+k})_{k \in \mathbb{Z}} &\mapsto (b_{k+1} \circ \varphi_{n+k+1} - (-1)^n \varphi_{n+k} \circ a_{n+k+1})_{k \in \mathbb{Z}} \nonumber
\end{align}
Let us define the bounded complex $\mathrm{Bdd}_{\mathcal{C}}(A_{\bullet}', B_{\bullet}')_{n \in \mathbb{Z}}$ as the subcomplex of $\mathrm{Hyp}_{\mathcal{C}}(A_{\bullet}', B_{\bullet}')_{n \in \mathbb{Z}}$ given by
\[ \mathrm{Bdd}_{\mathcal{C}}(A_{\bullet}', B_{\bullet}')_n = \bigoplus_{k \in \mathbb{Z}} \mathrm{Hom}_{\mathcal{C}}(A_{k+n}', B_k') \, . \]
Define the Vogel complex as the quotient complex~\cite[p.~209]{kro95}
\[ \mathrm{Vog}_{\mathcal{C}}(A_{\bullet}', B_{\bullet}')_{n \in \mathbb{Z}} := \Big(\mathrm{Hyp}_{\mathcal{C}}(A_{\bullet}', B_{\bullet}')_n / \mathrm{Bdd}_{\mathcal{C}}(A_{\bullet}', B_{\bullet}')_n \Big)_{n \in \mathbb{Z}} \, . \]
By this, we generalise Definition~1.2 from F.\! Goichot's paper~\cite{goi92} where he attributes it to P.\! Vogel on page~39.

\begin{defn}[Hypercohomology construction, Vogel]\label{defn:vogel}
For $n \in \mathbb{Z}$ we can define the $n^{\text{th}}$ completed unenriched Ext-functor as
\[ \widehat{\mathrm{Ext}}_{\mathcal{C}}^n(A, B) := H^n((\mathrm{Vog}_{\mathcal{C}}(A_{\bullet}', B_{\bullet}')_k, d^k)_{k \in \mathbb{Z}}) \, . \]
We can thus define complete unenriched cohomology as
\[ \widehat{H}_R^n(G, M) := H^n((\mathrm{Vog}_R(R_{\bullet}', B_{\bullet}')_k, d^k)_{k \in \mathbb{Z}})\, . \]
\end{defn}

Let us remark why it is unlikely that the previous two constructions could yield Mislin completions of more general enriched Ext-functors. Enriched Hom-functors are particular bifunctors of the form
\[ \underline{\mathrm{Hom}}_{\mathcal{C}}(-, -): \mathcal{C}^{\mathrm{op}} \times \mathcal{C} \rightarrow \mathcal{D} \]
where $\mathcal{D}$ is an abelian category that does not need to be $\mathbf{Ab}$. Accordingly, one can define enriched Ext-functors $\underline{\mathrm{Ext}}_{\mathcal{C}}^{\bullet}(-, -): \mathcal{C}^{\mathrm{op}} \times \mathcal{C} \rightarrow \mathcal{D}$, which are relevant in the context of condensed mathematics. For the naïve construction, one aims to find a morphism $\mathrm{Hom}_{\mathcal{C}}(\widetilde{A}_{n+k}, \widetilde{B}_k) \rightarrow \mathrm{Hom}_{\mathcal{C}}(\widetilde{A}_{n+k+1}, \widetilde{B}_{k+1})$. If one tries to lift along the morphisms induced by $A_{n+k} \rightarrow \widetilde{A}_{n+k}$ and $B_{n+k} \rightarrow \widetilde{B}_{n+k}$, one requires that $\mathrm{Hom}_{\mathcal{C}}(A_{n+k}, -)$ preserves epimorphisms. For Hom-functors, this role is exactly played by projective objects~\cite[Lemma~2.2.3]{wei94}. For the hypercohomology construction we require that the coproduct $\mathrm{Bdd}_{\mathcal{C}}(A_{\bullet}', B_{\bullet}')$ maps into the product $\mathrm{Hyp}_{\mathcal{C}}(A_{\bullet}', B_{\bullet}')$. As we shall see in Subsection~\ref{subsec:hyperindconn}, the cohomology groups of the Vogel complex $\mathrm{Vog}_{\mathcal{C}}(A_{\bullet}', B_{\bullet}')$ correspond exactly to almost chain maps modulo almost chain homotopy. However, this implies that the objects $\mathrm{Hom}_{\mathcal{C}}(A_p, B_q) \in \mathrm{obj}(\mathcal{D})$ describe morphisms of the form $A_p \rightarrow B_q$ in $\mathcal{C}$.

\begin{conv}
As we have provided explicit constructions for the approaches by G.\! Mislin, by D.\! J.\! Benson and J.\! F.\! Carlson and by P.\! Vogel, we establish the following convention for the remainder of the paper. We shall refer to each construction by its name and cede to refer to the author who has introduced it. For instance, we shall only refer to the satellite functor construction from now on and do not mention anymore that its general approach is due to G.\! Mislin. 
\end{conv}

\section{The satellite functor construction and Mislin completions}\label{sec:satellites}

By generalising the satellite functor construction in this section, we establish Mislin completions in full generality and present the very first explicit formulae for the resulting connecting homomorphisms and induced morphisms. More specifically, G.\! Mislin constructs in the proof of \cite[Theorem~2.2]{mis94} Mislin completions through satellite functors only for modules over a ring. We generalise his work such that one can define Mislin completions whenever the domain category $\mathcal{C}$ is an abelian category with enough projective objects and in the codomain category $\mathcal{D}$ all countable direct limits exist and are exact. The explicit formulae for the connecting homomorphisms and induced morphisms of the resulting Mislin completion $(\widehat{T}^{\bullet}, \widehat{\delta}^{\bullet})$ are fundamental for this paper because they form the base for analogous explicit formulae of all other constructions of Mislin completions. \\

For the reader's convenience, we summarise relevant results on left satellite functors from the third chapter of H.\! Cartan and S.\! Eilenberg's book~\cite{car56}. Because A.\! Grothendieck has generalised their satellite functors in his Tôhoku paper~\cite[p.~140--143]{gro57}, their arguments pertain to our categories $\mathcal{C}$ and $\mathcal{D}$. The first relevant result is an explicit description of induced morphisms for left satellite functors that we present in the following. Consider a morphism $f: C' \rightarrow C$ in $\mathcal{C}$. Using the assumption that the category $\mathcal{C}$ has enough projectives, we choose short exact sequences
\[ 0 \rightarrow M' \rightarrow P' \rightarrow C' \rightarrow 0 \quad \text{and} \quad 0 \rightarrow M \rightarrow P \rightarrow C \rightarrow 0 \]
with $P'$ and $P$ projective. Then there exists a lift $\overline{f}^{\ast}: P' \rightarrow P$ rendering the diagram
\begin{equation}\label{diag:satellitemorph}
\begin{tikzcd}
  0 \arrow[r]
    &M' \arrow[r] \arrow[d, "\overline{f}^{\ast}"]
    &P' \arrow[r] \arrow[d, "\overline{f}"]
    &C' \arrow[r] \arrow[d, "f"]
    &0 \\
 0 \arrow[r]
    &M \arrow[r]
    &P \arrow[r]
    &C \arrow[r]
    &0
\end{tikzcd}
\end{equation}
commutative. The morphism $T^n(\overline{f}^{\ast}): T^n(M') \rightarrow T^n(M)$ induces the induced morphism $S^{-1}T^n(f): S^{-1}T^n(C') \rightarrow S^{-1}T^n(C)$. Pictorially,
\begin{center}
\begin{tikzcd}
    S^{-1}T^n(C') \arrow[r, "\varepsilon^{-1}"] \arrow[d, dashed, "S^{-1}T^n(f)"] & T^n(M') \arrow[d, "T^n(\overline{f}^{\ast})"] \\
    S^{-1}T^n(C) \arrow[r, "\varepsilon^{-1}"] & T^n(M)
\end{tikzcd}
\end{center}
where the morphisms $\varepsilon^{-1}$ are the canonical monomorphisms from the kernels. The induced morphisms $S^{-k}T^n(f): S^{-k}T^n(C') \rightarrow S^{-k}T^n(C)$ for higher satellite functors are defined analogously. \\

There are connecting homomorphisms between left satellite functors that are indispensable for our results. For any exact sequence $0 \rightarrow A \rightarrow B \rightarrow C \rightarrow 0$ in $\mathcal{C}$ and any $k \in \mathbb{N}_0$ these connecting homomorphisms take the form $\varepsilon^{-k-1}: S^{-k-1}T^n(C) \rightarrow S^{-k}T^n(A)$. If $B$ is projective, then $\varepsilon^{-k-1}$ is the monomorphism from the kernel $S^{-k-1}T^n(C)$. Otherwise let $0 \rightarrow M \rightarrow P \rightarrow C \rightarrow 0$ be a short exact sequence with $P$ projective. Let $h: P \rightarrow B$ be a morphism lifting $\mathrm{id}_C: C \rightarrow C$ as in the commutative diagram
\begin{equation}\label{diag:delta}
\begin{tikzcd}
  0 \arrow[r]
    &M \arrow[r] \arrow[d, "h^{\ast}"]
    &P \arrow[r] \arrow[d, "h"]
    &C \arrow[r] \arrow[d, "\mathrm{id}_{C}"]
    &0 \\
 0 \arrow[r]
    &A \arrow[r]
    &B \arrow[r]
    &C \arrow[r]
    &0
\end{tikzcd}
\end{equation}
Then 
\[ \varepsilon^{-k-1}: S^{-k-1}T^n(C) \rightarrow S^{-k}T^n(M) \xrightarrow{S^{-k}T^n(h^{\ast})} S^{-k}T^n(A) \]
is a connecting homomorphism where the first morphism is the monomorphism of the kernel $S^{-n-1}T^n(C)$. The term connecting homomorphism is justified because the morphisms $\varepsilon^{-k}$ form natural transformations giving rise to a long exact sequence
\[ \dots {} \rightarrow S^{-k-1}T^n(C) \xrightarrow{\varepsilon^{-k-1}} S^{-k}T^n(A) \rightarrow S^{-k}T^n(B) \rightarrow S^{-k}T^n(C) \xrightarrow{\varepsilon^{-k}} S^{-k+1}T^n(A) \rightarrow {} \dots \]
The starting point of the satellite functor construction of Mislin completions is the following result that can be directly generalised from \cite[Proposition~III.5.2]{car56}.

\begin{lem}\label{lem:uniqueness}
Let $T^{\bullet}$ and $U^{\bullet}$ be cohomological functors and $\Phi^0: T^0 \rightarrow U^0$ be a natural transformation. If for every short exact sequence $0 \rightarrow K \rightarrow P \rightarrow M \rightarrow 0$ with $P$ projective and every $n \in \mathbb{N}$ the sequence $0 \rightarrow U^{-n-1}(M) \rightarrow U^{-n}(K)$ involving the connecting homomorphism is exact, then $\Phi^0$ extends uniquely to a (partial) morphism of cohomological functors $(\Phi^n: T^n \rightarrow U^n)_{n \leq 0}$ only defined for $n \leq 0$. In particular, if also $T^{\bullet}$ has the property that the sequence $0 \rightarrow T^{-n-1}(M) \rightarrow T^{-n}(K)$ is exact for $n \in \mathbb{N}$ and $\Phi^0$ is an equivalence, then the extension $\Phi^{\bullet}$ is an isomorphism.
\end{lem}

G.\! Mislin uses this lemma to perform his satellite functor construction in~\cite[pp.~295--296]{mis94} which we generalise in the following. For any $m \in \mathbb{N}_0$ we can define the cohomological functor
\begin{equation}\label{eq:directsyscohomfunctor}
T^k \langle m \rangle (M) := \begin{cases} S^{k-m}T^m(M) &\text{if } k < m \\ T^k(M) &\text{if } k \geq m \end{cases}
\end{equation}
where the connecting homomorphisms are given by $\delta^k \langle m \rangle = \delta^k: T^k \rightarrow T^{k+1}$ for $k \geq m$ and $\delta^k \langle m \rangle = \varepsilon^{k-m}: S^{k-m}T^m \rightarrow S^{k-m+1}T^m$ for $k < m$. Consider the partial morphism $(\Phi_m^k := \mathrm{id}: T^k \rightarrow T^k \langle m \rangle)_{m \leq k}$ of cohomological functors defined only in degrees $m \leq k$. According to the previous lemma we can extend $\Phi_m^{\bullet}$ uniquely to a morphism of cohomological functors defined in all degrees. Using the same arguments we have for $m \leq n$ a unique morphism $\Phi_{m, n}^{\bullet}: T^{\bullet} \langle m \rangle \rightarrow T^{\bullet} \langle n \rangle$ that is equal to the identity morphism in any degree $n \leq k \in \mathbb{Z}$. Using the uniqueness of all these morphisms, we conclude for $m \leq n \leq o \in \mathbb{N}_0$ that $\Phi_n^{\bullet} = \Phi_{m,n}^{\bullet} \circ \Phi_m^{\bullet}$ and $\Phi_{m,o}^{\bullet} = \Phi_{n,o}^{\bullet} \circ \Phi_{m,n}^{\bullet}$. This yields a morphism to a direct system $T^{\bullet} \rightarrow (T^{\bullet} \langle m \rangle, \Phi_{m,n}^{\bullet})_{m \leq n \in \mathbb{N}_0}$. This is where the existence of countable direct limits $\varinjlim_{\mathcal{D}, k \in \mathbb{N}_0}$ in the codomain category $\mathcal{D}$ comes into play. For every $n \in \mathbb{Z}$ we define the $n^{\text{th}}$ term of the satellite functor construction as
\[ \widehat{T}^n(M) := \varinjlim_{k \in \mathbb{N}_0} T^n \langle k \rangle(M) \, . \]
Unravelling the definitions and noting that $\lbrace k \in \mathbb{N}_0 \mid k \geq n \rbrace$ is cofinal in $\mathbb{N}_0$, we see that
\begin{equation}\label{eq:expliciterms}
\widehat{T}^n(M) =\!\!\!\! \varinjlim_{k \in \mathbb{N}_0, k \geq n} \big(S^{n-k}T^k(M), \Phi_{k, k+1}^n(M) \big) = \varinjlim_{k \in \mathbb{N}_0} \big( S^{-k}T^{n+k}(M), \Phi_{n+k, n+k+1}^n(M) \big) \, .
\end{equation}
One obtains the connecting homomorphisms of the above satellite functor construction as follows. By Axiom~\ref{axm:les}, every short exact sequence $0 \rightarrow A \rightarrow B \rightarrow C \rightarrow 0$ in the category $\mathcal{C}$ yields a long exact sequence in $T^{\bullet}$ and $T^{\bullet} \langle m \rangle$ for any $m \in \mathbb{N}_0$. We obtain a long exact sequence in $\widehat{T}^{\bullet} = \varinjlim_{\mathcal{D}, m \in \mathbb{N}_0} T^{\bullet} \langle m \rangle$ and in particular, a connecting homomorphism $\widehat{\delta}^{\bullet}$ only if the direct limit functor $\varinjlim_{\mathcal{D}, m \in \mathbb{N}_0}$ is exact. This proves the following generalisation of~\cite[Theorem~2.2]{mis94}.

\begin{thm}\label{thm:satellitemislin}
Let $\mathcal{C}$ be an abelian category with enough projective object and let $\mathcal{D}$ be an abelian category in which all direct limits exist and are exact. Let $(T^{\bullet}, \delta^{\bullet}): \mathcal{C} \rightarrow \mathcal{D}$ be any cohomological functor.  Then
\begin{itemize}
\item a Mislin completion $(\widehat{T}^{\bullet}, \widehat{\delta}^{\bullet})$ of $(T^{\bullet}, \delta^{\bullet})$ exists and
\item can be defined via the the satellite functor construction.
\end{itemize}
\end{thm}

\begin{notn}
For the rest of the paper we assume that $\mathcal{C}$ is an abelian category with enough projective object and that $\mathcal{D}$ is an abelian category in which all direct limits exist and are exact.
\end{notn}

Let us provide an explicit formula for the connecting homomorphisms of the satellite functor construction. For this, note that one can extend the concept of left satellite functors from~\cite[Corollary~5.3]{car56} to natural transformations. Let $\varphi: F \rightarrow G$ be a natural transformation between two functors $F, G: \mathcal{C} \rightarrow \mathcal{D}$. One can define the natural transformation $S^{-1}\varphi(C): S^{-1}F(C) \rightarrow S^{-1}G(C)$ in the same manner as we have defined induced morphisms $S^{-1}T^n(f): S^{-1}T^n(C') \rightarrow S^{-1}T^n(C)$. Write $S^0{\varphi} := \varphi$ and
\[ S^{-k}\varphi := S^{-1}(S^{-k+1}\varphi): S^{-k}F \rightarrow S^{-k}G \]
for any $k \in \mathbb{N}$. Recall from Section~\ref{sec:constructions} that connecting homomorphisms are natural transformation that factor as
\[ \delta^n: T^n \xrightarrow{\underline{\delta}^n} S^{-1}T^{n+1} \xrightarrow{\varepsilon^{-1}} T^{n+1} \, . \]

\begin{thm}\label{thm:satelliteconnmorphisms}
For a short exact sequence $0 \rightarrow A \rightarrow B \rightarrow C \rightarrow 0$ in $\mathcal{C}$ and $k \in \mathbb{N}_0$ denote by $\varepsilon^{-k}: S^{-k}T^{n+k}(C) \rightarrow S^{-k+1}T^{n+k}(A)$ the associated connecting homomorphism. Then the $n^{\text{th}}$ connecting homomorphism of the satellite functor construction is given by
\[ \widehat{\delta}^n = \varinjlim_{k \in \mathbb{N}} \varepsilon^{-k}: \widehat{T}^n(C) \rightarrow \widehat{T}^{n+1}(A) \]
where the morphisms $\varepsilon^{-k}$ are connected over $S^{-k}\underline{\delta}^{n+k}: S^{-k}T^{n+k} \rightarrow S^{-k-1}T^{n+k+1}$. In particular, $\underline{\delta}^n: T^n \rightarrow S^{-1}T^{n+1}$ is a natural transformation for any $n \in \mathbb{Z}$.
\end{thm}

\begin{proof}
The proof is based on determining all terms of each morphism of the cohomological functors $\Phi_{k, k+1}^{\bullet}: T^{\bullet}\langle k \rangle \rightarrow T^{\bullet}\langle k+1 \rangle$ that give rise to the satellite functor construction. We already know that $\Phi_{k, k+1}^n = \mathrm{id}_{\mathcal{D}}$ for $k < n$. To describe $\Phi_{k, k+1}^k$, consider a short exact sequence $0 \rightarrow M \rightarrow P \rightarrow C \rightarrow 0$ with $P$ projective. We obtain from Diagram~\ref{diag:delta} the commutative diagram
\begin{equation}\label{diag:factorise}
\begin{tikzcd}
  {} \dots {} \arrow[r]
    &T^k(P) \arrow[r] \arrow[d, "T^k(h)"]
    &T^k(C) \arrow[r, "\delta^k"] \arrow[d, "\mathrm{id}_{T^k(C)}"]
    &T^{k+1}(M) \arrow[r] \arrow[d, "T^{k+1}(h^{\ast})"]
    &T^{k+1}(P) \arrow[r] \arrow[d, "T^{k+1}(h)"]
    &{} \dots {} \\
  {} \dots {} \arrow[r]
    &T^k(B) \arrow[r]
    &T^k(C) \arrow[r, "\delta^k"]
    &T^{k+1}(A) \arrow[r]
    &T^{k+1}(B) \arrow[r]
    &{} \dots {}
\end{tikzcd}
\end{equation}
Recalling that $S^{-1}T^{k+1}(C)$ is a kernel and $\varepsilon^{-1}: S^{-1}T^{k+1}(C) \rightarrow T^{k+1}(M)$ its canonical monomorphism, we can write the middle square as
\begin{center}
\begin{tikzcd}
  T^k(C) \arrow[r, "\underline{\delta}^k"] \arrow[d, "\mathrm{id}_{T^k(C)}"]
    &S^{-1}T^{k+1}(C) \arrow[r, "\varepsilon^{-1}"]
    &T^{k+1}(M) \arrow[d, "T^{k+1}(h^{\ast})"] \\
  T^k(C) \arrow[rr, "\delta^n"]
    &{}
    &T^{k+1}(A)
\end{tikzcd}
\end{center}
By definition, $\delta^{k}\langle k+1 \rangle = \varepsilon^{-1}: S^{-1}T^{k+1}(C) \xrightarrow{\varepsilon^{-1}} T^{k+1}(M) \xrightarrow{T^{k+1}(h^{\ast})} T^{k+1}(A)$. Tilting the above diagram, we obtain the desired commutative diagram
\begin{equation}\label{diag:twoconnmorphisms}
\begin{tikzcd}
  T^k(C) \arrow[r, "\delta^k"] \arrow[d, "\underline{\delta}^k"]
    &T^{k+1}(A) \arrow[d, "\mathrm{id}_{T^{k+1}(A)}"] \\
  S^{-1}T^{k+1}(C) \arrow[r, "\varepsilon^{-1}"]
    &T^{k+1}(A)
\end{tikzcd}
\end{equation}
By the proof of Lemma~\ref{lem:uniqueness} found in~\cite[p.~46--47]{car56}, $\underline{\delta}^k = \Phi_{k, k+1}^k: T^k \rightarrow S^{-1}T^{k+1}$ is a natural transformation. Thus, $\Phi_{k, k+1}^n = S^{n-k}\underline{\delta}^k: S^{n-k}T^k \rightarrow S^{n-k+1}T^{k+1}$ for any $n < k$. In particular, the square
\begin{equation}\label{diag:takesatellitetrsfs}
\begin{tikzcd}
    S^{n-k}T^k(C) \arrow[r, "\varepsilon^{n-k}"] \arrow[d, "S^{n-k}\underline{\delta}^k"] & S^{-k+1}T^n(A) \arrow[d, "S^{n-k+1}\underline{\delta}^k"] \\
    S^{n-k-1}T^{k+1}(C) \arrow[r, "\varepsilon^{n-k-1}"] & S^{n-k}T^{k+1}(A)
\end{tikzcd}
\end{equation}
commutes. In the direct limit, these squares give rise to the connecting homomorphism $\widehat{\delta}^n$ and the formula in the lemma follows from reindexing as in Equation~\ref{eq:expliciterms}.
\end{proof}

Lastly, we provide an explicit formula for the induced morphisms in the satellite functor construction.

\begin{thm}\label{thm:inducedsatellitemorphs}
Let $f: A \rightarrow B$ be a morphism in $\mathcal{C}$. Then for any $n \in \mathbb{Z}$ the commuting squares
\begin{center}
\begin{tikzcd}
    S^{-k}T^{n+k}(A) \arrow[rrr, "S^{-k}T^{n+k}(f)"] \arrow[d, "S^{-k}\underline{\delta}^{n+k}"] & & & S^{-k}T^{n+k}(B) \arrow[d, "S^{-k}\underline{\delta}^{n+k}"] \\
    S^{-k-1}T^{n+k+1}(A) \arrow[rrr, "S^{-k-1}T^{n+k+1}(f)"] & & & S^{-k-1}T^{n+k+1}(B)
\end{tikzcd}
\end{center}
give rise to
\[ \widehat{T}^n(f) = \varinjlim_{\mathcal{D}, k \in \mathbb{N}_0} S^{-k}T^{n+k}(f): \widehat{T}(A) \rightarrow \widehat{T}(B) \, , \]
the induced morphism in the satellite functor construction.
\end{thm}

\begin{proof}
This follows from Equation~\ref{eq:expliciterms} and the description of the cohomological functors $\Phi_{n+k, n+k+1}^{\bullet}$ found in the proof of Theorem~\ref{thm:satelliteconnmorphisms}.
\end{proof}

\section{The resolution and the satellite functor construction}\label{sec:resols}

The main goal of this section is to prove that the resolution construction gives rise to Mislin completions of cohomological functors. For this purpose, we develop explicit formulae for its induced morphisms and connecting homomorphisms of the resolution construction. This is done by developing isomorphisms between the terms of the satellite functor construction and the ones of the resolution construction. As a consequence of these formulae, we demonstrate that there are countably many distinct constructions of Mislin  completions.

\subsection{Termwise isomorphisms}

To set up notation for the resolution construction, let $M \in \mathrm{obj}(\mathcal{C})$, $(M_n)_{n \in \mathbb{N}_0}$ be a projective resolution of $M$ and $\widetilde{M}_k$ be the $k^{\text{th}}$ syzygy of $M_{\bullet}$ for $k \in \mathbb{N}_0$. For $n \in \mathbb{Z}$ we denote the $n^{\text{th}}$ term of the satellite functor construction specifically by $\widehat{T}^n(M)$. If $\delta^{n+k}: T^{n+k}(\widetilde{M}_k) \rightarrow T^{n+k+1}(\widetilde{M}_{k+1})$ are taken as in Section~\ref{sec:constructions}, we write the $n^{\text{th}}$ term of the resolution construction specifically as $T_{Res}^n(M) := \varinjlim_{k \in \mathbb{N}_0} \big(T^{n+k}(\widetilde{M}_k), \delta^{n+k}\big)$. A priori, the terms of the resolution construction depend on a choice of projective resolution. This issue is resolved by demonstrating

\begin{lem}\label{lem:resolisoms}
For any $n \in \mathbb{Z}$ there is an isomorphism $\omega_n(M): \widehat{T}^n(M) \rightarrow T_{Res}^n(M)$ in $\mathcal{D}$.
\end{lem}

A proof of this lemma can be extracted from~\cite[p.~299]{mis94} for modules over rings and can be found in~\cite[p.~20]{guo23} for relative Ext-functors. However, the resulting isomorphisms are not in an explicit form in both cases. In order to develop explicit formulae for the induced morphisms and connecting homomorphisms of the resolution construction, we require an explicit formula for such an isomorphism. Hence, we remedy this issue by providing our own proof the above lemma.

\begin{proof}
The proof consists of a two step endeavor where we introduce the relevant notation throughout its course. In the first step, we construct an isomorphism
\[ \eta_n(M): \varinjlim_{k \in \mathbb{N}} (S^{-1}T^{n+k}(\widetilde{M}_{k-1}), \Sigma^{-1}\delta^{n+k}) \rightarrow T_{Res}^n(M) \]
and in the second step we extend it to an isomorphism $\omega_n(M): \widehat{T}^n(M) \rightarrow T_{Res}^n(M)$. \\

{\textbf{Step 1}} Defining an isomorphism $\eta_n(M): \varinjlim_{k \in \mathbb{N}} (S^{-1}T^{n+k}(\widetilde{M}_{k-1}), \Sigma^{-1}\delta^{n+k}) \rightarrow T_{Res}^n(M)$ \\

Denote by $\underline{\delta}^{n+k}: T^{n+k}(\widetilde{M}_k) \rightarrow S^{-1}T^{n+k+1}(\widetilde{M}_k)$ the morphism associated to the short exact sequence $0 \rightarrow \widetilde{M}_{k+1} \rightarrow M_k \rightarrow \widetilde{M}_k \rightarrow 0$ and by $\varepsilon^{-1}: S^{-1}T^{n+k}(\widetilde{M}_{k-1}) \rightarrow T^{n+k}(\widetilde{M}_k)$ the canonical morphism from the kernel. Define the morphism
\[ \Sigma^{-1}\delta^{n+k} := \underline{\delta}^{n+k} \circ \varepsilon^{-1}: \; \; S^{-1}T^{n+k}(\widetilde{M}_{k-1}) \rightarrow T^{n+k}(\widetilde{M}_k) \rightarrow S^{-1}T^{n+k+1}(\widetilde{M}_k) \, . \]
This morphism together with the cokernel of $\varepsilon^{-1}$ denoted by $q_{n+k}: T^{n+k}(\widetilde{M}_k) \rightarrow N_{n+k}$ extend Diagram~\ref{diag:twoconnmorphisms} in a commutative manner as
\begin{equation}\label{diag:cokernels}
\begin{tikzcd}
  0 \arrow[r]
    & S^{-1}T^{n+k}(\widetilde{M}_{k-1}) \arrow[r, "\varepsilon^{-1}"] \arrow[d, "\Sigma^{-1}\delta^{n+k}" near start]
    & T^{n+k}(\widetilde{M}_k) \arrow[r, "q_{n+k}"] \arrow[dl, "\underline{\delta}^{n+k}" near start] \arrow[d, "\delta^{n+k}" near start]
    & N_{n+k} \arrow[r] \arrow[d, "m_{n+k}"]
    & 0 \\
  0 \arrow[r]
    & S^{-1}T^{n+k+1}(\widetilde{M}_k) \arrow[r, "\varepsilon^{-1}"]
    & T^{n+k+1}(\widetilde{M}_{k+1}) \arrow[r, "q_{n+k+1}"]
    & N_{n+k+1} \arrow[r]
    & 0
\end{tikzcd}
\end{equation}
Since
\[ m_{n+k} \circ q_{n+k} = q_{n+k+1} \circ \varepsilon^{-1} \circ \underline{\delta}^{n+k}: T^{n+k}(\widetilde{M}_k) \rightarrow N_{n+k+1} \]
is the zero morphism, we deduce that $m_{n+k} = 0$. Note that these diagrams form a direct system of short exact sequences where countable direct limits in $\mathcal{D}$ are exact. Thus, we obtain the required isomorphism
\[ \eta_n(M) = \varinjlim \varepsilon^{-1}: \varinjlim_{k \in \mathbb{N}} (S^{-1}T^{n+k}(\widetilde{M}_{k-1}), \Sigma^{-1}\delta^{n+k}) \rightarrow T_{Res}^n(M) \, . \]

{\textbf{Step 2}} Defining an isomorphism $\omega_n(M): \widehat{T}^n(M) \rightarrow T_{Res}^n(M)$ \\

We connect the terms of the direct system $(S^{-k}T^{n+k}(M), S^{-k}\underline{\delta}^{n+k})_{k \in \mathbb{N}_0}$ to the ones of $(T^{n+k}(\widetilde{M}_k), \delta^{n+k})_{k \in \mathbb{N}_0}$ to construct this isomorphism. This requires a larger commutative diagram for which we introduce the necessary morphisms. As for left satellite functors, we define $\Sigma^0{\delta^{n+k}} := \delta^{n+k}$ and
\begin{align}
\Sigma^{-l}\delta^{n+k+l} &:= S^{-l+1}\underline{\delta}^{n+k+l} \circ \varepsilon^{-l}: \label{eq:sigmafunctors}  \\
&S^{-l}T^{n+k+l}(\widetilde{M}_k) \rightarrow S^{-l+1}T^{n+k+l}(\widetilde{M}_{k+1}) \rightarrow S^{-l}T^{n+k+l+1}(\widetilde{M}_{k+1}) \nonumber
\end{align}
for any $l \in \mathbb{N}$. We can extend Diagram~\ref{diag:takesatellitetrsfs} to another commutative diagram of the form
\begin{equation}\label{diag:commtriangles}
\begin{tikzcd}
    S^{-l}T^{n+k+l}(\widetilde{M}_k) \arrow[r, "\varepsilon^{-l}"] \arrow[dr, "\Sigma^{-l}\delta^{n+k+l}"] \arrow[d, "S^{-l}\underline{\delta}^{n+k+l}" near end] & S^{-l+1}T^{n+k+l}(\widetilde{M}_{k+1}) \arrow[d, "S^{-l+1}\underline{\delta}^{n+k+l}"] \\
    S^{-l-1}T^{n+k+l+1}(\widetilde{M}_k) \arrow[r, "\varepsilon^{-l-1}"] & S^{-l}T^{n+k+l+1}(\widetilde{M}_{k+1})
\end{tikzcd}
\end{equation}
Because left satellite functors vanish on projective objects such as $M_k$, the morphism $\varepsilon^{-l}: S^{-l}T^{n+k+l}(\widetilde{M}_k) \rightarrow S^{-l+1}T^{n+k+l}(\widetilde{M}_{k+1})$ is an isomorphism. Define the isomorphisms
\begin{align}
\widetilde{o}^k := \varepsilon^{-3} \circ {} \dots {} \circ \varepsilon^{-k+1} \circ \varepsilon^{-k}: \; \; {} &S^{-k}T^{n+k}(M) \rightarrow S^{-2}T^{n+k}(\widetilde{M}_{k-2}) \nonumber \\
\text{and} \quad o^{k} := \varepsilon^{-2} \circ \widetilde{o}^k: \; \; {} &S^{-k}T^{n+k}(M) \rightarrow S^{-1}T^{n+k}(\widetilde{M}_{k-1}) \, . \label{eq:themapsok}
\end{align}
Using the above morphisms together with Diagram~\ref{diag:twoconnmorphisms}, Diagram~\ref{diag:takesatellitetrsfs} and Diagram~\ref{diag:commtriangles}, we can write the commutative diagram

\begin{equation}\label{diag:finalisom}
\begin{tikzcd}
  {}
    & {}
    & {}
    & {\scriptstyle T^n(M)} \arrow[d, "{\scriptscriptstyle \delta^{n}}" near start] \arrow[dl, "{\scriptscriptstyle \underline{\delta}^{n}}" near start] \\
  {}
    & {}
    & {\scriptstyle S^{-1}T^{n+1}(M)} \arrow[r, "{\scriptscriptstyle \varepsilon^{-1}}"] \arrow[d, "{\scriptscriptstyle \Sigma^{-1}\delta^{n+1}}" near start] \arrow[dl, "{\scriptscriptstyle S^{-1}\underline{\delta}^{n+1}}" near start]
    & {\scriptstyle T^{n+1}(\widetilde{M}_1)} \arrow[d, "{\scriptscriptstyle \delta^{n+1}}" near start] \arrow[dl, "{\scriptscriptstyle \underline{\delta}^{n+1}}" near start] \\
  {}
    & {\scriptstyle S^{-2}T^{n+2}(M)} \arrow[r, "{\scriptscriptstyle \varepsilon^{-2}}"] \arrow[d, "{\scriptscriptstyle \Sigma^{-2}\delta^{n+2}}" near start] \arrow[dl, "{\scriptscriptstyle S^{-2}\underline{\delta}^{n+2}}" near start]
    & {\scriptstyle S^{-1}T^{n+2}(\widetilde{M}_1)} \arrow[r, "{\scriptscriptstyle \varepsilon^{-1}}"] \arrow[d, "{\scriptscriptstyle \Sigma^{-1}\delta^{n+2}}" near start] \arrow[dl, "{\scriptscriptstyle S^{-1}\underline{\delta}^{n+2}}" near start]
    & {\scriptstyle T^{n+2}(\widetilde{M}_2)} \arrow[d, "{\scriptscriptstyle \delta^{n+2}}" near start] \arrow[dl, "{\scriptscriptstyle \underline{\delta}^{n+2}}" near start] \\
  {\scriptstyle S^{-3}T^{n+3}(M)} \arrow[r, "{\scriptscriptstyle \widetilde{o}^3}"] \arrow[d, "{\scriptscriptstyle S^{-3}\underline{\delta}^{n+3}}" near start]
    & {\scriptstyle S^{-2}T^{n+3}(\widetilde{M}_1)} \arrow[r, "{\scriptscriptstyle \varepsilon^{-2}}"] \arrow[d, "{\scriptscriptstyle \Sigma^{-2}\delta^{n+3}}" near start]
    & {\scriptstyle S^{-1}T^{n+3}(\widetilde{M}_2)} \arrow[r, "{\scriptscriptstyle \varepsilon^{-1}}"] \arrow[d, "{\scriptscriptstyle \Sigma^{-1}\delta^{n+3}}" near start] \arrow[dl, "{\scriptscriptstyle S^{-1}\underline{\delta}^{n+3}}" near start]
    & {\scriptstyle T^{n+3}(\widetilde{M}_3)} \arrow[d, "{\scriptscriptstyle \delta^{n+3}}" near start] \arrow[dl, "{\scriptscriptstyle \underline{\delta}^{n+3}}" near start] \\
  {\scriptstyle S^{-4}T^{n+4}(M)} \arrow[r, "{\scriptscriptstyle \widetilde{o}^4}"] \arrow[d]
    & {\scriptstyle S^{-2}T^{n+4}(\widetilde{M}_2)} \arrow[r, "{\scriptscriptstyle \varepsilon^{-2}}"] \arrow[d]
    & {\scriptstyle S^{-1}T^{n+4}(\widetilde{M}_3)} \arrow[r, "{\scriptscriptstyle \varepsilon^{-1}}"] \arrow[d]
    & {\scriptstyle T^{n+4}(\widetilde{M}_4)} \arrow[d] \\
  {} {\scriptstyle \vdots} {} \arrow[r] \arrow[d]
    & {} {\scriptstyle \vdots} {} \arrow[r] \arrow[d]
    & {} {\scriptstyle \vdots} {} \arrow[r] \arrow[d]
    & {} {\scriptstyle \vdots} {} \arrow[d] \\
  {\scriptstyle S^{-k}T^{n+k}(M)} \arrow[r, "{\scriptscriptstyle \widetilde{o}^k}"] \arrow[rr, bend left = 15, crossing over, "{\scriptscriptstyle o^k}" near end] \arrow[d, "{\scriptscriptstyle S^{-k}\underline{\delta}^{n+k}}" near start]
    & {\scriptstyle S^{-2}T^{n+k}(\widetilde{M}_{k-2})} \arrow[r, "{\scriptscriptstyle \varepsilon^{-2}}"] \arrow[d, "{\scriptscriptstyle \Sigma^{-2}\delta^{n+k}}" near start]
    & {\scriptstyle S^{-1}T^{n+k}(\widetilde{M}_{k-1})} \arrow[r, "{\scriptscriptstyle \varepsilon^{-1}}"] \arrow[d, "{\scriptscriptstyle \Sigma^{-1}\delta^{n+k}}" near start] \arrow[dl, "{\scriptscriptstyle S^{-1}\underline{\delta}^{n+k}}" near start]
    & {\scriptstyle T^{n+k}(\widetilde{M}_k)} \arrow[d, "{\scriptscriptstyle \delta^{n+k}}" near start] \arrow[dl, "{\scriptscriptstyle \underline{\delta}^{n+k}}" near start] \\
  {\scriptstyle S^{-k-1}T^{n+k+1}(M)} \arrow[r, "{\scriptscriptstyle \widetilde{o}^{k+1}}"] \arrow[d]
    & {\scriptstyle S^{-2}T^{n+k+1}(\widetilde{M}_{k-1})} \arrow[r, "{\scriptscriptstyle \varepsilon^{-2}}"] \arrow[d]
    & {\scriptstyle S^{-1}T^{n+k+1}(\widetilde{M}_k)} \arrow[r, "{\scriptscriptstyle \varepsilon^{-1}}"] \arrow[d]
    & {\scriptstyle T^{n+k+1}(\widetilde{M}_{k+1})} \arrow[d] \\
  {} {\scriptstyle \vdots} {}
    & {} {\scriptstyle \vdots} {}
    & {} {\scriptstyle \vdots} {}
    & {} {\scriptstyle \vdots} {}
\end{tikzcd}
\end{equation}

We see that the right most column corresponds to the direct system $(T^{n+k}(\widetilde{M}_k), \delta^{n+k})_{k \in \mathbb{N}_0}$ while the top diagonal together with the left most column forms the direct system $(S^{-k}T^{n+k}(M), S^{-k}\underline{\delta}^{n+k})_{k \in \mathbb{N}_0}$. The right most squares in the diagram recover the isomorphism
\[ \eta_n(M) := \varinjlim \varepsilon^{-1}: \varinjlim_{k \in \mathbb{N}} (S^{-1}T^{n+k}(\widetilde{M}_{k-1}), \Sigma^{-1}\delta^{n+k}) \rightarrow T_{Res}^n(M) \, . \]
For $k \geq 3$, all the remaining squares yield together squares of the form
\begin{center}
\begin{tikzcd}
  S^{-k}T^{n+k}(M) \arrow[r, "o^k"] \arrow[d, "S^{-k}\underline{\delta}^{n+k}"]
    & S^{-1}T^{n+k}(\widetilde{M}_{k-1}) \arrow[d, "\Sigma^{-1}\delta^{n+k}"] \\
  S^{-k-1}T^{n+k+1}(M) \arrow[r, "o^{k+1}"]
    & S^{-1}T^{n+k+1}(\widetilde{M}_k)
\end{tikzcd}
\end{center}
As the morphisms $o^k$ are isomorphisms and countable direct limits in $\mathcal{D}$ are exact, this gives rise to the isomorphism
\[ \varinjlim o^k: \widehat{T}^n(M) \rightarrow \varinjlim_{k \in \mathbb{N}} (S^{-1}T^{n+k}(\widetilde{M}_{k-1}), \Sigma^{-1}\delta^{n+k}) \, . \]
Therefore,
\[ \omega_n(M) := (\varinjlim o^k) \circ \eta_n(M) = \varinjlim (o^k \circ \varepsilon^{-1}): \widehat{T}^n(M) \rightarrow T_{Res}^n(M) \]
is our desired isomorphism. Since the morphisms $o^k \circ \varepsilon^{-1}$ span the entire width of Diagram~\ref{diag:finalisom}, every single morphism was needed to construct this isomorphism.
\end{proof}

\subsection{Induced morphisms and connecting homomorphisms}

We present an explicit formula for the induced morphisms of the resolution construction.

\begin{defn}\label{defn:resolindmorph}
Let $f: M \rightarrow N$ be a morphism in $\mathcal{C}$ and $M_{\bullet}$, $N_{\bullet}$ be projective resolutions of $M$, $N$. We define a sequence $(\widetilde{f}_k: \widetilde{M}_k \rightarrow \widetilde{N}_k)_{k \in \mathbb{N}_0}$ inductively as follows. We set $\widetilde{f}_0 := f$ and impose that $\widetilde{f}_{k+1}$ is a lift $\widetilde{f}_k$, meaning that diagram
\begin{equation}\label{diag:complexlift}
\begin{tikzcd}
    0 \arrow[r] & \widetilde{M}_{k+1} \arrow[r] \arrow[d, "\widetilde{f}_{k+1}"] & M_k \arrow[r] \arrow[d, "f_k"] & \widetilde{M}_k \arrow[r] \arrow[d, "\widetilde{f}_k"] & 0 \\
    0 \arrow[r] & \widetilde{N}_{k+1} \arrow[r] & N_k \arrow[r] & \widetilde{N}_k \arrow[r] & 0
\end{tikzcd}
\end{equation}
commutes. Then the direct limit of the commutative diagrams
\begin{center}
\begin{tikzcd}
    T^{n+k}(\widetilde{M}_k) \arrow[rr, "T^{n+k}(\widetilde{f}_k)"] \arrow[d, "\delta^{n+k}"] & & T^{n+k}(\widetilde{N}_k) \arrow[d, "\delta^{n+k}"] \\
    T^{n+k+1}(\widetilde{M}_{k+1}) \arrow[rr, "T^{n+k+1}(\widetilde{f}_{k+1})"] & & T^{n+k+1}(\widetilde{N}_{k+1})
\end{tikzcd}
\end{center}
written as $T_{Res}^n(f) := \varinjlim_{k \in \mathbb{N}_0} T^{n+k}(\widetilde{f}_k)$ is the induced morphism in the resolution construction.
\end{defn}

Here, there is a similar issue as in the definition of the terms of the resolution construction. A priori, the morphism $T_{Res}^n(f): T_{Res}^n(M) \rightarrow T_{Res}^n(N)$ depends on choices of projective resolutions of $M$ and $N$. This can be resolved by demonstrating

\begin{lem}\label{lem:omeganatural}
The morphisms $\omega_n: \widehat{T}^n \rightarrow T_{Res}^n$  form a natural isomorphism. In other words, the square
\begin{center}
\begin{tikzcd}
    \widehat{T}^n(M) \arrow[r, "\omega_n(M)"] \arrow[d, "\widehat{T}^n(f)"] & T_{Res}^n(M) \arrow[d, "T_{Res}^n(f)"] \\
    \widehat{T}^n(N) \arrow[r, "\omega_n(N)"] & T_{Res}^n(N)
\end{tikzcd}
\end{center}
commutes for any morphism $f: M \rightarrow N$ in $\mathcal{C}$.
\end{lem}

\begin{proof}
This follows from the fact that the isomorphisms $\omega_n$ are constructed by the morphisms in Diagram~\ref{diag:finalisom}, which are all natural transformations.
\end{proof}

Hereby we provide an explicit formula for the connecting homomorphisms of the satellite functor construction.

\begin{defn}\label{defn:resolconnmorph}
Let $0 \rightarrow A \xrightarrow{f} B \xrightarrow{g} C \rightarrow 0$ be short exact sequence in $\mathcal{C}$. We construct a sequence of short exact sequences $(0 \rightarrow \widetilde{A}_k \xrightarrow{\widetilde{f}_k} \widetilde{B}_k \xrightarrow{\widetilde{g}_k} \widetilde{C}_k \rightarrow 0)_{k \in \mathbb{N}_0}$ in $\mathcal{C}$ inductively as follows. Set $\widetilde{D}_0 := D$ for $D \in \lbrace A, B, C \rbrace$ and $\widetilde{h}_0 := h$ for $h \in \lbrace f, g \rbrace$. We impose that the diagram as in the Horseshoe Lemma found in~\cite[p.~37]{wei94}
\begin{equation}\label{diag:horseshoe}
\begin{tikzcd}[row sep = scriptsize]
  {}
    & 0 \arrow[d]
    & 0 \arrow[d]
    & 0 \arrow[d]
    & {} \\
  0 \arrow[r]
    & \widetilde{A}_{k+1} \arrow[r, "\widetilde{f}_{k+1}"] \arrow[d]
    & \widetilde{B}_{k+1} \arrow[r, "\widetilde{g}_{k+1}"] \arrow[d]
    & \widetilde{C}_{k+1} \arrow[r] \arrow[d]
    & 0 \\
  0 \arrow[r]
    & A_k \arrow[r, "f_k"] \arrow[d]
    & B_k \arrow[r, "g_k"] \arrow[d]
    & C_k \arrow[r] \arrow[d]
    & 0 \\
  0 \arrow[r]
    & \widetilde{A}_k \arrow[r, "\widetilde{f}_k"] \arrow[d]
    & \widetilde{B}_k \arrow[r, "\widetilde{g}_k"] \arrow[d]
    & \widetilde{C}_k \arrow[r] \arrow[d]
    & 0 \\
  {}
    & 0
    & 0
    & 0
    & {}
\end{tikzcd}
\end{equation}
commutes. This means that $A_k$, $B_k$ and $C_k$ are projective and that all rows and columns are exact. Because Proposition~III.4.1 of~\cite{car56} also pertains to abelian categories, the square
\begin{equation}\label{diag:anticomm}
\begin{tikzcd}[row sep = scriptsize]
  T^{n+k}(\widetilde{C}_k) \arrow[r, "\delta^{n+k}"] \arrow[d, "\delta^{n+k}"]
    & T^{n+k+1}(\widetilde{A}_k) \arrow[d, "\delta^{n+k+1}"] \\
  T^{n+k+1}(\widetilde{C}_{k+1}) \arrow[r, "\delta^{n+k+1}"]
    & T^{n+k+2}(\widetilde{A}_{k+1})
\end{tikzcd}
\end{equation}
anticommutes for any $n \in \mathbb{Z}$ and $k \in \mathbb{N}_0$. This means that the composite morphisms in the abelian group $\mathrm{Hom}_{\mathcal{D}}(T^{n+k}(\widetilde{C}_k), T^{n+k+2}(\widetilde{A}_{k+1}))$ have opposite signs. Then the direct limit of these anticommutatitve squares
\[ \widetilde{\delta}^n := \varinjlim_{k \in \mathbb{N}_0} (-1)^k \delta^{n+k}: T_{Res}^n(C) \rightarrow T_{Res}^{n+1}(A) \]
is the connecting homomorphism of the resolution construction.
\end{defn}

A priori, this definition depends on specific choices of projective resolutions of the terms $A$, $B$ and $C$. This issue can be remedied by demonstrating

\begin{lem}\label{lem:commwithconn}
Denote by $\widehat{\delta}^{\bullet}$ the connecting homomorphism from the satellite functor construction $\widehat{T}^{\bullet}$. Then for any short exact sequence $0 \rightarrow A \rightarrow B \rightarrow C \rightarrow 0$ in $\mathcal{C}$ and any $n \in \mathbb{Z}$ the diagram
\begin{center}
\begin{tikzcd}[row sep = small]
  \widehat{T}^n(C) \arrow[r, "\omega_n"] \arrow[d, "\widehat{\delta}^n"]
  & T_{Res}^n(C) \arrow[d, "\widetilde{\delta}^n"] \\
  \widehat{T}^{n+1}(A) \arrow[r, "\omega_{n+1}"]
  & T_{Res}^{n+1}(A)
\end{tikzcd}
\end{center}
commutes.
\end{lem}

\begin{proof}
The proof is based on the diagram found on the next page. First, we show that it commutes. The front side is a copy of Diagram~\ref{diag:finalisom} as well as the back side up to its commuting right most column. The right hand side of the diagram gives rise to $\widetilde{\delta}^n: T_{Res}^n(C) \rightarrow T_{Res}^{n+1}(A)$ in the direct limit and the top and left hand side give rise to $\widehat{\delta}^n: \widehat{T}^n(C) \rightarrow \widehat{T}^{n+1}(A)$. In particular, all squares and triangles from these sides commute. Using the cohomological functors defined in Equation~\ref{eq:directsyscohomfunctor}, we see that the horizontal squares commute by~\cite[Proposition~III.4.1]{car56}. As in Diagram~\ref{diag:finalisom} we can insert commuting triangles of the form
\begin{center}
\begin{tikzcd}
  {} & S^{-l}T^{n+k}(\widetilde{M}_{k-l}) \arrow[d, "\Sigma^{-l}\delta^{n+k}"] \arrow[dl, "S^{-l}\underline{\delta}^{n+k}" near start] \\
  S^{-l-1}T^{n+k+1}(\widetilde{M}_{k-l}) \arrow[r, "\varepsilon^{-l-1}"]
  & S^{-l}T^{n+k+1}(\widetilde{M}_{k-l+1})
\end{tikzcd}
\end{center}
into the front and back side. We conclude by this and Diagram~\ref{diag:takesatellitetrsfs} that the vertical squares running from back to front also commute. In particular, Diagram~\ref{diag:largestdiag} commutes. \\

To complete the proof, we note that the right most column on the back side yields the identity morphism in the direct limit. Thus, the entire back side yields the morphism $\omega_{n+1}: \widetilde{T}^{n+1}(A) \rightarrow T_{Res}^{n+1}(A)$ while the front side yields $\omega_n: \widetilde{T}^n(A) \rightarrow T_{Res}^n(A)$. As mentioned before, the right hand side yields $\widetilde{\delta}^n: T_{Res}^n(C) \rightarrow T_{Res}^{n+1}(A)$ and left hand side $\widehat{\delta}^n: \widehat{T}^n(C) \rightarrow \widehat{T}^{n+1}(A)$. Therefore, the square in the statement of the lemma commutes.
\end{proof}

By Lemma~\ref{lem:omeganatural} and Lemma~\ref{lem:commwithconn}, we conclude that the resolution construction $(T_{Res}^{\bullet}, \widetilde{\delta}^{\bullet})$ with the formula for induced morphisms from Definition~\ref{defn:resolindmorph} and for connecting homomorphisms from Definition~\ref{defn:resolconnmorph} forms a Mislin completion of $(T^{\bullet}, \delta^{\bullet})$.

\begin{thm}\label{thm:resolsatellites}
The resolution construction $(T_{Res}^{\bullet}, \widetilde{\delta}^{\bullet})$ forms a Mislin completion of $(T^{\bullet}, \delta^{\bullet})$. More specifically, $(T_{Res}^{\bullet}, \widetilde{\delta}^{\bullet})$ forms a cohomological functor and $\omega_{\bullet}: (\widehat{T}^{\bullet}, \widehat{\delta}^{\bullet}) \rightarrow (T_{Res}^{\bullet}, \widetilde{\delta}^{\bullet})$ is an isomorphism of cohomological functors.
\end{thm}

\subsection{Countably many distinct constructions of Mislin completions}

In the remainder of this section, we demonstrate that there are countably many distinct constructions of Mislin completions. In particular, we shall clarify what we mean by two constructions of Mislin completions being distinct. First, we present a collection of constructions whose terms we define as follows.

\begin{samepage}
{\vspace{-15mm}
\begin{flushleft}
\includegraphics[height = 245mm]{giantdiagram.pdf}
\end{flushleft} }
{\vspace{-150mm}
\begin{equation}\label{diag:largestdiag}
{} \quad {}
\end{equation}
${} \quad {}$ \vspace{150mm}}
\end{samepage}

\begin{defn}
Let $a = (a_k)_{k \in \mathbb{N}} \in \lbrace 0, 1 \rbrace^{\mathbb{N}}$ be a sequence taking values in $0$ and $1$. Define
\[ (P(a)_k)_{k \in \mathbb{N}_0} := \begin{cases} 0 &\text{if } k = 0 \\ \sum_{i = 1}^k a_i &\text{if } k \in \mathbb{N} \end{cases} \quad \text{and} \qquad (D(a)_k)_{k \in \mathbb{N}_0} := \begin{cases} 0 &\text{if } k = 0 \\ \sum_{i = 1}^k (1-a_i) &\text{if } k \in \mathbb{N} \end{cases} \]
For any $M \in \mathrm{obj}(\mathcal{C})$ and $n \in \mathbb{Z}$ form the direct system $\big(S^{-P(a)_k}T^{n+k}(\widetilde{M}_{D(a)_k}), \delta_a^{n+k} \big)_{k \in \mathbb{N}_0}$ where
\[ \delta_a^{n+k} := \begin{cases} {\scriptstyle \Sigma^{-P(a)_k}\delta^{n+k}: \: S^{-P(a)_k}T^{n+k}(\widetilde{M}_{D(a)_k}) \rightarrow S^{-P(a)_k}T^{n+k+1}(\widetilde{M}_{D(a)_k+1})} &{\scriptstyle \text{if } P(a)_{k+1} = P(a)_k} \\ {\scriptstyle S^{-P(a)_k}\underline{\delta}^{n+k}: \: S^{-P(a)_k}T^{n+k}(\widetilde{M}_{D(a)_k}) \rightarrow S^{-P(a)_k-1}T^{n+k+1}(\widetilde{M}_{D(a)_k})} &{\scriptstyle \text{if } P(a)_{k+1} = P(a)_k+1} \end{cases} \]
with the morphisms $\Sigma^{-P(a)_k}\delta^{n+k}$ taken as in Equation~\ref{eq:sigmafunctors}. Write
\[ T_a^n(M) := \varinjlim_{\mathcal{D}, k \in \mathbb{N}_0} S^{-P(a)_k}T^{n+k}(\widetilde{M}_{D(a)_k}) \]
for the $n^{\text{th}}$ term of the construction $T_a^{\bullet}$.
\end{defn}

As a first step in showing that these constructions form Mislin completions, we prove

\begin{lem}\label{lem:oneconstrofcompletion}
There exists an isomorphism $\omega_{a, n}: T_a^n(M) \rightarrow T_{Res}^n(M)$ for any $M \in \mathrm{obj}(\mathcal{C})$.
\end{lem}

\begin{proof}
For any $k \in \mathbb{N}_0$ we define the morphism
\[ \widehat{o}_a^k := \begin{cases} \varepsilon^{-1} \circ {} \dots {} \circ \varepsilon^{-P(a)_k}: S^{-P(a)_k}T^{n+k}(\widetilde{M}_{D(a)_k}) \rightarrow T^{n+k}(\widetilde{M}_k) &\text{if } P(a)_k \geq 1 \\ \qquad \qquad \qquad \quad \mathrm{id}: T^{n+k}(\widetilde{M}_k) \rightarrow T^{n+k}(\widetilde{M}_k) &\text{if } P(a)_k = 0 \end{cases} \]
similarly to the morphism $o^k$ from Equation~\ref{eq:themapsok}. By Diagram~\ref{diag:finalisom} we infer that the square
\begin{center}
\begin{tikzcd}
    S^{-P(a)_k}T^{n+k}(\widetilde{M}_{D(a)_k}) \arrow[r, "\widehat{o}_a^k"] \arrow[d, "\delta_a^{n+k}"] & T^{n+k}(\widetilde{M}_k) \arrow[d, "\delta^{n+k}"] \\
    S^{-P(a)_{k+1}}T^{n+k+1}(\widetilde{M}_{D(a)_{k+1}}) \arrow[r, "\widehat{o}_a^{k+1}"] & T^{n+k+1}(\widetilde{M}_{k+1})
\end{tikzcd}
\end{center}
commutes. We conclude as in the proof of Lemma~\ref{lem:resolisoms} that $\omega_{a, n} := \varinjlim_{k \in \mathbb{N}_0} \widehat{o}_a^k$ is an isomorphism.
\end{proof}

We provide explicit formulae for the induced morphisms of the constructions $T_a^{\bullet}$.

\begin{defn}\label{defn:countmanyconstrsindmorph}
Let $f: M \rightarrow N$ be a morphism in $\mathcal{C}$ and $(\widetilde{f}_k: \widetilde{M}_k \rightarrow \widetilde{N}_k)_{k \in \mathbb{N}_0}$ be a sequence of morphism as in Definition~\ref{defn:resolindmorph}. Then the direct limit of the commutative diagrams
\begin{center}
\begin{tikzcd}
    S^{-P(a)_k}T^{n+k}(\widetilde{M}_{D(a)_k}) \arrow[rrrr, "S^{-P(a)_k}T^{n+k}(\widetilde{f}_{D(a)_k})"] \arrow[d, "\delta_a^{n+k}"] & & & & S^{-P(a)_k}T^{n+k}(\widetilde{N}_{D(a)_k}) \arrow[d, "\delta_a^{n+k}"] \\
    S^{-P(a)_{k+1}}T^{n+k+1}(\widetilde{M}_{D(a)_{k+1}}) \arrow[rrrr, "S^{-P(a)_{k+1}}T^{n+k+1}(\widetilde{f}_{D(a)_{k+1}})"] & & & & S^{P(a)_{k+1}}T^{n+k+1}(\widetilde{N}_{D(a)_{k+1}})
\end{tikzcd}
\end{center}
written as $T_a^n(f) := \varinjlim_{k \in \mathbb{N}_0} S^{-P(a)_k}T^{n+k}(\widetilde{f}_{D(a)_k})$ is the induced morphism in the construction $T_a^{\bullet}$.
\end{defn}

We can reiterate the proof of Lemma~\ref{lem:omeganatural} to demonstrate

\begin{lem}\label{lem:omegaanatural}
The homomorphisms $\omega_{a, n}: \widehat{T}^n \rightarrow T_{Res}^n$ form a natural isomorphism. In other words, the square
\begin{center}
\begin{tikzcd}
    T_a^n(M) \arrow[rr, "\omega_{a, n}(M)"] \arrow[d, "T_a^n(f)"] & & T_{Res}^n(M) \arrow[d, "T_{Res}^n(f)"] \\
    T_a^n(N) \arrow[rr, "\omega_{a, n}(N)"] & & T_{Res}^n(N)
\end{tikzcd}
\end{center}
commutes for any morphism $f: M \rightarrow N$ in $\mathcal{C}$.
\end{lem}

We provide explicit formulae for the connecting homomorphisms of the constructions $T_a^{\bullet}$.

\begin{lem}\label{lem:countmanyconstrsconnmorph}
Let $0 \rightarrow A \rightarrow B \rightarrow C \rightarrow 0$ be a short exact sequence in $\mathcal{C}$. Define for any $k \in \mathbb{N}_0$ the morphism
\begin{align*}
\delta_{a, k}^{n+k} &:= (-1)^{-P(a)_k-1} \varepsilon^{-P(a)_k-1} \circ S^{-P(a)_k}\underline{\delta}^{n+k}: \\
&S^{-P(a)_k}T^{n+k}(\widetilde{C}_{D(a)_k}) \rightarrow S^{-P(a)_k-1}T^{n+k+1}(\widetilde{C}_{D(a)_k}) \rightarrow S^{-P(a)_k}T^{(n+1)+k}(\widetilde{A}_{D(a)_k}) \, .
\end{align*}
Then the diagram
\begin{equation}\label{diag:levelconnmorphone}
\begin{tikzcd}
    S^{-P(a)_k}T^{n+k}(\widetilde{C}_{D(a)_k}) \arrow[r, "\delta_{a, k}^{n+k}"] \arrow[d, "\delta_a^{n+k}"] & S^{-P(a)_k}T^{n+k+1}(\widetilde{A}_{D(a)_k}) \arrow[d, "\delta_a^{n+k+1}"] \\
    S^{-P(a)_{k+1}}T^{n+k+1}(\widetilde{A}_{D(a)_{k+1}}) \arrow[r, "\delta_{a, k+1}^{n+k+1}"] & S^{-P(a)_{k+1}}T^{n+k+2}(\widetilde{A}_{D(a)_{k+1}})
\end{tikzcd}
\end{equation}
commutes for every $k \in \mathbb{N}_0$. The direct limit $\widetilde{\delta}_a^n := \varinjlim_{k \in \mathbb{N}_0} \delta_{a, k}^{n+k}: T_a^n(C) \rightarrow T_a^{n+1}(A)$ resulting from these diagrams is the connecting homomorphism in the construction $T_a^{\bullet}$.
\end{lem}

\begin{proof}
We need to prove that Diagram~\ref{diag:levelconnmorphone} commutes. On the one hand, the morphism
\[ \delta_{a, k}^{n+k}: S^{-P(a)_k}T^{n+k}(\widetilde{C}_{D(a)_k}) \rightarrow S^{-P(a)_k}T^{(n+1)+k}(\widetilde{A}_{D(a)_k}) \]
takes us first down the front side of Diagram~\ref{diag:largestdiag} and then over to its back side. On the other hand, the morphism
\[ \delta_a^{n+k}: S^{-P(a)_k}T^{n+k}(\widetilde{C}_{D(a)_k}) \rightarrow S^{-P(a)_{k+1}}T^{n+k+1}(\widetilde{A}_{D(a)_{k+1}}) \]
only takes us down in Diagram~\ref{diag:largestdiag}. The desired diagram commutes because Diagram~\ref{diag:largestdiag} does so too.
\end{proof}

The following provides an analogue of Lemma~\ref{lem:commwithconn} for the constructions $(T_a^{\bullet}, \widetilde{\delta}_a^{\bullet})$.

\begin{lem}\label{lem:commwithconna}
For any short exact sequence $0 \rightarrow A \rightarrow B \rightarrow C \rightarrow 0$ in $\mathcal{C}$ the diagrams
\begin{equation}\label{diag:levelconnmorphtwo}
\begin{tikzcd}
    S^{-P(a)_k}T^{n+k}(\widetilde{C}_{D(a)_k}) \arrow[r, "\widehat{o}_a^k"] \arrow[d, "\delta_a^{n+k}"] & T^{n+k}(\widetilde{C}_k) \arrow[d, "\delta^{n+k}"] \\
    S^{-P(a)_k}T^{(n+1)+k}(\widetilde{A}_{D(a)_k}) \arrow[r, "\widehat{o}_a^{k+1}"] & T^{(n+1)+k}(\widetilde{A}_k)
\end{tikzcd}
\end{equation}
commutes for any $k \in \mathbb{N}_0$. Therefore, the resulting diagram in the direct limit 
\begin{equation}\label{diag:limitconnmorph}
\begin{tikzcd}[row sep = small]
  T_a^n(C) \arrow[r, "\omega_{a, n}"] \arrow[d, "\widetilde{\delta}_a^n"]
  & T_{Res}^n(C) \arrow[d, "\widetilde{\delta}^n"] \\
  T_a^{n+1}(A) \arrow[r, "\omega_{a, n+1}"]
  & T_{Res}^{n+1}(A)
\end{tikzcd}
\end{equation}
also commutes.
\end{lem}

\begin{proof}
The proof is based on tracing the appropriate morphisms in Diagram~\ref{diag:largestdiag}. The term $S^{-P(a)_k}T^{n+k}(\widetilde{C}_{D(a)_k})$ found in Diagram~\ref{diag:levelconnmorphtwo} lies in the front side of Diagram~\ref{diag:largestdiag}. On the one hand, we reach the very right hand side of the latter diagram through the morphism $\widehat{o}_a^k: S^{-P(a)_k}T^{n+k}(\widetilde{C}_{D(a)_k}) \rightarrow T^{n+k}(\widetilde{C}_k)$. The morphism $\delta^{n+k}: T^{n+k}(\widetilde{C}_k) \rightarrow T^{(n+1)+k}(\widetilde{A}_k)$ leads us to the back side. We move by $\mathrm{id}: T^{(n+1)+k}(\widetilde{A}_k) \rightarrow T^{(n+1)+k}(\widetilde{A}_k)$ across a square in the very right hand column of the back side. On the other had, the morphism
\[ \delta_{a, k}^{n+k}: S^{-P(a)_k}T^{n+k}(\widetilde{C}_{D(a)_k}) \rightarrow S^{-P(a)_k}T^{(n+1)+k}(\widetilde{A}_{D(a)_k}) \]
takes us first down the front side of Diagram~\ref{diag:largestdiag} and then over to the back side. Lastly, we reach the left hand side of the right most column of the back side through the morphism $\widehat{o}_a^{k+1}: S^{-P(a)_k}T^{(n+1)+k}(\widetilde{A}_{D(a)_k}) \rightarrow T^{(n+1)+k}(\widetilde{A}_k)$. Since Diagram~\ref{diag:largestdiag} commutes, we conclude that Diagram~\ref{diag:levelconnmorphtwo} also commutes.
\end{proof}

By Lemma~\ref{lem:omegaanatural} and Lemma~\ref{lem:commwithconna}, we conclude that the constructions $(T_a^{\bullet}, \widetilde{\delta}_a^{\bullet})$ with the formula for induced morphisms from Definition~\ref{defn:countmanyconstrsindmorph} and for connecting homomorphisms from Lemma~\ref{lem:countmanyconstrsconnmorph} form Mislin completions of $(T^{\bullet}, \delta^{\bullet})$.

\begin{lem}
The constructions $(T_a^{\bullet}, \widetilde{\delta}_a^{\bullet})$ forms Mislin completions of $(T^{\bullet}, \delta^{\bullet})$. More specifically, $(T_a^{\bullet}, \widetilde{\delta}_a^{\bullet})$  form cohomological functors and $\omega_{a, \bullet}: (T_a^{\bullet}, \widetilde{\delta}_a^{\bullet}) \rightarrow (T_{Res}^{\bullet}, \widetilde{\delta}^{\bullet})$ are isomorphisms of cohomological functors.
\end{lem}

The following proposition shows that the satellite functor construction and the resolution construction can be seen as special cases of the Mislin completions $(T_a^{\bullet}, \widetilde{\delta}_a^{\bullet})$.

\begin{prop}
Let $e \in \lbrace 0, 1 \rbrace^{\mathbb{N}}$ denote the sequence all whose terms are set to $0$ and $f \in \lbrace 0, 1 \rbrace^{\mathbb{N}}$ the sequence all whose terms are set to $1$. Then $(\widehat{T}^{\bullet}, \widehat{\delta}^{\bullet}) = (T_f^{\bullet}, \widetilde{\delta}_f^{\bullet})$ and $(T_{Res}^{\bullet}, \widetilde{\delta}^{\bullet}) = (T_e^{\bullet}, \widetilde{\delta}_e^{\bullet})$.
\end{prop}

\begin{proof}
By definition, $\widehat{T}^n = T_f^n$ and $T_{Res}^n = T_e^n$ for every $n \in \mathbb{Z}$. It can be deduced from Diagram~\ref{diag:largestdiag} that $\widetilde{\delta}_f^n = \widehat{\delta}^n$
and that $\widetilde{\delta}_e^n = \widetilde{\delta}^n$.
\end{proof}

\begin{thm}\label{thm:countmanyconstrs}
For $a, b \in \lbrace 0, 1 \rbrace^{\mathbb{N}}$ the constructions of the Mislin completions $(T_a^{\bullet}, \widetilde{\delta}_a^{\bullet})$ and $(T_b^{\bullet}, \widetilde{\delta}_b^{\bullet})$ are said to agree if there is an infinite (cofinal) subset $\Lambda \subseteq \mathbb{N}$ such that $S^{-P(a)_{\lambda}}T^{n+{\lambda}}(\widetilde{M}_{D(a)_{\lambda}}) = S^{-P(b)_{\lambda}}T^{n+{\lambda}}(\widetilde{M}_{D(b)_{\lambda}})$ for every $\lambda \in \Lambda$, $n \in \mathbb{Z}$ and $M \in \mathrm{obj}(\mathcal{C})$. The basic idea is that the constructions agree whenever their terms
\[ T_a^n(M) = \varinjlim_{k \in \mathbb{N}_0} S^{-P(a)_k}T^{n+k}(\widetilde{M}_{D(a)_k}) \text{ and } T_b^n(M) = \varinjlim_{k \in \mathbb{N}_0} S^{-P(b)_k}T^{n+k}(\widetilde{M}_{D(b)_k}) \]
are isomorphic in a straightforward manner. If this is not the case, the constructions $(T_a^{\bullet}, \widetilde{\delta}_a^{\bullet})$ and $(T_b^{\bullet}, \widetilde{\delta}_b^{\bullet})$ are said to be distinct. Then there are countably many distinct constructions of Mislin completions of the form $(T_a^{\bullet}, \widetilde{\delta}_a^{\bullet})$.
\end{thm}

\begin{proof}
Define the binary relation `$\sim$' on $\lbrace 0, 1 \rbrace^{\mathbb{N}}$ by setting $a \sim b$ if for every $k \in \mathbb{N}$ there is $K \geq k$ such that $P(a)_K = P(b)_K$. If $\sim'$ denotes the transitive closure of $\sim$, then the proof reduces to showing that there are countably many elements in $\lbrace 0, 1 \rbrace^{\mathbb{N}}/\sim'$. Hence, define for any $m \in \mathbb{N}$ the sequences $e^m, f^m \in \lbrace 0, 1 \rbrace^{\mathbb{N}}$ by
\[ e_k^m := \begin{cases} 1 &\text{if } k < m \\ 0 &\text{if } k \geq m \end{cases} \qquad \text{and} \qquad f_k^m := \begin{cases} 0 &\text{if } k < m \\ 1 &\text{if } k \geq m \end{cases} \]
For any distinct $m, p \in \mathbb{N}$ we note that $e_m \nsim' e_p$, $e_m \nsim' f_p$ and $f_m \nsim' f_p$. In particular, there are at least countably many elements in $\lbrace 0, 1 \rbrace^{\mathbb{N}}/\sim'$. If $a \in \lbrace 0, 1 \rbrace^{\mathbb{N}}$ takes either the value $0$ or the value $1$ finitely many times, then there is $m \in \mathbb{N}$ such that $a \sim e_m$ or $a \sim f_m$. If $b, c \in \lbrace 0, 1 \rbrace^{\mathbb{N}}$ each take the values $0$ and $1$ infinitely often, one can construct $d \in \lbrace 0, 1 \rbrace^{\mathbb{N}}$ such that $b \sim d$ and $c \sim d$. Therefore, $\lbrace 0, 1 \rbrace^{\mathbb{N}}/\sim'$ is countable.
\end{proof}

\begin{notn}
In the rest of the paper, we write $(\widehat{T}^{\bullet}, \widehat{\delta}^{\bullet})$ for the Mislin completion of a cohomological functor $(T^{\bullet}, \delta^{\bullet})$ independently from which construction it arises. When required, we specify the construction or use a special notation for it.
\end{notn}

\section{The naïve and the resolution construction}\label{sec:naiveconstruction}

In this section we develop explicit formulae for the induced homomorphisms and connecting homomorphisms of the naïve construction. By virtue of these formulae, we prove that the naïve construction defines Mislin completions for Ext-functors by establishing an isomorphism of cohomological functor to the resolution construction. As already mentioned at the end of Section~\ref{sec:constructions}, the naïve constructions does not define Mislin completions for other cohomological functors.

\subsection{Construction and termwise isomorphisms}

We first perform the naïve construction rigorously and explicitly. Both A.\! Beligiannis and I.\! Reiten in~\cite[p.~37--38/163--165]{bel07} as well as S.\! Guo and L.\! Liang in~\cite[p.~14--15]{guo23} have established relative homological versions of the naïve construction in great generality. In contrast, our construction represents a faithful (absolute homological) generalisation of D.\! J.\! Benson and J.\! F.\! Carlson's original construction found in~\cite[p.~109]{ben92}. The main reason for this is because such a faithful generalisation allows us to develop explicit formulae for the induced homomorphisms and connecting homomorphisms as a Mislin completion. Moreover, this allows us to prove in~\cite{ghe24} that the complete cohomology group $\widehat{H}^0(G, -)$ detects finite cohomological dimension and that Yoneda products exist for completed Ext-functors $\widehat{\mathrm{Ext}}_{\mathcal{C}}^{\bullet}(-, -)$. \\

We define a bifunctor
\[ [-, -]_{\mathcal{C}}: \mathcal{C}^{\mathrm{op}} \times \mathcal{C} \rightarrow \mathbf{Ab} \]
that we have encountered in Subsection~\ref{sec:constructions}. Let $\mathcal{P}_{\mathcal{C}}(A, B) \subseteq \mathrm{Hom}_{\mathcal{C}}(M, N)$ be the subgroup of morphisms factoring through a projective. Define $[A, B]_{\mathcal{C}} := \mathrm{Hom}_{\mathcal{C}}(A, B) / \mathcal{P}_{\mathcal{C}}(A, B)$. It is straightforward to prove the following proposition that is relevant to the construction of Yoneda products in~\cite[Theorem~6.6]{ghe24}.

\begin{prop}\label{prop:biadditive}
The bifunctor $\circ: \mathrm{Hom}_{\mathcal{C}}(B, C) \times \mathrm{Hom}_{\mathcal{C}}(A, B) \rightarrow \mathrm{Hom}_{\mathcal{C}}(A, C)$ descends to a bifunctor $\circ: [B, C]_{\mathcal{C}} \times [A, B]_{\mathcal{C}} \rightarrow [A, C]_{\mathcal{C}}$ that is associative and additive in both variables.
\end{prop}

In order to perform the naïve construction we require the homomorphism from the following proposition.

\begin{prop}\label{prop:transitioning}
There exists a homomorphism $t_{A, B}: [A, B]_{\mathcal{C}} \rightarrow [\widetilde{A}_1, \widetilde{B}_1]_{\mathcal{C}}$ with the property that $t_{B, C}(-) \circ t_{A, B}(-) = t_{A, C}(- \circ -)$.
\end{prop}

\begin{proof}
Consider the following version of Diagram~\ref{diag:complexlift}
\begin{equation}\label{diag:simplelift}
\begin{tikzcd}
    0 \arrow[r] & \widetilde{A}_1 \arrow[r, "\iota_A"] \arrow[d, "\widetilde{f}_1"] & A_0 \arrow[r, "\pi_A"] \arrow[d, "f_0"] & A \arrow[r] \arrow[d, "f"] & 0 \\
    0 \arrow[r] & \widetilde{B}_1 \arrow[r, "\iota_B"] & B_0 \arrow[r, "\pi_B"] & B \arrow[r] & 0
\end{tikzcd}
\end{equation}
where $A_0$ and $B_0$ are assumed to be projective. If $f_0': A_0 \rightarrow B_0$ and $\widetilde{f}_1': \widetilde{A}_1 \rightarrow \widetilde{B}_1$ are different lifts, then there exists a morphism $e: A_0 \rightarrow \widetilde{B}_1$ such that $\widetilde{f}_1-\widetilde{f}_1' = e \circ \iota_A$. As $A_0$ is projective, $\widetilde{f}_1$ and $\widetilde{f}_1'$ agree in $[\widetilde{A}_1, \widetilde{B}_1]_{\mathcal{C}}$. Therefore, there is a well defined homomorphism
\[ s_{A, B}: \mathrm{Hom}_{\mathcal{C}}(A, B) \rightarrow [\widetilde{A}_1, \widetilde{B}_1]_{\mathcal{C}}, f \mapsto \widetilde{f_1}+\mathcal{P}_{\mathcal{C}}(\widetilde{A}_1, \widetilde{B}_1) \, . \]
If $h \in \mathrm{Hom}_{\mathcal{C}}(B, C)$, then we see that $s_{B, C}(h) \circ s_{A, B}(f) = s_{A, C}(h \circ f)$. In particular, if $f \in \mathcal{P}_{\mathcal{C}}(A, B)$, then $s_{A, B}(f)$ = 0. Thus, $s_{A, B}$ descends to the desired homomorphism $t_{A, B}: [A, B]_{\mathcal{C}} \rightarrow [\widetilde{A}_1, \widetilde{B}_1]_{\mathcal{C}}$. \\

Because $s_{B, C}(-) \circ s_{A, B}(-) = s_{A, C}(- \circ -)$, the respective maps factor through quotient homomorphisms as in the commutative diagram
\begin{center}
\begin{tikzcd}[column sep = scriptsize]
	{\scriptstyle \mathrm{Hom}_{\mathcal{C}}(B{,} \, C) \times \mathrm{Hom}_{\mathcal{C}}(A{,} \, B)} \arrow[r] \arrow[rrrr, bend right = 10, "{\scriptstyle s_{B, C} \times s_{A, B}}"] \arrow[ddd, "{\scriptstyle - \circ -}"] & {\scriptstyle \mathrm{Hom}_{\mathcal{C}}(B{,} \, C) \times {[}A{,} \, B{]}_{\mathcal{C}}} \arrow[r] \arrow[rrr, bend left = 15, "{\scriptstyle s_{B, C} \times t_{A, B}}"] & {\scriptstyle {[}B{,} \, C{]}_{\mathcal{C}} \times {[}A{,} \, B{]}_{\mathcal{C}}} \arrow[rr, "{\scriptscriptstyle t_{B, C} \times t_{A, B}}" near start] & & {\scriptstyle {[}\widetilde{B}_1{,} \, \widetilde{C}_1{]} \times {[}\widetilde{A}_1{,} \, \widetilde{B}_1{]}} \arrow[ddd, "{\scriptstyle - \circ -}"] \\ \\ \\
	{\scriptstyle \mathrm{Hom}_{\mathcal{C}}(A{,} \, C)} \arrow[rr] \arrow[rrrr, bend right = 10, "s_{A, C}"] & & {\scriptstyle {[}A{,} \, C{]}_{\mathcal{C}}} \arrow[rr, "{\scriptstyle t_{A, C}}"] & & {\scriptstyle {[}\widetilde{A}_1{,} \, \widetilde{C}_1{]}}
\end{tikzcd}
\end{center}
Its right-hand side and Proposition~\ref{prop:biadditive} imply that $t_{B, C}(-) \circ t_{A, B}(-) = t_{A, C}(- \circ -)$.
\end{proof}

\begin{defn}\label{defn:naiveconstr}
Let $A$ and $B$ be objects in $\mathcal{C}$ and denote by $(A_n)_{n \in \mathbb{N}_0}$, $(B_n)_{n \in \mathbb{N}_0}$ projective resolutions. Then for $n \in \mathbb{Z}$ we define the $n^{\text{th}}$ term of the naïve construction as
\[ BC_{\mathcal{C}}^n(A, B) := \varinjlim_{\mathbf{Ab}}([\widetilde{A}_{n+k}, \widetilde{B}_k]_{\mathcal{C}}, t_{\widetilde{A}_{n+k}, \widetilde{B}_k})_{k \in \mathbb{N}_0, n+k \geq 0} \]
where $t_{\widetilde{A}_{n+k}, \widetilde{B}_{n+k}}: [\widetilde{A}_{n+k}, \widetilde{B}_k]_{\mathcal{C}} \rightarrow [\widetilde{A}_{n+k+1}, \widetilde{B}_{k+1}]_{\mathcal{C}}$. 
\end{defn}

As a first step to prove that the naïve construction defines completed Ext-functors, we construct a homomorphism from the terms $\widehat{\mathrm{Ext}}_{\mathcal{C}}^n(A, B)$ arising from the resolution construction to the terms $BC_{\mathcal{C}}^n(A, B)$. The idea of the construction and of the proof that it yields an isomorphism can be traced back to G.\! Mislin's paper~\cite[p.~299]{mis94}. Again, this differs from S.\! Guo and L.\! Liang's account~\cite[pp.~18--19]{guo23} containing a similar construction and proof. The first step in constructing the desired homomorphism consists of the definition.

\begin{defn}\label{defn:betaprelim}
Let $A$ and $B$ be objects in $\mathcal{C}$ and denote by $(A_n)_{n \in \mathbb{N}_0}$, $(B_n)_{n \in \mathbb{N}_0}$ projective resolutions. Any morphism in the $A_{\bullet}$ can be factorised as $\partial_n: A_n \xrightarrow{\pi_n} \widetilde{A}_n \xrightarrow{\iota_n} A_{n-1}$. Then for every $f \in \mathrm{Ker}(\mathrm{Hom}_{\mathcal{C}}(\partial_{n+1}, B))$ there is a unique morphism $f': \widetilde{A}_n \rightarrow B$ such that $f = f' \circ \pi_n$. Therefore, there is an isomorphisms of abelian groups
\[ \alpha_n(B): \mathrm{Ker}(\mathrm{Hom}_{\mathcal{C}}(\partial_{n+1}, B)) \rightarrow \mathrm{Hom}_{\mathcal{C}}(\widetilde{A}_n, B),\; \; f \mapsto f' \, . \]
In the case where $n = 0$ this yields an isomorphism
\[ \alpha_0(B): \mathrm{Ext}_{\mathcal{C}}^0(A, B) \rightarrow \mathrm{Hom}_{\mathcal{C}}(A, B) \]
since $\mathrm{Ker}(\mathrm{Hom}_{\mathcal{C}}(\partial_0, B)) = \mathrm{Ext}_{\mathcal{C}}^0(A, B)$. If $n \geq 1$ and $f \in \mathrm{Im}(\mathrm{Hom}_{\mathcal{C}}(\partial_n, B))$, then $\alpha_n(B)(f) \in \mathcal{P}_R(\widetilde{A}_n, B)$. Thus $\alpha_n(B)$ descends to a homomorphism of abelian groups
\begin{align*}
\beta_n(B)\! : \, \mathrm{Ext}_{\mathcal{C}}^n(A, B) &\rightarrow [\widetilde{A}_n, B]_{\mathcal{C}} \\
f + \mathrm{Im}(\mathrm{Hom}_{\mathcal{C}}(\partial_n, B)) &\mapsto \alpha_n(B)(f)+ \mathcal{P}_R(\widetilde{A}_n, B) \, .    
\end{align*}
Note that $\beta^n(B)$ is only defined for $n \geq 1$.
\end{defn}

Since all Ext-functors in this and the next section map into the category of abelian groups, we require the following result.

\begin{prop}\label{prop:liminab}
(\cite[p.~261]{osb00}) In the category $\mathbf{Ab}$ of abelian groups, the direct limit of a direct system $\lbrace M_i, \psi_{i, j} \rbrace_{i \leq j \in I}$ can be given as $\bigoplus_{i \in I} M_i/\sim$ where $m_i \in M_i \sim m_j \in M_j$ if there is $k \in I$ such that $\psi_{i, k}(m_i) = \psi_{j, k}(m_j) \in M_k$. In particular, if $\lbrace m_i \in M_i \rbrace_{i \in I}$ is a collection of nonzero elements such that $\psi_{i, j}(m_i) = m_j$ for any $i \leq j$, then the element $(m_i)_{i \in I} \in \varinjlim_{i \in I} M_i$ is nonzero.
\end{prop}

We construct the desired isomorphism from the terms of the resolution construction to the ones of the naïve construction by generalising G.\! Mislin's Theorem~4.1 and its proof from~\cite{mis94}.

\begin{lem}\label{lem:betisomorphism}
Let $A$, $B$ be objects in $\mathcal{C}$. Then for every $n \in \mathbb{Z}$ and $k \in \mathbb{N}_0$ with $n+k \geq 1$ the square
\begin{equation}\label{diag:inducedbeta}
\begin{tikzcd}
  \mathrm{Ext}_{\mathcal{C}}^{n+k}(A, \widetilde{B}_k) \arrow[rr, "\beta_{n+k}(\widetilde{B}_k)"] \arrow[d, "\delta^{n+k}"]
    &{}
    &{[}\widetilde{A}_{n+k}{,} \, \widetilde{B}_k{]}_{\mathcal{C}} \arrow[d, "t_{\widetilde{A}_{n+k}{,} \, \widetilde{B}_k}"] \\
  \mathrm{Ext}_{\mathcal{C}}^{n+k+1}(A, \widetilde{B}_{k+1}) \arrow[rr, "\beta_{n+k+1}(\widetilde{B}_{k+1})"]
    &{}
    &{[}\widetilde{A}_{n+k+1}{,} \, \widetilde{B}_{k+1}{]}_{\mathcal{C}}
\end{tikzcd}
\end{equation}
commutes. The homomorphism resulting in the direct limit of the above diagrams
\[ \beta^n(B) := \varinjlim_{k \in \mathbb{N}_0, n+k \geq 1} \beta_{n+k}(\widetilde{B}_k): \widehat{\mathrm{Ext}}_{\mathcal{C}}^n(A, B) \rightarrow BC_{\mathcal{C}}^n(A, B) \]
is an isomorphism.
\end{lem}

\begin{proof}
To prove that Diagram~\ref{diag:inducedbeta} is commutative, let $f + \mathrm{Im}(\mathrm{Hom}_{\mathcal{C}}(\partial_{n+k}, \widetilde{A}_k)$ be an element in $\mathrm{Ext}_{\mathcal{C}}^{n+k}(A, \widetilde{B}_k)$. It suffices to show that the diagram
\begin{equation}\label{diag:mislinbenson}
\begin{tikzcd}
  A_{n+k+1} \arrow[r, "\pi_{n+k+1}"] \arrow[rd, "F"]
    &\widetilde{A}_{n+k+1} \arrow[r, "\iota_{n+k+1}"] \arrow[d, "\overline{f}^{\ast}"]
    &A_{n+k} \arrow[rr, "\pi_{n+k}"] \arrow[rrd, "f"] \arrow[d, "\overline{f}"]
    & &\widetilde{A}_{n+k} \arrow[d, "\alpha_{n+k}(\widetilde{B}_k)(f)"] \\
  {}
    &\widetilde{B}_{k+1} \arrow[r, "\iota_{k+1}"]
    &B_k \arrow[rr, "\pi_k"]
    & &\widetilde{B}_k
\end{tikzcd}
\end{equation}
commutes where we explain in the following how to construct it. We can lift $\alpha_{n+k}(\widetilde{B}_k)(f)$ to $\overline{f}^{\ast}$ as we have done in Diagram~\ref{diag:simplelift}. In particular, the morphism $\overline{f}^{\ast}$ is a representative of
\[ t_{\widetilde{A}_{n+k}, \widetilde{B}_k} \circ \beta_{n+k}(\widetilde{B}_k)(f + \mathrm{Im}(\mathrm{Hom}_{\mathcal{C}}(\partial_{n+k}, \widetilde{B}_k))) \, . \]
If we write $F$ for $\overline{f}^{\ast} \circ \pi_{n+k+1}$, then one can deduce that
\[ \delta^{n+k} \big(f + \mathrm{Im}(\mathrm{Hom}_{\mathcal{C}}(\partial_{n+k}, \widetilde{B}_k)) \big) = F + \mathrm{Im}(\mathrm{Hom}_{\mathcal{C}}(\partial_{n+k+1}, \widetilde{B}_{k+1})) \in \mathrm{Ext}_{\mathcal{C}}^{n+k+1}(A, \widetilde{B}_{k+1}) \]
by the construction of the connecting homomorphism as in~\cite[pp.~11--12]{wei94}. Thus, $\overline{f}^{\ast}$ is also a representative of $\beta_{n+k+1}(\widetilde{B}_{k+1}) \circ \delta^{n+k}(f + \mathrm{Im}(\mathrm{Hom}_{\mathcal{C}}(\partial_{n+k}, \widetilde{B}_k)))$ and Diagram~\ref{diag:inducedbeta} commutes. \\

Let us show that $\beta^{n}(B)$ is an isomorphism. As $\beta_{n+k}(\widetilde{B}_k)$ is surjective and direct limits in the category of abelian groups are exact, $\beta^n(B)$ is surjective. To demonstrate that $\beta^n(B)$ is injective, let $\varphi \in \widehat{\mathrm{Ext}}_{\mathcal{C}}^n(A, B)$ be an element such that $\beta^n(B)(\varphi) = 0$. By the resolution construction and  Proposition~\ref{prop:liminab}, there exists $k \in \mathbb{N}_0$ and $g \in \mathrm{Ker}\big(\mathrm{Hom}_{\mathcal{C}}(\partial_{n+k+1}, \widetilde{B}_k))\big)$ such that
\[ g + \mathrm{Im}(\mathrm{Hom}_{\mathcal{C}}(\partial_{n+k}, \widetilde{B}_k)) \in \mathrm{Ext}_{\mathcal{C}}^{n+k}(A, \widetilde{B}_k) \]
is mapped to $\varphi \in \widehat{\mathrm{Ext}}_{\mathcal{C}}^n(A, B)$ in the direct limit. Since we may assume that
\[ \beta_{n+k}(\widetilde{B}_k)\big(g + \mathrm{Im}(\mathrm{Hom}_{\mathcal{C}}(\partial_{n+k}, \widetilde{B}_k))\big) = 0 \, , \]
the morphism $\alpha_{n+k}(\widetilde{B}_k)(g): \widetilde{A}_{n+k} \rightarrow \widetilde{B}_k$ from Definition~\ref{defn:betaprelim} factors through a projective. We infer from this and Diagram~\ref{diag:mislinbenson} that there is a morphism $h: \widetilde{A}_{n+k} \rightarrow \widetilde{B}_{k+1}$ rendering the diagram
\begin{center}
\begin{tikzcd}
  A_{n+k+1} \arrow[r, "\pi_{n+k+1}"] \arrow[rd, "G"]
    & \widetilde{A}_{n+k+1} \arrow[rr, "\iota_{n+k+1}"] \arrow[d, "\overline{g}^{\ast}"] & & \widetilde{A}_{n+k} \arrow[lld, "h" near start] \arrow[d, "\overline{g}"] \\
  {}
    & \widetilde{B}_{k+1} \arrow[rr, "\iota_{k+1}" near end] & & B_k
\end{tikzcd}
\end{center}
commutative. Because $G = h \circ \partial_{n+k+1}$ is a representative in $\mathrm{Im}(\delta^{n+k})$, we infer that $\delta^{n+k}(g + \mathrm{Im}(\mathrm{Hom}_{\mathcal{C}}(\partial_{n+k}, \widetilde{B}_k))) = 0$. Therefore, $\varphi = 0$ and $\beta^n(B)$ is injective.
\end{proof}

\subsection{Induced homomorphisms and connecting homomorphisms}

We develop an explicit formula for induced homomorphisms of the naïve construction.

\begin{defn}\label{defn:inducednaivemorph}
Let $A$, $B$, $C$ be objects in $\mathcal{C}$ and let $A_{\bullet}$, $B_{\bullet}$, $C_{\bullet}$ be projective resolutions. Let $f: B \rightarrow C$ be a morphism and let $(\widetilde{f}_k: \widetilde{B}_k \rightarrow \widetilde{C}_k)_{k \in \mathbb{N}_0}$ be a sequence of morphisms as in Definition~\ref{defn:resolindmorph}. By Proposition~\ref{prop:transitioning}, the square of homomorphisms of abelian groups
\begin{equation}\label{diag:inducenaivemorph}
\begin{tikzcd}
    {[}\widetilde{A}_{n+k}, \widetilde{B}_k{]}_{\mathcal{C}} \arrow[rrrrr, "\big(\widetilde{f}_k+\mathcal{P}_{\mathcal{C}}(\widetilde{B}_k{,} \widetilde{C}_k)\big) \circ -"] \arrow[d, "t_{\widetilde{A}_{n+k}, \widetilde{B}_k}"] & & & & & {[}\widetilde{A}_{n+k}, \widetilde{C}_k{]}_{\mathcal{C}} \arrow[d, "t_{\widetilde{A}_{n+k}, \widetilde{C}_k}"] \\
    {[}\widetilde{A}_{n+k+1}, \widetilde{B}_{k+1}{]}_{\mathcal{C}} \arrow[rrrrr, "\big(\widetilde{f}_{k+1}+\mathcal{P}_{\mathcal{C}}(\widetilde{B}_{k+1}{,} \widetilde{C}_{k+1})\big) \circ -"] & & & & & {[}\widetilde{A}_{n+k+1}, \widetilde{C}_{k+1}{]}_{\mathcal{C}}
\end{tikzcd}
\end{equation}
commutes. Then
\[ BC_{\mathcal{C}}^n(A, f) := \varinjlim_{k \in \mathbb{N}_0, n+k \geq 0}\big(\big(\widetilde{f}_k+\mathcal{P}_{\mathcal{C}}(\widetilde{B}_k, \widetilde{C}_k)\big)\circ - \big): BC_{\mathcal{C}}^n(A, B) \rightarrow BC_{\mathcal{C}}^n(A, C) \]
is the induced homomorphism in the naïve construction. \\

Analogously, if $g: C \rightarrow A$ is another morphism with a corresponding sequence of morphisms $(\widetilde{g}_k: \widetilde{B}_k \rightarrow \widetilde{C}_k)_{k \in \mathbb{N}_0}$, then
\[ BC_{\mathcal{C}}^n(g, B) := \varinjlim_{k \in \mathbb{N}_0, n+k \geq 0}\big(- \circ \big(\widetilde{g}_{n+k}+\mathcal{P}_{\mathcal{C}}(\widetilde{C}_{n+k}, \widetilde{A}_{n+k})\big)\big): BC_{\mathcal{C}}^n(A, B) \rightarrow BC_{\mathcal{C}}^n(C, B) \]
defines another induced homomorphism in the naïve construction.
\end{defn}

One can deduce from this definition the following proposition.

\begin{prop}\label{prop:naivebifunctor}
For every $n \in \mathbb{Z}$, $BC_{\mathcal{C}}^n(-, -): \mathcal{C}^{\mathrm{op}} \times \mathcal{C} \rightarrow \mathbf{Ab}$ forms a well defined bifunctor.
\end{prop}

As a step in showing that the naïve construction defines completed Ext-functors, we demonstrate

\begin{lem}\label{lem:betanatural}
For every $n \in \mathbb{Z}$ the isomorphism $\beta^n(-)$ is natural. That is, for every $f \in \mathrm{Hom}_{\mathcal{C}}(B, C)$ the square
\begin{center}
\begin{tikzcd}
  \widehat{\mathrm{Ext}}_{\mathcal{C}}^n(A{,}\, B) \arrow[r, "\beta^n(B)"] \arrow[d, "\widehat{\mathrm{Ext}}_{\mathcal{C}}^n(A{,}\, f)"] & BC_{\mathcal{C}}^n(A{,}\, B) \arrow[d, "BC_{\mathcal{C}}^n(A{,}\, f)"] \\
  \widehat{\mathrm{Ext}}_{\mathcal{C}}^n(A{,}\, C) \arrow[r, "\beta^n(C)"] & BC_{\mathcal{C}}^n(A{,}\, C)
\end{tikzcd}
\end{center}
commutes.
\end{lem}

\begin{proof}
By construction of $\widehat{\mathrm{Ext}}_{\mathcal{C}}^n(A, f)$ from Definition~\ref{defn:resolindmorph} and of $BC_{\mathcal{C}}^n(A, f)$ from Definition~\ref{defn:inducednaivemorph}, it suffices to consider the cube
\begin{center}
\begin{tikzcd}[row sep = scriptsize, column sep = tiny]
  & \mathrm{Ext}_{\mathcal{C}}^{n+k}(A, \widetilde{B}_k) \arrow[dl, "\mathrm{Ext}_{\mathcal{C}}^{n+k}(A{,} \widetilde{f}_k)"] \arrow[rr, "\beta_{n+k}(\widetilde{B}_k)" near end] \arrow[dd, shift left = 2mm, "\delta^{n+k}" near end] & & {[}\widetilde{A}_{n+k}{,} \, \widetilde{B}_k{]}_{\mathcal{C}} \arrow[dl, "{[}\widetilde{A}_{n+k}{,} \, \widetilde{f}_k{]}_{\mathcal{C}}" near end] \arrow[dd, "t_{\widetilde{A}_{n+k}, \widetilde{B}_k}"] \\
  \mathrm{Ext}_{\mathcal{C}}^{n+k}(A, \widetilde{C}_k) \arrow[rr, crossing over, "\beta_{n+k}(\widetilde{C}_k)" near end] \arrow[dd, "\delta^{n+k}"] & & {[}\widetilde{A}_{n+k}{,} \, \widetilde{C}_k{]}_{\mathcal{C}} \\
  & \mathrm{Ext}_{\mathcal{C}}^{n+k+1}(A, \widetilde{B}_{k+1}) \arrow[dl, "\mathrm{Ext}_{\mathcal{C}}^{n+k+1}(A{,} \widetilde{f}_{k+1})" near start] \arrow[rr, "\beta_{n+k+1}(\widetilde{B}_{k+1})" near end] & & {[}\widetilde{A}_{n+k+1}{,} \, \widetilde{B}_{k+1}{]}_{\mathcal{C}} \arrow[dl, "{[}\widetilde{A}_{n+k+1}{,} \, \widetilde{f}_{k+1}{]}_{\mathcal{C}}"] \\
  \mathrm{Ext}_{\mathcal{C}}^{n+k+1}(A, \widetilde{C}_{k+1}) \arrow[rr, "\beta_{n+k+1}(\widetilde{C}_{k+1})" near end] & & {[}\widetilde{A}_{n+k+1}{,} \, \widetilde{C}_{k+1}{]}_{\mathcal{C}} \arrow[from=uu, crossing over, shift right = 3mm, "t_{\widetilde{A}_{n+k}, \widetilde{C}_k}" near end]
\end{tikzcd}
\end{center}
for any $k \in \mathbb{N}_0$ with $n+k \geq 1$. The left hand side commutes and right hand side, too, due to the choice of the sequence of morphisms $(\widetilde{f}_k: \widetilde{B}_k \rightarrow \widetilde{C}_k)_{k \in \mathbb{N}_0}$. The front and back side correspond to Diagram~\ref{diag:inducedbeta}. Lastly, the top and bottom side commute by Definition~\ref{defn:resolindmorph}, Definition~\ref{defn:inducednaivemorph} and Definition~\ref{defn:betaprelim} and form a direct system of commuting squares. The statement of the proposition follows by passing to the direct limit.
\end{proof}

We require the following definition in order to develop an explicit formula for connecting homomorphisms of the naïve construction.

\begin{defn}\label{defn:naiveconnmorph}
Let $0 \rightarrow B \rightarrow C \rightarrow D \rightarrow 0$ be a short exact sequence in $\mathcal{C}$ and denote by $(0 \rightarrow \widetilde{B}_k \rightarrow \widetilde{C}_k \rightarrow \widetilde{D}_k \rightarrow 0)_{k \in \mathbb{N}_0}$ a sequence of short exact sequences as in Definition~\ref{defn:resolconnmorph}. Let $A \in \mathrm{obj}(\mathcal{C})$ and denote by $A_{\bullet}$ a projective resolution. If $f: \widetilde{A}_{n+k} \rightarrow \widetilde{D}_k$ is a morphism in $\mathcal{C}$, let $\overline{f}^{\ast}: \widetilde{A}_{n+k+1} \rightarrow \widetilde{D}_{k+1}$ be a lift of $f$ as in Diagram~\ref{diag:satellitemorph}. Using Diagram~\ref{diag:delta}, we obtain the commutative diagram
\begin{center}
\begin{tikzcd}
  0 \arrow[r]
    & \widetilde{A}_{n+k+1} \arrow[r] \arrow[d, "\overline{f}^{\ast}"]
    & A_{n+k} \arrow[r] \arrow[d, "\overline{f}"]
    & \widetilde{A}_{n+k} \arrow[r] \arrow[d, "f"]
    & 0 \\
  0 \arrow[r]
    & \widetilde{D}_{k+1} \arrow[r] \arrow[d, "h^{\ast}"]
    & D_k \arrow[r] \arrow[d, "h"]
    & \widetilde{D}_k \arrow[r] \arrow[d, "\mathrm{id}_{\widetilde{D}_k}"]
    & 0 \\
  0 \arrow[r]
    & \widetilde{B}_k \arrow[r]
    & \widetilde{C}_k \arrow[r]
    & \widetilde{D}_k \arrow[r]
    & 0
\end{tikzcd}
\end{center}
Then by Proposition~\ref{prop:biadditive} and~\ref{prop:transitioning} there is a homomorphism
\begin{align*}
\tau_{n+k, k} := [\widetilde{A}_{n+k+1}, h^{\ast}]_{\mathcal{C}} \circ t_{\widetilde{A}_{n+k}, \widetilde{D}_k}: \; \; [\widetilde{A}_{n+k}, \widetilde{D}_k]_{\mathcal{C}} &\rightarrow [\widetilde{A}_{n+k+1}, \widetilde{B}_k]_{\mathcal{C}} \\
f + \mathcal{P}_{\mathcal{C}}(\widetilde{A}_{n+k}, \widetilde{D}_k) &\mapsto h^{\ast} \circ \overline{f}^{\ast} + \mathcal{P}_{\mathcal{C}}(\widetilde{A}_{n+k+1}, \widetilde{B}_k)
\end{align*}
which is well defined in the sense that it does not depend on the choice of the morphism $h^{\ast}: \widetilde{D}_{k+1} \rightarrow \widetilde{B}_k$.
\end{defn}

\begin{thm}\label{thm:naiveconnmorph}
Using the terminology from the previous definition, both the squares
\begin{center}
\begin{tikzcd}[column sep = small, row sep = scriptsize]
  {\scriptstyle \mathrm{Ext}_{\mathcal{C}}^{n+k}(A{,} \widetilde{D}_k)} \arrow[rrr, "{\scriptscriptstyle \beta_{n+k}(\widetilde{D}_k)}"] \arrow[dd, "{\scriptscriptstyle \delta^{n+k}}"]
    & & & {\scriptstyle {[}\widetilde{A}_{n+k}{,} \, \widetilde{D}_k{]}_{\mathcal{C}}} \arrow[dd, "{\scriptscriptstyle \tau_{n+k, k}}"] \arrow[ddrr, bend right = 15, white, "\text{{\normalsize and}}" black]
    & & {\scriptstyle {[}\widetilde{A}_{n+k}{,} \, \widetilde{D}_k{]}_{\mathcal{C}}} \arrow[rrrrr, "{\scriptscriptstyle (-1)^k \tau_{n+k, k}}"] \arrow[dd, "{\scriptscriptstyle t_{\widetilde{A}_{n+k}, \widetilde{D}_k}}"]
    & & & & & {\scriptstyle {[}\widetilde{A}_{n+k+1}{,} \, \widetilde{B}_k{]}_{\mathcal{C}}} \arrow[dd, "{\scriptscriptstyle t_{\widetilde{A}_{n+k+1}, \widetilde{B}_k}}"] \\ \\
  {\scriptstyle \mathrm{Ext}_{\mathcal{C}}^{n+k+1}(A{,} \widetilde{B}_k)} \arrow[rrr, "{\scriptscriptstyle \beta_{n+k+1}(\widetilde{B}_k)}"]
    & & & {\scriptstyle {[}\widetilde{A}_{n+k+1}{,} \, \widetilde{B}_k{]}_{\mathcal{C}}}
    & & {\scriptstyle {[}\widetilde{A}_{n+k+1}{,} \, \widetilde{D}_{k+1}{]}_{\mathcal{C}}} \arrow[rrrrr, "{\scriptscriptstyle (-1)^{k+1}\tau_{n+k+1, k+1}}"]
    & & & & & {\scriptstyle {[}\widetilde{A}_{n+k+2}{,} \, \widetilde{B}_{k+1}{]}_{\mathcal{C}}}
\end{tikzcd}
\end{center}
commute for every $k \in \mathbb{N}_0$ with $n+k \geq 1$. The homomorphism in the direct limit of the above right hand square written as
\[ \tau^n := \varinjlim_{k \in \mathbb{N}_0, n+k \geq 1} (-1)^k \tau_{n+k, k}: BC_{\mathcal{C}}^n(A, D) \rightarrow BC_{\mathcal{C}}^{n+1}(A, B) \]
defines the connecting homomorphism in the naïve construction. It fits into the commuting square
\begin{center}
\begin{tikzcd}
     \widehat{\mathrm{Ext}}_{\mathcal{C}}^n(A{,} \, D) \arrow[r, "\beta^n(D)"] \arrow[d, "\widehat{\delta}^n"] & BC_{\mathcal{C}}^n(A{,} \, D) \arrow[d, "\tau^n"] \\
     \widehat{\mathrm{Ext}}_{\mathcal{C}}^{n+1}(A{,} \, B) \arrow[r, "\beta^{n+1}(B)"] & BC_{\mathcal{C}}^{n+1}(A{,} \, B)
\end{tikzcd}
\end{center}
\end{thm}

\begin{proof}
Taking a morphism $h^{\ast}: \widetilde{D}_{k+1} \rightarrow \widetilde{B}_k$ as in Definition~\ref{defn:naiveconnmorph}, we see that the square
\begin{equation}\label{diag:liftfactoring}
\begin{tikzcd}
    \mathrm{Ext}_{\mathcal{C}}^{n+k}(A{,} \, \widetilde{D}_k) \arrow[r, "\delta^{n+k}"] \arrow[d, "\mathrm{id}"] & \mathrm{Ext}_{\mathcal{C}}^{n+k+1}(A{,} \, \widetilde{D}_{k+1}) \arrow[d, "\mathrm{Ext}_{\mathcal{C}}^{n+k+1}(A{,} \, h^{\ast})"] \\
    \mathrm{Ext}_{\mathcal{C}}^{n+k}(A{,} \, \widetilde{D}_k) \arrow[r, "\delta^{n+k}"] & \mathrm{Ext}_{\mathcal{C}}^n(A{,} \, \widetilde{B}_k)
\end{tikzcd}
\end{equation}
commutes. Using this and Definition~\ref{defn:naiveconnmorph}, we can factorise the top left square in the statement of the Lemma as
\begin{center}
\begin{tikzcd}
  \mathrm{Ext}_{\mathcal{C}}^{n+k}(A{,} \, \widetilde{D}_k) \arrow[rrr, "\beta_{n+k}(\widetilde{D}_k)"] \arrow[d, "\delta^{n+k}"] & & & {[}\widetilde{A}_{n+k}{,} \, \widetilde{D}_k{]}_{\mathcal{C}} \arrow[d, "t_{\widetilde{A}_{n+k}, \widetilde{D}_k}"] \\
  \mathrm{Ext}_{\mathcal{C}}^{n+k+1}(A{,} \, \widetilde{D}_{k+1}) \arrow[rrr, "\beta_{n+k+1}(\widetilde{D}_{k+1})"] \arrow[d, "\mathrm{Ext}_{\mathcal{C}}^{n+k+1}(A{,} \, h^{\ast})"] & & & {[}\widetilde{A}_{n+k+1}{,} \, \widetilde{D}_{k+1}{]}_{\mathcal{C}} \arrow[d, "{[}\widetilde{A}_{n+k+1}{,}\, h^{\ast}{]}_{\mathcal{C}}"] \\
  \mathrm{Ext}_{\mathcal{C}}^{n+k+1}(A{,} \, \widetilde{B}_k) \arrow[rrr, "\beta_{n+k+1}(\widetilde{B}_k)"] & & & {[}\widetilde{A}_{n+k+1}{,} \, \widetilde{B}_k{]}_{\mathcal{C}}
\end{tikzcd}
\end{center}
It commutes by Lemma~\ref{lem:betisomorphism} and Lemma~\ref{lem:betanatural}. By this square, Diagram~\ref{diag:anticomm} and Lemma~\ref{lem:betisomorphism} together with its proof, we conclude that the top right square in the statement of the lemma commutes. The latter squares form a direct system in whose direct limit we obtain the bottom square in the lemma.
\end{proof}

By Lemma~\ref{lem:betanatural} and Theorem~\ref{thm:naiveconnmorph}, we conclude that the naïve construction $(BC_{\mathcal{C}}^{\bullet}(A, -), \tau^{\bullet})$ with the formula for induced homomorphisms from Definition~\ref{defn:inducednaivemorph} and for connecting homomorphisms from Theorem~\ref{thm:naiveconnmorph} forms a Mislin completion of $(\mathrm{Ext}_{\mathcal{C}}^{\bullet}(A, -), \delta^{\bullet})$.

\begin{thm}\label{thm:resolbenson}
The naïve construction $(BC_{\mathcal{C}}^{\bullet}(A, -), \tau^{\bullet})$ forms a Mislin completion of $(\mathrm{Ext}_{\mathcal{C}}^{\bullet}(A, -), \delta^{\bullet})$. More specifically, $(BC_{\mathcal{C}}^{\bullet}(A, -), \tau^{\bullet})$ forms a cohomological functor and $\beta^{\bullet}: (\widehat{\mathrm{Ext}}_{\mathcal{C}}^{\bullet}(A, -), \widehat{\delta}^{\bullet}) \rightarrow (BC_{\mathcal{C}}^{\bullet}(A, -), \tau^{\bullet})$ is an isomorphism of cohomological functors.
\end{thm}

\begin{rem}\label{rem:mislinsmistake}
In the case of modules over a ring, G.\! Mislin constructs in~\cite[p.~299]{mis94} an isomorphism
\[ \varinjlim_{k \in \mathbb{N}_0} S^{-k}\mathrm{Ext}^{n+k}(A, -) \rightarrow BC^n(A, -) \]
for every $n \in \mathbb{Z}$. Then he asserts that one can turn $BC^{\bullet}(A, -)$ into a Mislin completion by transferring the structure of $\widehat{\mathrm{Ext}}^{\bullet}(A, -)$. He claims that homomorphisms of the form $[\widetilde{A}_{n+k}, \widetilde{D}_k] \rightarrow [\widetilde{A}_{n+k+1}, \widetilde{B}_k]$ induce the connecting homomorphism
\[ \tau^n: BC^n(A, D) \rightarrow BC^{n+1}(A, B) \, . \]
Since he does not explain how to construct $\tau^n$, he erroneously misses the alternating signs involved in the direct system that gives rise to $\tau^n$ in Theorem~\ref{thm:naiveconnmorph}.
\end{rem}

\section{The hypercohomology and the resolution construction}\label{sec:hypercohom}

We demonstrate in this section that the hypercohomology construction gives rise to completed Ext-functors in the same way as the naïve construction. We achieve this by constructing isomorphisms from the terms of the resolution construction to the ones of the hypercohomology construction. In order to do this, we relate the hypercohomology complex with chain maps and the Vogel complex with almost chain maps. More specifically, we use almost chain maps to develop explicit formulae for induced homomorphisms and connecting homomorphisms in the hypercohomology construction. \\

As the hypercohomology constructions has been generalised previously, let us provide pointers to the literature. Although S.\! Guo and L.\! Liang have turned the cohomology groups of the Vogel complex into a cohomological functor in~\cite[Proposition~4.8]{guo23}, we do not know whether their construction yields completed Ext-functors (Question~\ref{ques:comparewithguoliang}). Similarly, J.\! Hu et al.\! introduce the Vogel complex in~\cite[p.~7]{hu21} in a vastly general setting, but we also do not know whether their construction yields completed Ext-functors.

\subsection{Induced homomorphisms and connecting homomorphisms}\label{subsec:hyperindconn}

We reformulate the hypercohomology construction in terms of almost chain maps in order to develop the above mentioned formulae for induced homomorphisms and connecting homomorphisms. D.\! J.\! Benson and J.\! F.\! Carlson's hint in their paper~\cite[p.~109]{ben92} at such a reformulation in the case of modules over a ring. We require the following notation. If $(M_k, \mu_k)_{k \in \mathbb{Z}}$ is a chain complex in $\mathcal{C}$ with boundary maps $\mu_k: M_k \rightarrow M_{k-1}$ and $n \in \mathbb{Z}$, then we define $(M[n]_k, \mu[n]_k)_{k \in \mathbb{Z}}$ to be $M[n]_k := M_{k+n}$ and $\mu[n]_k := (-1)^n \mu_{k+n}$~\cite[p.~154]{gel03}, \cite[p.~9--10]{wei94}. If $(N_k, \nu_k)_{k \in \mathbb{Z}}$ denotes another chain complex, then we have constructed the hypercohomoology complex $\mathrm{Hyp}_{\mathcal{C}}(M_{\bullet}, N_{\bullet})_{\bullet}$ in Equation~\ref{eq:hypercohombd}. According to~\cite[p.~62--63]{wei94}, an $n$-cocycle of $\mathrm{Hyp}_{\mathcal{C}}(M_{\bullet}, N_{\bullet})_{\bullet}$ is exactly a chain map of the form $M[n]_{\bullet} \rightarrow N_{\bullet}$ where $n$-coboundaries are nullhomotopic chain maps of this form. In the same way, we observe that an $n$-cocycle of the Vogel complex $\mathrm{Vog}_{\mathcal{C}}(M_{\bullet}, N_{\bullet})_{\bullet}$ is a collection of morphisms $(f_{k+n}: M[n]_k \rightarrow N_k)_{k \in \mathbb{Z}}$ such that all but finitely many of the squares
\begin{center}
\begin{tikzcd}
    M_{k+n} \arrow[r, "f_{k+n}"] \arrow[d, "(-1)^n \mu_{n+k}"] & N_k \arrow[d, "\nu_k"] \\
    M_{k-1+n} \arrow[r, "f_{k-1+n}"]  & N_{k-1}
\end{tikzcd}
\end{center}
commute. We follow the convention in~\cite[p.~109]{ben92} by calling this an almost chain map of degree $n$. We call $f_{\bullet+n}: M[n]_{\bullet} \rightarrow N_{\bullet}$ nullhomotopic if there is a collection of morphisms $(e_{k+n}: M[n]_k \rightarrow N_{k+1})_{k \in \mathbb{Z}}$ such that $f_{k+n} = e_{k-1+n} \circ \mu[n]_k + \mu[n]_{k+1} \circ e_{k+n}$ for all but finitely many $k \in \mathbb{Z}$. We observe that an $n$-coboundary of $\mathrm{Vog}_{\mathcal{C}}(M_{\bullet}, N_{\bullet})_{\bullet}$ is exactly such a nullhomotopic almost chain map. Denote by $\mathrm{Hom}_{\mathrm{Ch}(\mathcal{C})}(M[n]_{\bullet}, N_{\bullet})$ the set of chain maps $M[n]_{\bullet} \rightarrow N_{\bullet}$ and by $\mathrm{Null}_{\mathrm{Ch}(\mathcal{C})}(M[n]_{\bullet}, N_{\bullet})$ the subset of nullhomotopic maps. Analogously, write $\widehat{\mathrm{Hom}}_{\mathrm{Ch}(\mathcal{C})}(M[n]_{\bullet}, N_{\bullet})$ for the set of almost chain maps $M[n]_{\bullet} \rightarrow N_{\bullet}$ and by $\widehat{\mathrm{Null}}_{\mathrm{Ch}(\mathcal{C})}(M[n]_{\bullet}, N_{\bullet})$ the subset of nullhomotopic almost chain maps. Above, we have just seen that
\begin{align*}
H^n(\mathrm{Hyp}_{\mathcal{C}}(M_{\bullet}, N_{\bullet})_{\bullet}) &= \mathrm{Hom}_{\mathrm{Ch}(\mathcal{C})}(M[n]_{\bullet}, N_{\bullet})/ \mathrm{Null}_{\mathrm{Ch}(\mathcal{C})}(M[n]_{\bullet}, N_{\bullet}) \text{ and} \\
H^n(\mathrm{Vog}_{\mathcal{C}}(M_{\bullet}, N_{\bullet})_{\bullet}) &= \widehat{\mathrm{Hom}}_{\mathrm{Ch}(\mathcal{C})}(M[n]_{\bullet}, N_{\bullet})/ \widehat{\mathrm{Null}}_{\mathrm{Ch}(\mathcal{C})}(M[n]_{\bullet}, N_{\bullet}) \, .
\end{align*}
Having reformulated the hypercohomology construction, we assume from now on that $M_{\bullet}$, $N_{\bullet}$ are projective resolutions of an object $M, N \in \mathrm{obj}(\mathcal{C})$. \\

The induced homomorphisms and connecting homomorphisms in the hypercohomology construction are modelled on the corresponding homomorphisms for Ext-functors. To substantiate this assertion, we construct an explicit isomorphism from the Ext-group $\mathrm{Ext}_{\mathcal{C}}^n(M, N)$ to the term $H^n(\mathrm{Hyp}_{\mathcal{C}}(M_{\bullet}, N_{\bullet})_{\bullet})$. It is well known that the latter these terms are isomorphic. The classical setting to prove this are derived categories as can be seen in \cite[Chapter~III]{gel03} and \cite[Sections~2.4--2.5]{wei94}. However, one cannot retrieve an explicit isomorphism from this setting. The construction of our isomorphism is similar to a construction in C.\! A.\! Weibel's book~\cite[pp.~63--64]{wei94} where he compares the chain complex $\mathrm{Hom}_{\mathcal{C}}(A_{\bullet}, B)$ to $\mathrm{Hyp}_{\mathcal{C}}(A_{\bullet}, B_{\bullet})_{\bullet, \bullet}$ seen as a double complex. As our techniques differ and we have not found this isomorphism in the literature otherwise, we provide full details. First, we need to distinguish the two emerging notions of Ext-functors whence we introduce the following notation.

\begin{notn}\label{notn:chainext}
For the remainder of this section, if $(A_k)_{k \in \mathbb{Z}}$ is a chain complex in $\mathcal{C}$, then it is assumed that $(A_k)_{k \in \mathbb{N}_0}$ is a projective resolution of an object $A$ in $\mathcal{C}$ and $A_k = 0$ for $k < 0$. We write
\[ \mathcal{E}xt_{\mathcal{C}}^n(A, B) := \mathrm{Hom}_{\mathrm{Ch}(\mathcal{C})}(A[n]_{\bullet}, B_{\bullet})/ \mathrm{Null}_{\mathrm{Ch}(\mathcal{C})}(A[n]_{\bullet}, B_{\bullet}) \]
whenever we consider the $n^{\text{th}}$ Ext-group as arising from chain maps modulo chain homotopy. Analogously, we write
\[ \widehat{\mathcal{E}xt}_{\mathcal{C}}^n(A, B) := \widehat{\mathrm{Hom}}_{\mathrm{Ch}(\mathcal{C})}(A[n]_{\bullet}, B_{\bullet})/ \widehat{\mathrm{Null}}_{\mathrm{Ch}(\mathcal{C})}(A[n]_{\bullet}, B_{\bullet}) \, . \]
If $a_k: A_k \rightarrow A_{k-1}$ denote the boundary maps of $A_{\bullet}$, then we continue to write
\[ \mathrm{Ext}_{\mathcal{C}}^n(A, B) = \mathrm{Ker}(\mathrm{Hom}_{\mathcal{C}}(a_{n+1}, B))/ \mathrm{Im}(\mathrm{Hom}_{\mathcal{C}}(a_n, B)) \]
for the $n^{\text{th}}$ Ext-group defined as the $n^{\text{th}}$ derived functor of $\mathrm{Hom}_{\mathcal{C}}(-, B)$.
\end{notn}

The required isomorphism is given by

\begin{defn}\label{defn:extisoms}
Let $A, B \in \mathrm{obj}(\mathcal{C})$, $(A_{\bullet}, a_{\bullet})$, $(B_{\bullet}, b_{\bullet})$ be projective resolutions and $n \in \mathbb{N}_0$. By Definition~\ref{defn:betaprelim} and the Comparison Theorem~\cite[Theorem~2.2.6]{wei94}, there exists for every $\varphi \in \mathrm{Ker}(\mathrm{Hom}_{\mathcal{C}}(a_{k+1}, B))$ a chain map $(\varphi[n]_k: A[n]_k \rightarrow B_k)_{k \in \mathbb{Z}}$ that is unique up to chain homotopy and that renders the diagram
\begin{equation}\label{diag:lifttochainmap}
\begin{tikzcd}
    {} \dots {} \arrow[r] & A_{n+2} \arrow[rr, "(-1)^n a_{n+2}"] \arrow[d, "\varphi{[}n{]}_2"] & & A_{n+1} \arrow[rr, "(-1)^n a_{n+1}"] \arrow[d, "\varphi{[}n{]}_1"] & & A_n \arrow[rr, "\pi_n"] \arrow[drr, "\varphi"] \arrow[d, "\varphi{[}n{]}_0" near end] & & \widetilde{A}_n \arrow[d, "\alpha_n(B)(\varphi)"] \\
    {} \dots {} \arrow[r] & B_2 \arrow[rr, "b_2"] & & B_1 \arrow[rr, "b_1"] & & B_0 \arrow[rr, "b"] & & B
\end{tikzcd}
\end{equation}
commutative. If $\varphi \in \mathrm{Im}(\mathrm{Hom}_{\mathcal{C}}(a_n, B))$, then one can deduce that $\varphi[n]_{\bullet}$ is chain homotopic to the zero chain map. Therefore, there is a well defined homomorphism
\begin{align*}
\zeta_n: \; \mathrm{Ext}_{\mathcal{C}}^n(A, B) &\rightarrow \mathcal{E}xt_{\mathcal{C}}^n(A, B) \\
\varphi + \mathrm{Im}(\mathrm{Hom}_{\mathcal{C}}(a_n, B)) &\mapsto \varphi[n]_{\bullet} + \mathrm{Null}_{\mathrm{Ch}(\mathcal{C})}(A[n]_{\bullet}, B_{\bullet}) \, .
\end{align*}
Given that Ext-functors of negative degree vanish, we set $\zeta_n := 0$ in the case where $n < 0$.
\end{defn}

The following proposition follows from the above definition and justifies the notation introduced above.

\begin{prop}
For every $n \in \mathbb{Z}$, $\zeta_n: \mathrm{Ext}^n(A, B) \rightarrow \mathcal{E}xt_{\mathcal{C}}^n(A, B)$ is an isomorphism.
\end{prop}

With this notation at hand, we define the induced homomorphisms for the hypercohomology construction. Using the above isomorphism, we shall see in the next subsection how these induced homomorphisms are a reformulation of the corresponding homomorphisms in Ext-functors.

\begin{defn}\label{defn:extchainmaps}
Let $f: B \rightarrow C$ and $g: C \rightarrow A$ be morphisms in $\mathcal{C}$. Assume that $A_{\bullet}$, $B_{\bullet}$, $C_{\bullet}$ are projective resolutions. By the Comparison Theorem~\cite[Theorem~2.2.6]{wei94}, there is a chain map $f_{\bullet}: B_{\bullet} \rightarrow C_{\bullet}$ that is unique up to chain homotopy and that renders the diagram
\begin{center}
\begin{tikzcd}
    {} \dots {} \arrow[r] & B_2 \arrow[r] \arrow[d, "f_2"] & B_1 \arrow[r] \arrow[d, "f_1"] & B_0 \arrow[r] \arrow[d, "f_0"] & B \arrow[d, "f"] \\
    {} \dots {} \arrow[r] & C_2 \arrow[r] & C_1 \arrow[r] & C_0 \arrow[r] & C
\end{tikzcd}
\end{center}
commutative. Then the homomorphisms
\begin{align*}
\widehat{\mathcal{E}xt}_{\mathcal{C}}^n(A, f): \;  \widehat{\mathcal{E}xt}_{\mathcal{C}}^n(A, B) &\rightarrow \widehat{\mathcal{E}xt}_{\mathcal{C}}^n(A, C) \\
\varphi_{\bullet} + \widehat{\mathrm{Null}}_{\mathrm{Ch}(\mathcal{C})}(A[n]_{\bullet}, B_{\bullet}) &\mapsto f_{\bullet} \circ \varphi_{\bullet} + \widehat{\mathrm{Null}}_{\mathrm{Ch}(\mathcal{C})}(A[n]_{\bullet}, C_{\bullet})
\end{align*}
and
\begin{align*}
\widehat{\mathcal{E}xt}_{\mathcal{C}}^n(g, B): \;  \widehat{\mathcal{E}xt}_{\mathcal{C}}^n(A, B) &\rightarrow \widehat{\mathcal{E}xt}_{\mathcal{C}}^n(C, B) \\
\varphi_{\bullet} + \widehat{\mathrm{Null}}_{\mathrm{Ch}(\mathcal{C})}(A[n]_{\bullet}, B_{\bullet}) &\mapsto \varphi_{\bullet} \circ g[n]_{\bullet} + \widehat{\mathrm{Null}}_{\mathrm{Ch}(\mathcal{C})}(C[n]_{\bullet}, B_{\bullet}) \, .
\end{align*}
are the induced homomorphisms of the hypercohomology construction. The induced homomorphisms
\[ \mathcal{E}xt_{\mathcal{C}}^n(A, f): \mathcal{E}xt_{\mathcal{C}}^n(A, B) \rightarrow \mathcal{E}xt_{\mathcal{C}}^n(A, C) \text{ and } \mathcal{E}xt_{\mathcal{C}}^n(g, B): \mathcal{E}xt_{\mathcal{C}}^n(A, B) \rightarrow \mathcal{E}xt_{\mathcal{C}}^n(C, B) \]
are defined analogously.
\end{defn}

We deduce from this definition

\begin{prop}\label{prop:hyperbifunctor}
For any $n \in \mathbb{Z}$, both
\[ \mathcal{E}xt_{\mathcal{C}}^n(-, -): \mathcal{C}^{\mathrm{op}} \times \mathcal{C} \rightarrow \mathbf{Ab} \text{ and } \, \widehat{\mathcal{E}xt}_{\mathcal{C}}^n(-, -): \mathcal{C}^{\mathrm{op}} \times \mathcal{C} \rightarrow \mathbf{Ab} \]
form well defined bifunctors.
\end{prop}

Finally, we define connecting homomorphisms of the hypercohomology construction.

\begin{defn}\label{defn:hyperconnmorph}
Let $0 \rightarrow B \rightarrow C \rightarrow D \rightarrow 0$ and $0 \rightarrow \widetilde{D}_1 \rightarrow D_0 \rightarrow D \rightarrow 0$ be short exact sequences in $\mathcal{C}$ where $D_0$ is projective.Let $(D_{\bullet}, d_{\bullet})$ be a projective resolution of $D$. Define the chain complex $(D_{\bullet}^{+1}, d_{\bullet}^{+1})$ by setting $(D_m^{+1}, d_m^{+1})_{m \in \mathbb{N}} := (D_{m+1}, d_{m+1})_{m \in \mathbb{N}}$, $D_0^{+1} := D_1$ and $D_m^{+1} := 0$ for $m < 0$. \\

Note that $D_{\bullet}^{+1}$ is a projective resolution of $\widetilde{D}_1$. Define the chain map $\pi_{\bullet}^n: D_{\bullet} \rightarrow D^{+1}[-1]_{\bullet}$ by setting
\begin{equation}
\pi_m^n := \begin{cases} (-1)^{m-1+n}\mathrm{id}_{D_m}: D_m \rightarrow D_m &\text{if } m \geq 1 \\ {} \; \; 0 &\text{if } m \leq 0 \end{cases}
\end{equation}
According to~\cite[p.~166]{gel03},
\[ \mathcal{E}xt_{\mathcal{C}}^{n+1}(A, D) = \mathrm{Hom}_{\mathrm{Ch}(\mathcal{C})}(A[n]_{\bullet}, D[-1]_{\bullet})/ \mathrm{Null}_{\mathrm{Ch}(\mathcal{C})}(A[n]_{\bullet}, D[-1]_{\bullet}) \]
where an analogous statement holds true for $\widehat{\mathcal{E}xt}_{\mathcal{C}}^{n+1}(A, D)$. Then
\begin{align*}
\Delta^n: \mathcal{E}xt_{\mathcal{C}}^n(A, D) &\rightarrow \mathcal{E}xt_{\mathcal{C}}^{n+1}(A, \widetilde{D}_1) \\
\varphi_{\bullet+n}+\mathrm{Null}_{\mathrm{Ch}(\mathcal{C})}(A[n]_{\bullet}, D_{\bullet}) &\mapsto \pi_{\bullet}^n \circ \varphi_{\bullet+n}+\mathrm{Null}_{\mathrm{Ch}(\mathcal{C})}(A[n]_{\bullet}, D^{+1}[-1]_{\bullet})
\end{align*}
is the connecting homomorphism associated to $0 \rightarrow \widetilde{D}_1 \rightarrow D_0 \rightarrow D \rightarrow 0$. Define the connecting homomorphism $\widehat{\Delta}^n: \widehat{\mathcal{E}xt}_{\mathcal{C}}^n(A, D) \rightarrow \widehat{\mathcal{E}xt}_{\mathcal{C}}^{n+1}(A, \widetilde{D}_1)$ analogously. \\

If $h^{\ast}: \widetilde{D}_1 \rightarrow B$ is a lift of $\mathrm{id}_D: D \rightarrow D$ as in Diagram~\ref{diag:delta} and $h^{\ast}_{\bullet+1}: D_{\bullet}^{+1} \rightarrow B_{\bullet}$ a chain map as in the Comparison Theorem, then
\begin{align*}
\Delta^n: \mathcal{E}xt_{\mathcal{C}}^n(A, D) &\rightarrow \mathcal{E}xt_{\mathcal{C}}^{n+1}(A, B) \\
\varphi_{\bullet+n}+\mathrm{Null}_{\mathrm{Ch}(\mathcal{C})}(A[n]_{\bullet}, D_{\bullet}) &\mapsto h^{\ast}[-1]_{\bullet+1} \circ \pi_{\bullet}^n \circ \varphi_{\bullet+n}+\mathrm{Null}_{\mathrm{Ch}(\mathcal{C})}(A[n]_{\bullet}, B[-1]_{\bullet})
\end{align*}
is the connecting homomorphism associated to $0 \rightarrow B \rightarrow C \rightarrow D \rightarrow 0$. Define the connecting homomorphism $\widehat{\Delta}^n: \widehat{\mathcal{E}xt}_{\mathcal{C}}^n(A, D) \rightarrow \widehat{\mathcal{E}xt}_{\mathcal{C}}^{n+1}(A, B)$ analogously.
\end{defn}

\subsection{Termwise isomorphisms}

We prove that the Ext-groups $\mathcal{E}xt_{\mathcal{C}}^n(A, B)$ defined through chain maps form cohomological functors as a step in constructing an isomorphism from the terms of the resolution construction to the ones of the hypercohomology construction.

\begin{lem}\label{lem:ordinaryhypercohom}
The isomorphisms $\zeta_{\bullet}: (\mathrm{Ext}_{\mathcal{C}}^{\bullet}(A, -), \delta^{\bullet}) \rightarrow (\mathcal{E}xt_{\mathcal{C}}^{\bullet}(A, -), \Delta^{\bullet})$ from Definition~\ref{defn:extisoms} form an isomorphism of cohomological functors. In particular, $(\mathcal{E}xt_{\mathcal{C}}^{\bullet}(A, -), \Delta^{\bullet})$ forms a cohomological functor.
\end{lem}

\begin{proof}
It follows from Definition~\ref{defn:extchainmaps} that the isomorphisms $\zeta_n: \mathrm{Ext}_{\mathcal{C}}^n(A, -) \rightarrow \mathcal{E}xt_{\mathcal{C}}^n(A, -)$ are natural. Hence, it suffices to prove that $\zeta_n$ commutes with the connecting homomorphisms. Let $0 \rightarrow B \rightarrow C \rightarrow D \rightarrow 0$ and $0 \rightarrow \widetilde{D}_1 \rightarrow D_0 \rightarrow D \rightarrow 0$ be short exact sequences in $\mathcal{C}$ where $D_0$ is projective. If $h^{\ast}: \widetilde{D}_1 \rightarrow B$ denotes a lift of $\mathrm{id}_D$ as in Definition~\ref{defn:hyperconnmorph}, then by Diagram~\ref{diag:liftfactoring} the connecting homomorphism associated to the former sequence factors through the one associated to the latter as
\[ \delta^n: \mathrm{Ext}_{\mathcal{C}}^n(A, D) \xrightarrow{\delta^n} \mathrm{Ext}_{\mathcal{C}}^{n+1}(A, \widetilde{D}_1) \xrightarrow{\mathrm{Ext}_{\mathcal{C}}^{n+1}(A, h^{\ast})} \mathrm{Ext}_{\mathcal{C}}^{n+1}(A, B) \, . \]
Thus, the left hand triangle in the diagram
\begin{equation}\label{diag:trapezia}
\begin{tikzcd}
  \mathrm{Ext}_{\mathcal{C}}^n(A{,} \, D) \arrow[rrr, "\zeta_n"] \arrow[dr, "\delta^n"] \arrow[dd, "\delta^n"]
    & & & \mathcal{E}xt_{\mathcal{C}}^n(A{,} \, D) \arrow[dl, "\Delta^n"] \arrow[dd, shift left = 2mm, "\Delta^n"] \\
  & \mathrm{Ext}_{\mathcal{C}}^{n+1}(A{,} \, \widetilde{D}_1) \arrow[r, "\zeta_{n+1}"] \arrow[dl, "\mathrm{Ext}_{\mathcal{C}}^{n+1}(A{,} \, h^{\ast})" near start]
    & \mathcal{E}xt_{\mathcal{C}}^{n+1}(A{,} \, \widetilde{D}_1) \arrow[dr, "\mathcal{E}xt_{\mathcal{C}}^{n+1}(A{,} \, h^{\ast})" near start] & \\
  \mathrm{Ext}_{\mathcal{C}}^{n+1}(A{,} \, B) \arrow[rrr, "\zeta_{n+1}"]
    & & & \mathcal{E}xt_{\mathcal{C}}^{n+1}(A{,} \, B)
\end{tikzcd}
\end{equation}
commutes. The bottom trapezium is commutes and the right hand triangle triangle, too, by Definition~\ref{defn:extchainmaps} and Definition~\ref{defn:hyperconnmorph}. The top trapzium can be seen as a generalisation of \cite[Proposition~1.1]{emm17} where we use a different method to prove that it commutes. Let $\varphi \in \mathrm{Ker}(\mathrm{Hom}_{\mathcal{C}}(a_n, D))$. We explain how one can combine Diagram~\ref{diag:mislinbenson} and Diagram~\ref{diag:lifttochainmap} as
\begin{equation}\label{diag:zetanddelta}
\begin{tikzcd}
    {} \dots {} \arrow[r]
    & A_{n+2} \arrow[rr, "(-1)^n a_{n+2}"] \arrow[d, "\varphi{[}n{]}_2"]
    & & A_{n+1} \arrow[r, "\pi_{n+1}" near end] \arrow[rr, bend left = 26, "a_{n+1}"] \arrow[d, "\varphi{[}n{]}_1" near end] \arrow[dr, "\Phi"]
    & \widetilde{A}_{n+1} \arrow[r, "\iota_{n+1}" near start] \arrow[d]
    & A_n \arrow[r, "\pi_n"] \arrow[dr, bend left = 15, "\varphi"] \arrow[d, "\varphi{[}n{]}_0"]
    & \widetilde{A}_n \arrow[d, "\alpha_{n+k}(B)(\varphi)"] \\
    {} \dots {} \arrow[r]
    & D_2 \arrow[rr, "d_2"]
    & & D_1 \arrow[r, "\pi_1"] \arrow[rr, bend right = 28, "d_1"]
    & \widetilde{D}_1 \arrow[r, "\iota_1"]
    & D_0 \arrow[r, "\pi_0"]
    & D
\end{tikzcd}
\end{equation}
The right most morphisms up to $\Phi: A_{n+1} \rightarrow \widetilde{B}_1$ are taken from Diagram~\ref{diag:mislinbenson} while $\varphi: A_n \rightarrow D$ and $\varphi[n]_{\bullet}: A[n]_{\bullet} \rightarrow D_{\bullet}$ are the morphisms as they occur in Diagram~\ref{diag:lifttochainmap}. In particular,
\[ \delta^n \big(\varphi+\mathrm{Im}(\mathrm{Hom}_{\mathcal{C}}(a_n, D))\big) = \Phi+\big(\varphi+\mathrm{Im}(\mathrm{Hom}_{\mathcal{C}}(a_n, \widetilde{D}_1)) \, . \]
However, $\iota_{n+1} \circ \pi_{n+1} = a_{n+1}$ might not agree with $(-1)^n a_{n+1}: A_{n+1} \rightarrow A_n$. This together Diagram~\ref{diag:mislinbenson} and the Comparison from~\cite[p.~36]{wei94} implies that $\Phi = (-1)^n \pi_1 \circ \varphi[n]_1$. Therefore, the chain map $\Phi[n+1]_{\bullet}: A[n+1] \rightarrow \widetilde{D}_{\bullet}^{+1}$ defined by
\[ \Phi[n+1]_m := \begin{cases} (-1)^{n+m} \varphi[n]_{m+1}: A_{m+n+1} \rightarrow D_{m+1} &\text{if } m \geq 0 \\ \; \; 0 &\text{if } m < 0 \end{cases} \]
is a lift of $\Phi$ as in Diagram~\ref{diag:lifttochainmap}. The alternating signs are due to the fact that the boundary maps of $A[n]_{\bullet}$ and $A[n+1]$ have opposite signs. We conclude by Diagram~\ref{diag:zetanddelta} and Definition~\ref{defn:hyperconnmorph} that
\[ \zeta_{n+1}\Big(\delta^n \big(\varphi+\mathrm{Im}(\mathrm{Hom}_{\mathcal{C}}(a_n, D))\big)\Big) = \Delta^n \Big(\zeta^n \big(\varphi+\mathrm{Im}(\mathrm{Hom}_{\mathcal{C}}(a_n, D))\big)\Big) \, . \]
The top trapezium and thus all of Diagram~\ref{diag:trapezia} commutes, proving that $\zeta_{\bullet}$ commutes with the connecting homomorphisms.
\end{proof}

We reformulate the resolution construction in terms of direct limit of chain maps as another step in constructing an isomorphism from the terms of the resolution construction to the ones of the hypercohomology construction.

\begin{defn}
Let $B \in \mathrm{obj}(\mathcal{C})$ and denote by $B_{\bullet}$ a projective resolution. For any $n \in \mathbb{Z}$ define
\[ \mathcal{E}xt_{Res, \mathcal{C}}^n(A, B) := \varinjlim_{k \in \mathbb{N}_0} (\mathcal{E}xt_{\mathcal{C}}^{n+k}(A, \widetilde{B}_k), \Delta^{n+k}) \, . \]
Let $f: B \rightarrow C$ be a morphism in $\mathcal{C}$ and $(\widetilde{f}_k: \widetilde{B}_k \rightarrow \widetilde{C}_k)_{k \in \mathbb{N}_0}$ be a sequence of morphisms as in Definition~\ref{defn:resolindmorph}. Then define the induced homomorphism by
\[ \mathcal{E}xt_{Res, \mathcal{C}}^n(A, f) := \varinjlim_{k \in \mathbb{N}_0} \mathcal{E}xt_{\mathcal{C}}^{n+k}(A, \widetilde{f}_k): \mathcal{E}xt_{\mathcal{C}}^n(A, B) \rightarrow \mathcal{E}xt_{\mathcal{C}}^n(A, C) \, . \]
Let $0 \rightarrow B \rightarrow C \rightarrow D \rightarrow 0$ be a short exact sequence and $(0 \rightarrow \widetilde{A}_k \rightarrow \widetilde{B}_k \rightarrow \widetilde{C}_k \rightarrow 0)_{k \in \mathbb{N}_0}$ be a sequence of short exact sequences as in Definition~\ref{defn:resolconnmorph}. Then define the connecting homomorphism by
\[ \widetilde{\Delta}^n := \varinjlim_{k \in \mathbb{N}_0} (-1)^k \Delta^{n+k}: \mathcal{E}xt_{Res, \mathcal{C}}^n(A, D) \rightarrow \mathcal{E}xt_{Res, \mathcal{C}}^{n+1}(A, B) \]
where $\Delta^{n+k}: \mathcal{E}xt_{\mathcal{C}}^{n+k}(A, \widetilde{D}_k) \rightarrow \mathcal{E}xt_{\mathcal{C}}^{n+k+1}(A, \widetilde{B}_k)$ are taken as in Diagram~\ref{diag:anticomm}.
\end{defn}

\begin{defn}\label{defn:resolinchains}
By virtue of the isomorphisms of cohomological functors 
\[ \zeta_{n+k}: \mathrm{Ext}_{\mathcal{C}}^{n+k}(A, -) \rightarrow \mathcal{E}xt_{\mathcal{C}}^{n+k}(A, -) \]
from Lemma~\ref{lem:ordinaryhypercohom}, define the natural isomorphisms
\[ \zeta^n := \varinjlim_{k \in \mathbb{N}_0} \zeta_{n+k}: \mathrm{Ext}_{Res, \mathcal{C}}^n(A, -) \rightarrow \mathcal{E}xt_{Res, \mathcal{C}}^n(A, -) \, . \]
Then
\[ \zeta^{\bullet}: (\mathrm{Ext}_{Res, \mathcal{C}}^{\bullet}(A, -), \widetilde{\delta}^{\bullet}) \rightarrow (\mathcal{E}xt_{Res, \mathcal{C}}^{\bullet}(A, -), \widetilde{\Delta}^{\bullet}) \]
forms an isomorphism of cohomological functors.
\end{defn}


We introduce the following terminology in order ease notation when constructing a homomorphism from $\mathcal{E}xt_{Res, \mathcal{C}}^n(A, B)$ to $\widehat{\mathcal{E}xt}_{\mathcal{C}}^n(A, B)$.

\begin{conv}
For the remainder of this section, we consider the cofinal system
\[ (\mathcal{E}xt_{\mathcal{C}}^{n+2k}(A, \widetilde{B}_{2k}), \Delta^{n+2k+1} \circ \Delta^{n+2k})_{k \in \mathbb{N}_0} \]
and specifically form $\mathcal{E}xt_{Res, \mathcal{C}}^n(A, B)$ as its direct limit.
\end{conv}

\begin{notn}
If $(B_{\bullet}, b_{\bullet})$ is a projective resolution of $B$ and $k \in \mathbb{N}_0$, then we define the projective resolution $(B_{\bullet}^{+k}, b_{\bullet}^{+k})_{m \in \mathbb{Z}}$ of $\widetilde{B}_k$ by setting $(B_m^{+k}, b_m^{+k})_{m \in \mathbb{N}} := (B_{m+k}, b_{m+k})_{m \in \mathbb{N}}$, $B_0^{+k} := B_k$ and $B_m^{+k} := 0$ for $m < 0$.
\end{notn}

\begin{lem}\label{lem:vogelresol}
There is an isomorphism $\vartheta^n(B): \mathcal{E}xt_{Res, \mathcal{C}}^n(A, B) \rightarrow \widehat{\mathcal{E}xt}_{\mathcal{C}}^n(A, B)$ for any $B \in \mathrm{obj}(\mathcal{C})$ and $n \in \mathbb{Z}$.
\end{lem}

\begin{proof}
For any $k \in \mathbb{N}_0$ we construct a homomorphism
\begin{align*}
\vartheta_{n}^{2k}(B): \mathcal{E}xt_{\mathcal{C}}^{n+2k}(A, \widetilde{B}_{2k}) &\rightarrow \widehat{\mathcal{E}xt}_{\mathcal{C}}^n(A, B) \\
\varphi_{\bullet+n+2k} + \mathrm{Null}_{\mathrm{Ch}(\mathcal{C})}(A[n+2k]_{\bullet}, B_{\bullet}^{+2k}) &\mapsto \Phi_{\bullet+n} + \widehat{\mathrm{Null}}_{\mathrm{Ch}(\mathcal{C})}(A[n]_{\bullet}, B_{\bullet})
\end{align*}
as follows. For a representative $\varphi_{\bullet+n+2k} \in \mathrm{Hom}_{\mathrm{Ch}(\mathcal{C})}(A[n+2k]_{\bullet}, B_{\bullet}^{+2k})$ of an element in $\mathcal{E}xt_{\mathcal{C}}^{n+2k}(A, \widetilde{B}_{2k})$ we define $\Phi_{\bullet+n}: A[n]_{\bullet} \rightarrow B_{\bullet}$ by
\[ \Phi_{m+n} := \begin{cases} \varphi_{m+n}: A_{m+n} \rightarrow B_m &\text{if } m \geq 2k \\ \; \; 0 &\text{if } m < 2k \end{cases} \]
As the boundary maps of $A[n]$ and $A[n+2k]$ agree modulo a shift of $2k$ in their indices, all squares of the form
\begin{center}
\begin{tikzcd}
  A_{m+n} \arrow[rr, "\Phi_{m+n}"] \arrow[d, "(-1)^n a_{m+n}"]
    & & B_m \arrow[d, "b_m"] \\
  A_{m-1+n} \arrow[rr, "\Phi_{n+m-1}"]
    & & B_{m-1}
\end{tikzcd}
\end{center}
commute for $m \in \mathbb{Z} \setminus \lbrace 2k \rbrace$. We conclude that $\Phi_{\bullet+n} \in \widehat{\mathrm{Hom}}_{\mathrm{Ch}(\mathcal{C})}(A[n]_{\bullet}, B_{\bullet})$. If the chain map $\varphi_{\bullet+n+2k}$ is nullhomotopic, then $\Phi_{\bullet+n}$ is genuinely nullhomotopic in any degree $m \in \mathbb{Z} \setminus \lbrace 2k-1 \rbrace$, hence nullhomotopic as an almost chain map. Thus, $\vartheta_{n}^{2k}$ is a well defined homomorphism. \\

By definition of $\vartheta_{n}^{2k}$ and by Definition~\ref{defn:hyperconnmorph}, the triangle
\begin{equation}\label{diag:thetadeltatriangle}
\begin{tikzcd}
  \mathcal{E}xt_{\mathcal{C}}^{n+2k}(A{,} \, \widetilde{B}_{2k}) \arrow[ddrr, "\vartheta_{n}^{2k}(B)"] \arrow[dddd, shift right = 6mm, "\Delta^{n+2k+1} \circ \Delta^{n+2k}"] & & {} \\ \\
  {} & & \widehat{\mathcal{E}xt}_{\mathcal{C}}^n(A{,} \, B) \\ \\
  \mathcal{E}xt_{\mathcal{C}}^{n+2k+2}(A{,} \, \widetilde{B}_{2k+2)}) \arrow[uurr, "\vartheta_{n}^{2k+2)}(B)" near end] & & {}
\end{tikzcd}
\end{equation}
commutes. We write the resulting homomorphism in the direct limit of the above triangles as
\[ \vartheta^n(B) := \varinjlim_{k \in \mathbb{N}_0} \vartheta_{n}^{2k}(B): \; \mathcal{E}xt_{Res, \mathcal{C}}^n(A, B) \rightarrow \widehat{\mathcal{E}xt}_{\mathcal{C}}^n(A, B) \, . \]
To show that it is surjective, let $\Phi_{\bullet+n}+\widehat{\mathrm{Null}}_{\mathrm{Ch}(\mathcal{C})}(A[n]_{\bullet}, B_{\bullet})$ be in $\widehat{\mathcal{E}xt}_{\mathcal{C}}^n(A, B)$. Since it is an almost chain map, there is $k \in \mathbb{N}_0$ such that $(\Phi_{m+n}: A_{m+n} \rightarrow B_m)_{m \geq 2k}$ is a chain map. Defining $\varphi_{\bullet+n+2k}: A[n+2k]_{\bullet} \rightarrow B_{\bullet}^{+2k}$ by setting
\[ \varphi_{m+n+2k} := \begin{cases}\Phi_{m+n+2k}: A_{m+n+2k} \rightarrow B_{m+2k} &\text{if } m \geq 0 \\ \; \; 0 &\text{if } m < 0 \end{cases} \]
gives an element in $\mathcal{E}xt_{\mathcal{C}}^{n+2k}(A, \widetilde{B}_{2k})$ that $\vartheta_{n}^{2k}(B)$ maps to $\Phi_{\bullet+n} + \widehat{\mathrm{Null}}_{\mathrm{Ch}(\mathcal{C})}(A[n]_{\bullet}, B_{\bullet})$. To prove that $\vartheta^n(B)$ is injective, let $x \in \varinjlim  \mathcal{E}xt_{\mathcal{C}}^{n+2k}(A, \widetilde{B}_{2k})$ with $\vartheta^n(B)(x) = 0$. According to Proposition~\ref{prop:liminab}, there exists $k \in \mathbb{N}_0$ and an element
\[  \psi_{\bullet+n+2k} \in \mathrm{Hom}_{\mathrm{Ch}(\mathcal{C})}(A[n+2k]_{\bullet}, B_{\bullet}^{+2k}) \]
that maps to $x$ in the direct limit and that lies in $\mathrm{Ker}(\vartheta_n^{2k}(B))$. One can choose the natural number $k$ such that the representative $\psi_{\bullet+n+2k}$ is nullhomotopic. As $x = 0$, $\vartheta^n(B)$ is injective and thus an isomorphism.
\end{proof}

The desired isomorphisms from the terms of the resolution construction to the ones of the hypercohomology construction is given by

\begin{defn}\label{defn:hyperresolcohom}
For every $n \in \mathbb{Z}$ there is an isomorphism
\[ \sigma^n := \vartheta^n \circ \zeta^n: \widehat{\mathrm{Ext}}_{\mathcal{C}}^n(A, -) \rightarrow \mathcal{E}xt_{Res, \mathcal{C}}^n(A, -) \rightarrow \widehat{\mathcal{E}xt}_{\mathcal{C}}^n(A, -) \]
where $\zeta^n$ is taken from Definition~\ref{defn:resolinchains} and $\vartheta^n$ is taken from Lemma~\ref{lem:vogelresol}.
\end{defn}

\subsection{Preliminaries to the isomorphisms of cohomological functors}

We require preliminary results to demonstrate that the isomorphisms 
\[ \vartheta^n(-): \mathcal{E}xt_{Res, \mathcal{C}}^n(A, -) \rightarrow \widehat{\mathcal{E}xt}_{\mathcal{C}}^n(A, -) \quad \text{and} \quad \sigma^n: \widehat{\mathrm{Ext}}_{\mathcal{C}}^n(A, -) \rightarrow \widehat{\mathcal{E}xt}_{\mathcal{C}}^n(A, -) \]
from Lemma~\ref{lem:vogelresol} and Definition~\ref{defn:hyperresolcohom} extend to isomorphisms of cohomological functors. We need to to prove that $\vartheta^n(-)$ is natural on the one hand while we need to develop a particular lift of an identity homomorphism on the other hand.

\begin{prop}\label{prop:thetanatural}
For every $n \in \mathbb{Z}$ the isomorphism $\vartheta^n(-)$ from Lemma~\ref{lem:vogelresol} is natural. That is, for every morphism $f: B \rightarrow C$ in $\mathcal{C}$ the square
\begin{center}
\begin{tikzcd}
  \mathcal{E}xt_{Res, \mathcal{C}}^n(A{,}\, B) \arrow[r, "\vartheta^n(B)"] \arrow[d, "\mathcal{E}xt_{Res, \mathcal{C}}^n(A{,}\, f)"] & \widehat{\mathcal{E}xt}_{\mathcal{C}}^n(A{,}\, B) \arrow[d, "\widehat{\mathcal{E}xt}_{\mathcal{C}}^n(A{,}\, f)"] \\
  \mathcal{E}xt_{Res, \mathcal{C}}^n(A{,}\, C) \arrow[r, "\vartheta^n(C)"] & \widehat{\mathcal{E}xt}_{\mathcal{C}}^n(A{,}\, C)
\end{tikzcd}
\end{center}
commutes.
\end{prop}

\begin{proof}
Let $(\widetilde{f}_k: \widetilde{M}_k \rightarrow \widetilde{N}_k)_{k \in \mathbb{N}_0}$ be a sequence of morphism as in Definition~\ref{defn:resolindmorph}. It suffices to consider the square
\begin{center}
\begin{tikzcd}
    \mathcal{E}xt_{\mathcal{C}}^{n+2k}(A{,} \, \widetilde{B}_{2k}) \arrow[rr, "\mathcal{E}xt_{\mathcal{C}}^{n+2k}(A{,} \, \widetilde{f}_{2k})"] \arrow[d, "\vartheta_n^{2k}(B)"] & & \mathcal{E}xt_{\mathcal{C}}^{n+2k}(A{,} \, \widetilde{C}_{2k}) \arrow[d, "\vartheta_n^{2k}(C)"] \\
    \widehat{\mathcal{E}xt}_{\mathcal{C}}^n(A{,} \, B) \arrow[rr, "\widehat{\mathcal{E}xt}_{\mathcal{C}}^n(A{,} \, f)"] & & \widehat{\mathcal{E}xt}_{\mathcal{C}}^n(A{,} \, C)
\end{tikzcd}
\end{center}
for any $k \in \mathbb{N}_0$ because Diagram~\ref{diag:thetadeltatriangle} commutes and the connecting homomorphisms $\Delta^{\bullet}: \mathcal{E}xt_{\mathcal{C}}^{\bullet}(A, -) \rightarrow \mathcal{E}xt_{\mathcal{C}}^{\bullet+1}(A, -)$ are natural. It commutes as we can lift $\widetilde{f}_{2k}: \widetilde{B}_{2k} \rightarrow \widetilde{C}_{2k}$ to a chain map $f_{\bullet}^{+2k}: B_{\bullet}^{+2k} \rightarrow C_{\bullet}^{+2k}$ given by
\[ f_m^{+2k} := \begin{cases} f_{m+2k}: B_{m+2k} \rightarrow C_{m+2k} &\text{if } m \geq 0 \\ \; \; 0 &\text{if } m < 0  \end{cases} \]
The above commutative squares form a direct system whose direct limit is the square in the statement of the proposition.
\end{proof}

We present a technical criterion of a lift of an identity homomorphism that we have promised above.

\begin{prop}\label{prop:syzygyislift}
Let $0 \rightarrow B \xrightarrow{f} C \xrightarrow{g} D \rightarrow 0$ be a short exact sequence in $\mathcal{C}$ and let $B_{\bullet}$, $D_{\bullet}$ be projective resolutions. Let $h^{\ast}: \widetilde{D}_1 \rightarrow B$ be a morphism lifting $\mathrm{id}_D: D \rightarrow D$ as in Diagram~\ref{diag:delta} and let $h_{\bullet+1}: D_{\bullet}^{+1} \rightarrow B_{\bullet}$ be a chain map lifting $h^{\ast}$ as in the Comparison Theorem. Denote by $\widetilde{h}_{k+1}^{\ast}: \widetilde{D}_{k+1} \rightarrow \widetilde{B}_k$ the induced homomorphism between the corresponding syzygies where $\widetilde{h}_1^{\ast} = h^{\ast}$. \\

Then there is a projective resolution $\underline{D}_{\bullet}'$ of $D$ such that for any $k \in \mathbb{N}_0$, $\widetilde{D}_k = \widetilde{\underline{D}}_k$ and there is a commutative diagram
\begin{equation}\label{diag:syzygyislift}
\begin{tikzcd}
    0 \arrow[r] & \widetilde{D}_{k+1} \arrow[r, "\iota_{k+1}"] \arrow[d, "\widetilde{h}_{k+1}^{\ast}"] & \underline{D}_k \arrow[r, "\pi_k"] \arrow[d, "h_k"] & \widetilde{D}_k \arrow[r] \arrow[d, "\mathrm{id}"] & 0 \\
    0 \arrow[r] & \widetilde{B}_k \arrow[r, "\widetilde{f}_k"] & \widetilde{C}_k \arrow[r, "\widetilde{g}_k"] & \widetilde{D}_k \arrow[r] & 0
\end{tikzcd}
\end{equation}
where the morphisms $\iota_{k+1}$, $\pi_k$ stem from $\underline{D}_{\bullet}$. In particular, any $\widetilde{h}_{k+1}^{\ast}$ is a lift of $\mathrm{id}_{\widetilde{D}_k}$ as in Diagram~\ref{diag:delta}.
\end{prop}

Let us explain the context in which we need this result. In Definition~\ref{defn:hyperconnmorph}, we have first constructed a connecting homomorphism $\widehat{\Delta}^n: \widehat{\mathcal{E}xt}_{\mathcal{C}}^n(A, D) \rightarrow \widehat{\mathcal{E}xt}_{\mathcal{C}}^{n+1}(A, \widetilde{D}_1)$ associated to the short exact sequence $0 \rightarrow \widetilde{D}_1 \rightarrow D_0 \rightarrow D \rightarrow 0$. Then we defined
\[ \widehat{\mathcal{E}xt}_{\mathcal{C}}^n(A, D) \xrightarrow{\widehat{\Delta}^n} \widehat{\mathcal{E}xt}_{\mathcal{C}}^{n+1}(A, \widetilde{D}_1) \xrightarrow{\widehat{\mathcal{E}xt}_{\mathcal{C}}^{n+1}(A, h^{\ast})} \widehat{\mathcal{E}xt}_{\mathcal{C}}^{n+1}(A, B) \]
to be the connecting homomorphism associated to $0 \rightarrow B \rightarrow C \rightarrow D \rightarrow 0$. Later, we require
\[ \mathcal{E}xt_{\mathcal{C}}^{n+2k}(A, \widetilde{D}_{2k}) \xrightarrow{\Delta^n} \mathcal{E}xt_{\mathcal{C}}^{n+1+2k}(A, \widetilde{D}_{2k+1}) \xrightarrow{\mathcal{E}xt_{\mathcal{C}}^{n+1+2k}(A, \widetilde{h}_{2k+1}^{\ast})} \mathcal{E}xt_{\mathcal{C}}^{n+1+2k}(A, \widetilde{B}_{2k}) \]
to be the connecting homomorphism associated to $0 \rightarrow \widetilde{B}_{2k} \rightarrow \widetilde{C}_{2k} \rightarrow \widetilde{D}_{2k} \rightarrow 0$. This is where Proposition~\ref{prop:syzygyislift} comes into play.

\begin{proof}[Proposition~\ref{prop:syzygyislift}]
We proceed by induction over $k \in \mathbb{N}_0$. If $k = 0$, then we take $\underline{D}_0 := D_0$ where Diagram~\ref{diag:delta} yields the desired diagram. Assume that $\underline{D}_0, \dots, \underline{D}_k$ have been defined and that there is a diagram of the required form for a number $k \geq 0$. In order to construct $\underline{D}_{k+1}$ and establish Diagram~\ref{diag:syzygyislift} in the inductive step, we consider two instance of Diagram~\ref{diag:horseshoe} from the Horseshoe Lemma. Adhering to the construction found in~\cite[p.~37--38]{wei94}, we explain how to obtain the first diagram:
\begin{equation}\label{diag:horseshoeback}
\begin{tikzcd}[row sep = scriptsize]
  {}
    & 0 \arrow[d]
    & & 0 \arrow[d]
    & & 0 \arrow[d]
    & {} \\
  0 \arrow[r]
    & \widetilde{D}_{k+2} \arrow[rr, "\iota_{k+2}''"] \arrow[dd, "\iota_{k+2}'"]
    & & \mathrm{Ker}(\varepsilon) \arrow[rr, "\pi_{k+1}''"] \arrow[dd]
    & & \widetilde{D}_{k+1} \arrow[r] \arrow[dd, "\iota_{k+1}"]
    & 0 \\ \\
  0 \arrow[r]
    & D_{k+1} \arrow[rr, "i"] \arrow[ddrr, "\iota_{k+1} \circ \pi_{k+1}'"] \arrow[dd, "\pi_{k+1}'"]
    & & D_{k+1} \oplus \underline{D}_k \arrow[rr, "p"] \arrow[dd, "\epsilon"]
    & & \underline{D}_k \arrow[r] \arrow[ddll, "\mathrm{id}"] \arrow[dd, "\pi_k"]
    & 0 \\ \\
  0 \arrow[r]
    & \widetilde{D}_{k+1} \arrow[rr, "\iota_{k+1}"] \arrow[d]
    & & \underline{D}_k \arrow[rr, "\pi_k"] \arrow[d]
    & & \widetilde{D}_k \arrow[r] \arrow[d]
    & 0 \\
  {}
    & 0
    & & 0
    & & 0
    & {}
\end{tikzcd}
\end{equation}
The morphisms $\pi_k$, $\iota_{k+1}$ stem from the projective resolution $(\underline{D}_m)_{m = 1}^k$ and the morphisms $\pi_{k+1}'$, $\iota_{k+2}'$ from $D_{\bullet}$. In particular, the morphisms $\iota_{k+2}'$ and $\iota_{k+1}$ from the top can be taken as kernels of the morphisms $\pi_{k+1}'$ and $\pi_k$. The morphisms of the middle row are the ones of a direct product and induce the morphisms of the top row. All rows and columns are exact. Because $\underline{D}_k$ and $D_{k+1} \oplus \underline{D}_k$ are projective, the middle column is a split exact sequence and $\mathrm{Ker}(\varepsilon)$ is also projective. Analogously, we construct with the morphisms $\widetilde{f}_k$ and $\widetilde{g}_k$ from Diagram~\ref{diag:syzygyislift} the second diagram of this form:
\begin{equation}\label{diag:horseshoefront}
\begin{tikzcd}[row sep = scriptsize]
  {}
    & 0 \arrow[d]
    & & 0 \arrow[d]
    & & 0 \arrow[d]
    & {} \\
  0 \arrow[r]
    & \widetilde{B}_{k+1} \arrow[rr, "\widetilde{f}_{k+1}"] \arrow[dd]
    & & \widetilde{C}_{k+1} \arrow[rr, "\widetilde{g}_{k+1}"] \arrow[dd]
    & & \widetilde{D}_{k+1} \arrow[r] \arrow[dd, "\iota_{k+1}"]
    & 0 \\ \\
  0 \arrow[r]
    & B_k \arrow[rr, "i = f_k"] \arrow[ddrr, "\widetilde{f}_k \circ \check{\pi}_k"] \arrow[dd, "\check{\pi}_k"]
    & & B_k \oplus \underline{D}_k \arrow[rr, "p = g_k"] \arrow[dd, "\eta"]
    & & \underline{D}_k \arrow[r] \arrow[ddll, "h_k"] \arrow[dd, "\pi_k"]
    & 0 \\ \\
  0 \arrow[r]
    & \widetilde{B}_k \arrow[rr, "\widetilde{f}_k"] \arrow[d]
    & & C_k \arrow[rr, "\widetilde{g}_k"] \arrow[d]
    & & \widetilde{D}_k \arrow[r] \arrow[d]
    & 0 \\
  {}
    & 0
    & & 0
    & & 0
    & {}
\end{tikzcd}
\end{equation}
The triangle at the bottom right corner commutes by Diagram~\ref{diag:syzygyislift}. We incorporate the lower halves of Diagram~\ref{diag:horseshoefront} and Diagram~\ref{diag:horseshoeback} into the front and back side of the diagram \\

\begin{equation}\label{diag:doesntcommingen}
\begin{tikzcd}
  & {\scriptstyle D_{k+1}} \arrow[rr, "{\scriptstyle i_1^1}"] \arrow[ddrr, bend right = 15, "{\scriptstyle \iota_{k+1} \circ \pi_{k+1}'}" near end] \arrow[dl, "{\scriptstyle h_{k+1}^{\ast}}"] \arrow[dd, "{\scriptstyle \pi_{k+1}'}" near end]
    & & {\scriptstyle D_{k+1} \oplus \underline{D}_k} \arrow[rr, "{\scriptstyle p_2^1}"] \arrow[dl, "{\scriptstyle h_{k+1}^{\ast} \oplus \mathrm{id}}"] \arrow[dd, shift left = 2mm, "{\scriptstyle \epsilon}" near start]
    & & {\scriptstyle \underline{D}_k} \arrow[dl, "{\scriptstyle \mathrm{id}}" near end] \arrow[ddll, bend left = 32, "{\scriptstyle \mathrm{id}}" near start] \arrow[dd, "{\scriptstyle \pi_k}"] \\
  {\scriptstyle B_k} \arrow[rr, crossing over, "{\scriptstyle i_1^2}" near end] \arrow[dd, "{\scriptstyle \check{\pi}_k}"]
    & & {\scriptstyle B_k \oplus \underline{D}_k} \arrow[rr, crossing over, "{\scriptstyle p_2^2}" near end]
    & & {\scriptstyle \underline{D}_k} & \\
  & {\scriptstyle \widetilde{D}_{k+1}} \arrow[rr, "{\scriptstyle \iota_{k+1}}" near start] \arrow[dl, "{\scriptstyle \widetilde{h}_{k+1}^{\ast}}"]
    & & {\scriptstyle \underline{D}_k} \arrow[rr, "{\scriptstyle \pi_k}" near end] \arrow[dl, "{\scriptstyle h_k}" near start]
    & & {\scriptstyle \widetilde{D}_k} \arrow[dl, "{\scriptstyle \mathrm{id}}"] \\
  {\scriptstyle \widetilde{B}_k} \arrow[rr, "{\scriptstyle \widetilde{f}_k}"]
    & & {\scriptstyle C_k} \arrow[rr, "{\scriptstyle \widetilde{g}_k}" near end] \arrow[from = uull, bend right = 20, crossing over, "{\scriptstyle \widetilde{f}_k \circ \check{\pi}_k}" near end] \arrow[from = uu, crossing over, "{\scriptstyle \eta}" near end] \arrow[from = uurr, bend left = 25, crossing over, "{\scriptstyle h_k}"]
    & & {\scriptstyle \widetilde{D}_k} \arrow[from = uu, crossing over, "{\scriptstyle \pi_k}" near end] &
\end{tikzcd}
\end{equation}
The bottom side stems from Diagram~\ref{diag:syzygyislift}, which commutes by our induction hypothesis. The left hand side arises from the chain map $h_{\bullet}^{\ast}: D_{\bullet}^{+1} \rightarrow B_{\bullet}$. The morphism $h_{k+1}^{\ast} \oplus \mathrm{id}$ is constructed as in~\cite[p.~251]{mac95} and renders the top side commutative. In particular, all squares on the outside of Diagram~\ref{diag:doesntcommingen} commute. The slanted squares in the interior give rise to the commuting diagram
\begin{center}
\begin{tikzcd}
    D_{k+1} \arrow[rr, "\iota_{k+1} \circ \pi_{k+1}'"] \arrow[dd, "h_{k+1}^{\ast}"] & & \underline{D}_k \arrow[dd, "h_k"] & & \underline{D}_k  \arrow[dd, "\mathrm{id}"] \arrow[ll, "\mathrm{id}"] \\ \\
    B_k \arrow[rr, "\widetilde{f}_k \circ \check{\pi}_k"] & & C_k & & \underline{D}_k \arrow[ll, "h_k"]
\end{tikzcd}
\end{center}
We can use it together with Diagram~\ref{diag:horseshoeback}, Diagram~\ref{diag:horseshoefront} and the content of~\cite[p.~251]{mac95} to prove that that the middle square in the interior of Diagram~\ref{diag:doesntcommingen} also commutes. Therefore, all squares in Diagram~\ref{diag:doesntcommingen} commute. Its top side is connected to its bottom side by epimorphisms according to Diagram~\ref{diag:horseshoeback} and~\ref{diag:horseshoefront}. We take kernels to obtain the commuting diagram
\begin{center}
\begin{tikzcd}
    0 \arrow[r] & \widetilde{D}_{k+2} \arrow[r, "\iota_{k+2}''"] \arrow[d, "\widetilde{h}_{k+2}^{\ast}"] & \mathrm{Ker}(\varepsilon) \arrow[r, "\pi_{k+1}''"] \arrow[d, "H"] & \widetilde{D}_{k+1} \arrow[r] \arrow[d, "\mathrm{id}"] & 0 \\
    0 \arrow[r] & \widetilde{B}_{k+1} \arrow[r, "\widetilde{f}_{k+1}"] & \widetilde{C}_{k+1} \arrow[r, "\widetilde{g}_{k+1}"] & \widetilde{D}_{k+1} \arrow[r] & 0
\end{tikzcd}
\end{center}
where the rows are short exact sequences. Having seen that $\mathrm{Ker}(\varepsilon)$ is projective, we complete the inductive step by setting $\underline{D}_{k+1} := \mathrm{Ker}(\varepsilon)$, $\iota_{k+2} : = \iota_{k+2}''$, $\pi_{k+1} := \pi_{k+1}''$ and $h_{k+1} := H$.
\end{proof}

\subsection{Isomorphisms of cohomological functors}

We establish the first isomorphism of cohomological functors to the terms of the hypercohomology construction.

\begin{lem}\label{lem:hyperresol}
For any short exact sequence $0 \rightarrow B \rightarrow C \rightarrow D \rightarrow 0$ in $\mathcal{C}$ the square
\begin{center}
\begin{tikzcd}
    \mathcal{E}xt_{Res, \mathcal{C}}^n(A{,} \, D) \arrow[r, "\vartheta^n"] \arrow[d, "\widetilde{\Delta}^n"] & \widehat{\mathcal{E}xt}_{\mathcal{C}}^n(A{,} \, D) \arrow[d, "\widehat{\Delta}^n"] \\
    \mathcal{E}xt_{Res, \mathcal{C}}^{n+1}(A{,} \, B) \arrow[r, "\vartheta^{n+1}"] & \widehat{\mathcal{E}xt}_{\mathcal{C}}^{n+1}(A{,} \, B)
\end{tikzcd}
\end{center}
commutes for every $n \in \mathbb{Z}$. In particular, $\vartheta^{\bullet}: \mathcal{E}xt_{Res, \mathcal{C}}^{\bullet}(A, -) \rightarrow \widehat{\mathcal{E}xt}_{\mathcal{C}}^{\bullet}(A, -)$ is an isomorphism of cohomological functors and $(\widehat{\mathcal{E}xt}_{\mathcal{C}}^{\bullet}(A, -), \widehat{\Delta}^{\bullet})$ forms a cohomological functor.
\end{lem}

\begin{proof}
It suffices to prove that the square in the statement of the theorem commutes since $\vartheta^n$ is already a natural isomorphism by Proposition~\ref{prop:thetanatural}. Let $k \in \mathbb{N}_0$, $N := n+2k$ and $K := 2k$. Then we need to argue that the below diagram commutes.
\begin{equation}\label{diag:unfoldedprism}
\begin{tikzcd}
    {\scriptscriptstyle \mathcal{E}xt_{\mathcal{C}}^N(A{,} \, \widetilde{D}_K)} \arrow[rr, "{\scriptscriptstyle \Delta^N}" near end] \arrow[rrrrr, bend right = 12, "\Delta^N"] \arrow[dd, shift right = 4mm, "{\scriptscriptstyle \Delta^N}"] & & {\scriptscriptstyle \mathcal{E}xt_{\mathcal{C}}^{N+1}(A{,} \, \widetilde{D}_{K+1})} \arrow[rrr, "{\scriptscriptstyle \mathcal{E}xt_{\mathcal{C}}^{N+1}(A{,} \, \widetilde{h}_{K+1}^{\ast})}" near end] \arrow[dd, shift right = 4mm, crossing over, "{\scriptscriptstyle \Delta^{N+1}}" near end] & & & {\scriptscriptstyle \mathcal{E}xt_{\mathcal{C}}^{N+1}(A{,} \, \widetilde{B}_K)} \arrow[dd, shift right = 4mm, "{\scriptscriptstyle \Delta^{N+1}}"] \arrow[dddddd, shift left = 6mm, bend left = 25, "{\scriptscriptstyle \vartheta_{n+1}^K(B)}"] \\ \\
    {\scriptscriptstyle \mathcal{E}xt_{\mathcal{C}}^{N+1}(A{,} \, \widetilde{D}_{K+1})} \arrow[rr, "{\scriptscriptstyle -\Delta^{N+1}}" near end] \arrow[rrrrr, bend right = 12, "-\Delta^{N+1}"] \arrow[dd, shift right = 4mm, "{\scriptscriptstyle \Delta^{N+1}}"] & & {\scriptscriptstyle \mathcal{E}xt_{\mathcal{C}}^{N+2}(A{,} \, \widetilde{D}_{K+2})} \arrow[rrr, "{\scriptscriptstyle \mathcal{E}xt_{\mathcal{C}}^{N+2}(A{,} \, \widetilde{h}_{K+2}^{\ast})}" near end] \arrow[dd, shift right = 4mm, crossing over, "{\scriptscriptstyle \Delta^{N+2}}" near end] & & & {\scriptscriptstyle \mathcal{E}xt_{\mathcal{C}}^{N+2}(A{,} \, \widetilde{B}_{K+1})} \arrow[dd, shift right = 4mm, "{\scriptscriptstyle \Delta^{N+2}}"] \\ \\
    {\scriptscriptstyle \mathcal{E}xt_{\mathcal{C}}^{N+2}(A{,} \, \widetilde{D}_{K+2})} \arrow[rr, "{\scriptscriptstyle \Delta^{N+2}}" near end] \arrow[rrrrr, bend right = 12, "\Delta^{N+2}"] \arrow[dd, shift right = 4mm, "{\scriptscriptstyle \vartheta_n^{K+2}(D)}"] & & {\scriptscriptstyle \mathcal{E}xt_{\mathcal{C}}^{N+3}(A{,} \, \widetilde{D}_{K+3})} \arrow[rrr, "{\scriptscriptstyle \mathcal{E}xt_{\mathcal{C}}^{N+2}(A{,} \, \widetilde{h}_{K+3}^{\ast})}" near end] \arrow[dd, shift right = 4mm, crossing over, "{\scriptscriptstyle \vartheta_{n+1}^{K+2}(\widetilde{D}_1)}" near end] & & & {\scriptscriptstyle \mathcal{E}xt_{\mathcal{C}}^{N+3}(A{,} \, \widetilde{B}_{K+2})} \arrow[dd, shift right = 4mm, "{\scriptscriptstyle \vartheta_{n+1}^{K+2}(B)}"] \\ \\
    {\scriptscriptstyle \widehat{\mathcal{E}xt}_{\mathcal{C}}^n(A{,} \, D)} \arrow[rr, "{\scriptscriptstyle \widehat{\Delta}^n}" near end] \arrow[rrrrr, bend right = 12, "\widehat{\Delta}^n"] \arrow[from = uuuuuu, shift left = 6mm, bend left = 25, crossing over, "{\scriptscriptstyle \vartheta_n^K(D)}"] & & {\scriptscriptstyle \widehat{\mathcal{E}xt}_{\mathcal{C}}^{n+1}(A{,} \, \widetilde{D}_1)} \arrow[rrr, "{\scriptscriptstyle \widehat{\mathcal{E}xt}_{\mathcal{C}}^{n+1}(A{,} \, h^{\ast})}" near end] \arrow[from = uuuuuu, shift left = 7mm, bend left = 25, crossing over, "{\scriptscriptstyle \vartheta_{n+1}^K(\widetilde{D}_1)}"] & & & {\scriptscriptstyle \widehat{\mathcal{E}xt}_{\mathcal{C}}^{n+1}(A{,} \, B)}
\end{tikzcd}
\end{equation}
We explain how to construct it. According to Definition~\ref{defn:hyperconnmorph}, the connecting homomorphism $\widehat{\Delta}^n: \widehat{\mathcal{E}xt}_{\mathcal{C}}^n(A, D) \rightarrow \widehat{\mathcal{E}xt}_{\mathcal{C}}^{n+1}(A, B)$ associated to $0 \rightarrow B \rightarrow C \rightarrow D \rightarrow 0$ is rendered in the very bottom row. The top left and middle left square arise from Diagram~\ref{diag:anticomm} and Lemma~\ref{lem:ordinaryhypercohom}. More specifically, their connecting homomorphisms are associated to the short exact sequences at the margins of Diagram~\ref{diag:horseshoeback}. It follows from Definition~\ref{defn:hyperconnmorph} and Lemma~\ref{lem:vogelresol} that the bottom left square and thus, the entire left hand side commutes. The induced homomorphisms in the right hand side are taken from Proposition~\ref{prop:syzygyislift}. More specifically, the homomorphism $h^{\ast}: \widetilde{D}_1 \rightarrow B$ is a lift of $\mathrm{id}_D: D \rightarrow D$ and the homomorphisms $\widetilde{h}_{K+1}^{\ast}$ stem from the chain map $h_{\bullet+1}^{\ast}: D_{\bullet}^{+1} \rightarrow B_{\bullet}$ lifting $h^{\ast}$ as in the Comparison Theorem. By Proposition~\ref{prop:thetanatural}, the entire right hand side commutes. According to Proposition~\ref{prop:syzygyislift}, the homomorphism
\[ \mathcal{E}xt_{\mathcal{C}}^N(A, \widetilde{D}_K) \xrightarrow{\Delta^N} \mathcal{E}xt_{\mathcal{C}}^{N+1}(A, \widetilde{D}_{K+1}) \xrightarrow{\mathcal{E}xt_{\mathcal{C}}^{N+1}(A, \widetilde{h}_{K+1}^{\ast})} \mathcal{E}xt_{\mathcal{C}}^{N+1}(A, \widetilde{B}_K) \]
at the very top of the diagram is the connecting homomorphism $\Delta^N$ associated to the short exact sequence $0 \rightarrow \widetilde{B}_K \rightarrow \widetilde{C}_K \rightarrow \widetilde{D}_K \rightarrow 0$. Analogously, the middle row to the bottom represents the connecting homomorphism $\Delta^{N+2}$. Therefore, the homomorphisms at the margins of Diagram~\ref{diag:unfoldedprism} together with the homomorphisms in its lower third form a direct system of commuting squares. We obtain the square in statement of the theorem when passing to the direct limi.
\end{proof}

\begin{rem}
For any short exact sequence $0 \rightarrow B \rightarrow C \rightarrow D \rightarrow 0$ in $\mathcal{C}$ and any $n \in \mathbb{Z}$, Lemma~\ref{lem:hyperresol} implies that $\widehat{\Delta}^n: \widehat{\mathcal{E}xt}_{\mathcal{C}}^n(A, D) \rightarrow \widehat{\mathcal{E}xt}_{\mathcal{C}}^{n+1}(A, B)$ is well defined in the sense that it does not depend on a choice of a morphism $h^{\ast}: \widetilde{D}_1 \rightarrow B$ as in Definition~\ref{defn:hyperconnmorph}. We have not been able to prove this directly.
\end{rem}

From Definition~\ref{defn:resolinchains} and Lemma~\ref{lem:hyperresol} we deduce

\begin{thm}\label{thm:hyperresol}
The hypercohomology construction $(\widehat{\mathcal{E}xt}_{\mathcal{C}}^{\bullet}(A, -), \widehat{\Delta}^{\bullet})$ forms a Mislin completion of $(\mathrm{Ext}_{\mathcal{C}}^{\bullet}(A, -), \delta^{\bullet})$. More specifically,
\[ \sigma^{\bullet} = \vartheta^{\bullet} \circ \zeta^{\bullet}: \big(\widehat{\mathrm{Ext}}_{\mathcal{C}}^{\bullet}(A, -), \widehat{\delta}^{\bullet}\big) \rightarrow \big(\mathcal{E}xt_{Res, \mathcal{C}}^{\bullet}(A, -), \widetilde{\Delta}^{\bullet}\big) \rightarrow \big(\widehat{\mathcal{E}xt}_{\mathcal{C}}^{\bullet}(A, -), \widehat{\Delta}^{\bullet}\big) \]
is an isomorphism of cohomological functors where the isomorphisms $\sigma^n$ are taken from Definition~\ref{defn:hyperresolcohom}
\end{thm}

In~\cite[p.~21--22]{guo23}, S.\! Guo and L.\! Liang establish an isomorphism $\widehat{\mathcal{E}xt}_{\mathcal{C}}^n(A, B) \rightarrow BC_{\mathcal{C}}^n(A, B)$ from the terms of the hypercohomology construction to the ones of the naïve construction for every $n \in \mathbb{Z}$. As it is used in the construction of Yoneda products in~\cite[Theorem~6.6]{ghe24}, we present it below and prove that it specifically forms an isomorphism of cohomological functors.

\begin{defn}\label{defn:hypernaive}
(\cite[p.~21--22]{guo23}) For any $A, B \in \mathrm{obj}(\mathcal{C})$ and $n \in \mathbb{Z}$ let the map
\[ \rho^n: \widehat{\mathcal{E}xt}_{\mathcal{C}}^n(A, B) \rightarrow BC_{\mathcal{C}}^n(A, B) \]
be given as follows. If $x = \varphi_{\bullet+n}+\widehat{\mathrm{Null}}_{\mathrm{Ch}(\mathcal{C})}(A[n]_{\bullet}, B_{\bullet})$ is an element in $\widehat{\mathcal{E}xt}_{\mathcal{C}}^n(A, B)$, then there exists $K \in \mathbb{N}_0$ such that $(\varphi_{m+n}: A_{m+n} \rightarrow B_m)_{m \geq 2K}$ forms a chain map. Denote by $\widetilde{\varphi}_{2K+n}: \widetilde{A}_{2K+n} \rightarrow \widetilde{B}_{2K}$ the morphism induced between the corresponding syzygies. We set $\rho^n(x) \in BC_{\mathcal{C}}^n(A, B)$ to be the element to which the element $\widetilde{\varphi}_{2K}+\mathcal{P}_{\mathcal{C}}(\widetilde{A}_{2K+n}, \widetilde{B}_{2K})$ in $[\widetilde{A}_{2K+n}, B_{2K}]_{\mathcal{C}}$ is mapped in the direct limit.
\end{defn}

\begin{lem}\label{lem:hypernaive}
The maps $\rho^n$ form an isomorphism of cohomological functors such that
\[ \rho^{\bullet} = \beta^{\bullet} \circ (\sigma^{\bullet})^{-1}: (\widehat{\mathcal{E}xt}_{\mathcal{C}}^{\bullet}(A, -), \widehat{\Delta}^{\bullet}) \rightarrow (\mathrm{Ext}_{Res, \mathcal{C}}^n(A, -), \widetilde{\delta}^{\bullet}) \rightarrow (BC_{\mathcal{C}}^{\bullet}(A, -), \tau^{\bullet}) \]
where $\beta^{\bullet}$ is taken from Theorem~\ref{thm:resolbenson} and $\sigma^{\bullet}$ from Theorem~\ref{thm:hyperresol}.
\end{lem}

\begin{proof}
Since $\beta^{\bullet}$ and $\sigma^{\bullet}$ are already isomorphisms of cohomological functors, it suffices to show that $\rho^n = \beta^n \circ (\sigma^{n})^{-1}$ for every $n \in \mathbb{Z}$. This follows from Definition~\ref{defn:hyperresolcohom} and Definition~\ref{defn:hypernaive}.
\end{proof}

We close our account on constructions of Mislin completions by posing a question. At the start of Subsection~\ref{subsec:hyperindconn} we have seen that any completed Ext-group $\widehat{\mathcal{E}xt}_{\mathcal{C}}^n(A, B)$ can be defined as the $n^{\text{th}}$ cohomology group of the Vogel complex $\mathrm{Vog}_{\mathcal{C}}^n(A_{\bullet}, B_{\bullet})_{\bullet}$. Therefore, one can define connecting homomorphisms for the hypercohomology construction in a different manner as the following lemma demonstrates.

\begin{lem}
(\cite[Proposition~4.8]{guo23}) Let $0 \rightarrow B \rightarrow C \rightarrow D \rightarrow 0$ be a short exact sequence in $\mathcal{C}$ and let $0 \rightarrow B_{\bullet} \rightarrow C_{\bullet} \rightarrow D_{\bullet} \rightarrow 0$ be a short exact sequence of chain complexes obtained by the Horseshoe Lemma~\cite[p.~37]{wei94}. According to the Snake Lemma~\cite[p.~11--12]{wei94}, there are associated connecting homomorphisms
\[ \overline{\delta}^n: \widehat{\mathcal{E}xt}_{\mathcal{C}}^n(A, D) \rightarrow \widehat{\mathcal{E}xt}_{\mathcal{C}}^{n+1}(A, B) \quad \text{and} \quad \overline{\delta}^n: \widehat{\mathcal{E}xt}_{\mathcal{C}}^n(B, A) \rightarrow \widehat{\mathcal{E}xt}_{\mathcal{C}}^{n+1}(D, A) \, . \]
In particular, $\big(\widehat{\mathcal{E}xt}_{\mathcal{C}}^{\bullet}(A, -), \overline{\delta}^{\bullet}\big)$ and $\big(\widehat{\mathcal{E}xt}_{\mathcal{C}}^{\bullet}(-, A), \overline{\delta}^{\bullet}\big)$ form cohomological functors.
\end{lem}

\begin{ques}\label{ques:comparewithguoliang}
Does the connecting homomorphism $\widehat{\Delta}^{\bullet}$ constructed in Definition~\ref{defn:hyperconnmorph} agree with $\overline{\delta}^{\bullet}$? More specifically, is $\big(\widehat{\mathcal{E}xt}_{\mathcal{C}}^{\bullet}(A, -), \overline{\delta}^{\bullet}\big)$ a Mislin completion of the Ext-functors $\big(\mathrm{Ext}_{\mathcal{C}}^{\bullet}(A, -), \delta^{\bullet}\big)$?
\end{ques}

\section{Complete cohomology in condensed mathematics}\label{sec:condmaths}

As a noteworthy feature, the uniform theory developed in this paper is applicable to condensed mathematics, a powerful novel theory developed by D.\! Clausen and P.\! Scholze. Thus, we first provide an outline of the latter theory. Then we introduce Tate cohomology for all $T1$ topological groups using condensed mathematics. We also present an example of completed Ext-functors in condensed mathematics that can be established only through our theory and through no previous framework. \\

The building blocks of condensed mathematics are condensed sets that we define following D.\! Clausen and P.\! Scholze's account~\cite[Lecture~I and~II]{sch22}. In short, condensed sets can be thought as particular sheaves of sets. The reason for using sheaves is because they have excellent category theoretic properties that are fundamental for group cohomology. Further below we provide a plain and straightforward definition of $\kappa$-condensed sets after introducing the relevant Grothendieck sites of our sheaves. A profinite space is an inverse limit of finite discrete spaces, meaning that it can be assembled from finite sets in a particular manner. Let $\mathbf{Pro}$ denote the category of profinite spaces. A finite collection of (continuous) maps $\lbrace f_i: S_i \rightarrow S \rbrace_{i = 1}^n$ in $\mathbf{Pro}$ is called a covering if the induced map $\bigsqcup_{i = 1}^n S_i \rightarrow S$ is surjective. Although this turns $\mathbf{Pro}$ into a Grothendieck site, it is not well suited for sheaf theoretic purposes as it is not essentially small. A cardinal number $\kappa$ is called a strong limit cardinal if for every $\lambda < \kappa$ we also have $2^{\lambda} < \kappa$. For an uncountable strong limit cardinal $\kappa$ we define the site of $\kappa$-small profinite spaces $\mathbf{Pro}_{\kappa}$ as the full subcategory of $\mathbf{Pro}$ consisting of all profinite spaces with less than $\kappa$ clopen subsets. Note that the Grothendieck site $\mathbf{Pro}_{\kappa}$ is essentially small and thus suited for sheaf theoretic purposes. Then a $\kappa$-condensed set is a sheaf of sets on $\mathbf{Pro}_{\kappa}$. In more down-to-earth terms, a (contravariant) functor
\[ X: \mathbf{Pro}_{\kappa}^{\mathrm{op}} \rightarrow \mathbf{Set} \]
is a $\kappa$-condensed set if the following two conditions are satisfied.
\begin{enumerate}
    \item For every finite collection $\lbrace S_i \rbrace_{i = 1}^n$ of objects in the category $\mathbf{Pro}_{\kappa}$ there is a bijection $X(\bigsqcup_{i = 1}^n S_i) \rightarrow \prod_{i = 1}^n X(S_i)$.
    \item Let $f: T \rightarrow S$ be a (continuous) surjection in $\mathbf{Pro}_{\kappa}$ and denote the projections from the fiber product by $p_1, p_2: T \times_S T \rightarrow T$. Then the function $X(f): X(S) \rightarrow X(T)$ is injective with image
\[ \mathrm{Im}(X(f)) = \lbrace x \in X(T) \mid X(p_1)(x) = X(p_2)(x) \in X(T \times_S T) \rbrace \, . \]
\end{enumerate}
One can think of $X$ as encoding a topological space and of $X(S)$ as the continuous maps ``$S \rightarrow X$'' from a $\kappa$-profinite space. If $\kappa < \kappa'$ is another strong limit cardinal, then there is a functor
\[ \mathrm{Cond}_{\kappa}(\mathbf{Set}) \rightarrow \mathrm{Cond}_{\kappa'}(\mathbf{Set}) \, . \]
It is fully faithful because we have used strong limit cardinals. This means that one can extend any $\kappa$-condensed set to a $\kappa'$-condensed set. If $K$ denotes the class of uncountable strong limit cardinals, then the category of condensed sets is defined as the colimit category
\[ \mathrm{Cond}(\mathbf{Set}) := \varinjlim_{\kappa \in K} \mathrm{Cond}_{\kappa}(\mathbf{Set}) \, . \]
Hence, a condensed set is an equivalence class of $\kappa$-condensed sets and thus, an equivalence class of sheaves of sets. Although $\mathrm{Cond}(\mathbf{Set})$ does not form the category of sheaves over any site~\cite[Remark~2.12]{sch19}, it nevertheless inherits good category theoretic properties from the categories $\mathrm{Cond}_{\kappa}(\mathbf{Set})$~\cite[p.~13]{sch22}. \\

We define condensed groups and modules, we demonstrate how topological groups and modules can be turned into condensed objects and we present relevant category theoretic properties. By~\cite[p.~7]{sch19} condensed groups are defined as equivalence classes of $\kappa$-condensed groups, meaning sheaves of groups on the site $\mathbf{Pro}_{\kappa}$. One defines other condensed algebraic structures such as condensed rings and condensed modules analogously. According to~\cite[p.~15/16 and Proposition~1.7]{sch19} we can establish the following relation with topological spaces. If $T1\text{-}\mathbf{Top}$ denotes the category of $T1$ topological spaces, then there is a faithful functor $T1\text{-}\mathbf{Top} \rightarrow \mathrm{Cond}(\mathbf{Set})$ taking any $Y$ to its condensate
\[ \underline{Y}: \mathbf{Pro}_{\kappa}^{\mathrm{op}} \rightarrow \mathbf{Set}, \; S \mapsto \lbrace f: S \rightarrow Y \text{ continuous} \rbrace \]
where $\kappa$ is chosen to be a sufficiently large strong limit cardinal. This functor does not depend on $\kappa$ as evidenced by the proofs of Proposition~2.9 and Proposition~2.15 in~\cite{sch19}. Thus, one can replace any $T1$ topological space by its condensate and only work with the latter. Condensates of $T1$ topological groups (resp.\! rings, modules etc) are condensed groups (resp.\! rings, modules etc)~\cite[p.~8]{sch19}. In the category of condensed abelian groups $\mathrm{Cond}(\mathbf{Ab})$ all limits and colimits exist where arbitrary products, arbitrary direct sums and direct limits are exact~\cite[Theorem~1.10]{sch19}. According to~\cite[p.~12]{sch19}, $\mathrm{Cond}(\mathbf{Ab})$ has enough projectives. Following the source, these assertions also hold for the category of condensed $\mathcal{R}$-modules $\mathrm{Cond}(Mod(\mathcal{R}))$ where $\mathcal{R}$ is a condensed ring. This is relevant to our theory as it constructs Mislin completions for functors $T^{\bullet}: \mathcal{C} \rightarrow \mathcal{D}$ where $\mathcal{C}$ has enough projectives and in $\mathcal{D}$ all countable direct limits exists and are exact. Therefore, one can construct Mislin completions using condensed modules. \\

We construct free condensed modules in order to define the above mentioned completed Ext-functors that can be established only through our theory. Let $S$ be a projective object in $\mathbf{Pro}$. Then the free condensed $\mathcal{R}$-module $\mathcal{R}[\underline{S}]$ is the sheafification of the presheaf of $\mathcal{R}$-modules
\[ \mathbf{Pro}_{\kappa}^{\mathrm{op}} \rightarrow \mathbf{Ab}, T \mapsto \mathcal{R}(T)[\underline{S}(T)] \]
where $\mathcal{R}(T)[\underline{S}(T)]$ denotes the free $\mathcal{R}(T)$-module over the set $\underline{S}(T)$ and $\kappa$ is a sufficiently large strong limit cardinal~\cite[p.~12]{sch19}. Since the Grothendieck site $\mathbf{Pro}_{\kappa}$ is essentially small, sheafification exists and the construction of $\mathcal{R}[\underline{S}]$ is independent of $\kappa$ by the proof of~\cite[Proposition~2.9]{sch19}. The condensed modules $\mathcal{R}[\underline{S}]$ are compact projective generators, meaning that for every condensed $\mathcal{R}$-module $A$ there is a collection of projective profinite spaces $S_i$ such that there is an epimorphism $\bigoplus_{i \in I} \mathcal{R}[\underline{S_i}] \rightarrow A$~\cite[p.~12]{sch19}. For condensed $\mathcal{R}$-modules $A$, $B$ one does not only have the ``usual'' Hom-set
\[ \mathrm{Hom}_{\mathrm{Cond}(Mod(\mathcal{R}))}(A, B) \in \mathrm{obj}(\mathbf{Ab}) \, , \]
but also an internal Hom-set which is constructed in~\cite[p.~13]{sch19}. For any condensed $\mathcal{R}$-module $C$ the condensed abelian group $C \otimes_{\mathcal{R}} A$ is the sheafification of
\[ \mathbf{Pro}_{\kappa}^{\mathrm{op}} \rightarrow \mathbf{Ab}, T \mapsto C(T) \otimes_{\mathcal{R}(T)} A(T) \, . \]
Then define the internal Hom-set
\[ \underline{\mathrm{Hom}}_{\mathrm{Cond}(Mod(\mathcal{R}))}(A, B) \in \mathrm{obj}(\mathrm{Cond}(\mathbf{Ab})) \]
by
\[ \mathbf{Pro}_{\kappa}^{\mathrm{op}} \rightarrow \mathbf{Ab}, T \mapsto \mathrm{Hom}_{\mathrm{Cond}(Mod(\mathcal{R}))}(A \otimes_{\mathcal{R}} \mathcal{R}[\underline{T}], B) \, . \]
By definition, it satisfies the tensor-Hom adjunction
\[ \mathrm{Hom}_{\mathrm{Cond}(Mod(\mathcal{R}))}(C, \underline{\mathrm{Hom}}_{\mathrm{Cond}(Mod(\mathcal{R}))}(A, B)) \cong \mathrm{Hom}_{\mathrm{Cond}(Mod(\mathcal{R}))}(C \otimes_{\mathcal{R}} A, B) \, . \]
Because the category $\mathrm{Cond}(Mod(\mathcal{R}))$ has enough projectives, we can define the unenriched Ext-functors
\[ \mathrm{Ext}_{\mathrm{Cond}(Mod(\mathcal{R}))}^{\bullet}(A, -): \mathrm{Cond}(Mod(\mathcal{R})) \rightarrow \mathbf{Ab} \]
and the internal Ext-functors
\[ \underline{\mathrm{Ext}}_{\mathrm{Cond}(Mod(\mathcal{R}))}^{\bullet}(A, -): \mathrm{Cond}(Mod(\mathcal{R})) \rightarrow \mathrm{Cond}(\mathbf{Ab}) \]
to be the derived functors of $\mathrm{Hom}_{\mathrm{Cond}(Mod(\mathcal{R}))}(A, -)$ and $\underline{\mathrm{Hom}}_{\mathrm{Cond}(Mod(\mathcal{R}))}(A, -)$. \\

These Ext-functors can be used to define cohomology of a condensed group $\mathcal{G}$. Following~\cite[p.~2--3]{ans20} one can form the condensed group ring $\mathcal{R}[\mathcal{G}]$ such that the category $\mathrm{Cond}(Mod(\mathcal{R}[\mathcal{G}]))$ is equivalent to the subcategory of $\mathrm{Cond}(\mathcal{R})$ of all $\mathcal{R}$-modules $M$ with a $\mathcal{G}$-action $\mathcal{G} \times M \rightarrow M$. By~\cite[p.~5/8]{ans20}, one can define condensed unenriched group cohomology $H_{\mathcal{R}}^{\bullet}(\mathcal{G}, -)$ and the condensed internal group cohomology $\underline{H}_{\mathcal{R}}^{\bullet}(\mathcal{G}, -)$ as the Ext-functors $\mathrm{Ext}_{\mathrm{Cond}(Mod(\mathcal{R}[\mathcal{G}]))}^{\bullet}(\mathcal{R}, -)$ and $\underline{\mathrm{Ext}}_{\mathrm{Cond}(Mod(\mathcal{R}[\mathcal{G}]))}^{\bullet}(\mathcal{R}, -)$. By the above, we conclude

\begin{thm}\label{thm:condgroupcohom}
Let $\mathcal{R}$ be a condensed ring and $A$ a condensed $\mathcal{R}$-module. Then there exist completed condensed unenriched Ext-functors $\widehat{\mathrm{Ext}}_{\mathrm{Cond}(Mod(\mathcal{R}))}^{\bullet}(A, -)$ and completed condensed internal Ext-functors $\widehat{\underline{\mathrm{Ext}}}_{\mathrm{Cond}(Mod(\mathcal{R}))}^{\bullet}(A, -)$. If $\mathcal{G}$ is a condensed group, then there exists complete condensed unenriched group cohomology $\widehat{H}_{\mathcal{R}}^{\bullet}(\mathcal{G}, -)$ and complete condensed internal group cohomology $\widehat{\underline{H}}_{\mathcal{R}}^{\bullet}(\mathcal{G}, -)$. In particular, there exists group cohomology and complete cohomology for any $T1$ topological group $G$ defined over its condensate $\underline{G}$.
\end{thm}

\begin{rem}\label{rem:differencetoguoliang}
Because both $\widehat{\underline{\mathrm{Ext}}}_{\mathrm{Cond}(Mod(\mathcal{R}))}^{\bullet}(A, -)$ and $\widehat{\underline{H}}_{\mathcal{R}}^{\bullet}(\mathcal{G}, -)$ have as codomain category $\mathrm{Cond}(\mathbf{Ab})$, they are an example of completed Ext-functors that are not covered under the previous frameworks such as the one by A.\! Beligiannis and I.\! Reiten in~\cite{bel07}, the one by S.\! Guo and L.\! Liang's in~\cite{guo23} and the one by J.\! Hu et al.\! in~\cite{hu21}.
\end{rem}

We conclude this paper by presenting an instance where the complete cohomology of a profinite group $G$ is isomorphic to the complete cohomology of its condensate $\underline{G}$. For this purpose, we introduce a special class of condensed modules. Namely, in the same manner as a completed group ring of a profinite group over a profinite ring can be seen as a completion of an ordinary group ring~\cite[Proposition~7.1.2]{wil98}, one can introduce a notion of completion for condensed modules over a condensed ring $\mathcal{R}$. We call a condensed ring admitting such a form of completion an analytic ring and a condensed module satisfying this completion a solid module. More formally, let $\mathbf{ProjPro}$ denote the full subcategory of $\mathbf{Pro}$ consisting of projective objects. Then an analytic ring $\mathcal{A}$ comes with a particular functor
\[ c_{\mathcal{A}}: \mathbf{ProjPro} \rightarrow \mathrm{Cond}(Mod(\mathcal{A})), \; S \mapsto \mathcal{A}[S]_{\blacksquare}  \, . \]
and a condensed $\mathcal{A}$-module homomorphism $C_{\mathcal{A}}(S): \mathcal{A}[\underline{S}] \rightarrow \mathcal{A}[\underline{S}]_{\blacksquare}$ for every projective profinite space $S$~\cite[p.~44]{sch19}. One can think of $c_{\mathcal{A}}$ as assigning to every projective profinite space $S$ a canonical choice of free solid $\mathcal{A}$-module $\mathcal{A}[S]_{\blacksquare}$. Let $\mathrm{Solid}(\mathcal{A})$ denote full subcategory of $\mathrm{Cond}(Mod(\mathcal{A}))$ for all whose objects $M$ the homomorphism
\[ \mathrm{Hom}_{\mathrm{Cond}(Mod(\mathcal{A}))}(C_{\mathcal{A}}(S), M): \mathrm{Hom}_{\mathrm{Cond}(Mod(\mathcal{A}))}(\mathcal{A}[S]_{\blacksquare}, M) \rightarrow M(S) \]
is an isomorphism for every $S \in \mathbf{ProjPro}$. According to~\cite[Proposition~7.5]{sch19}, $\mathrm{Solid}(\mathcal{A})$ is an abelian category stable under limits, colimits and extensionsin which the objects $\mathcal{A}[S]_{\blacksquare}$ are compact projective generators. In particular, we can define the Ext functors
\[ \mathrm{Ext}_{\mathrm{Solid}(\mathcal{A})}^{\bullet}(M, -): \mathrm{Solid}(\mathcal{A}) \rightarrow \mathbf{Ab} \]
as the derived functors of $\mathrm{Hom}_{\mathrm{Solid}(\mathcal{A})}(M, -)$. Moreover, there exists an extension of the morphisms $C_{\mathcal{A}}(S): \mathcal{A}[\underline{S}] \rightarrow \mathcal{A}[\underline{S}]_{\blacksquare}$ to a functor $\mathrm{Cond}(Mod(\mathcal{A})) \rightarrow \mathrm{Solid}(\mathcal{A})$ according to~\cite[Proposition~7.5]{sch19}. We call this functor solidification and need it to establish the isomorphism from the complete cohomology of a profinite group to the complete cohomology of its condensate.

\begin{exl}\label{exl:analyticrings}
Let $S = \varprojlim_{i \in I} S_i$ be a profinite space and define the condensed $\underline{\mathbb{Z}}$-module $\mathbb{Z}[S]^{\blacksquare} := \varprojlim_{i \in I} \underline{\mathbb{Z}}[\underline{S_i}]$. Then the condensed ring $\underline{\mathbb{Z}}$ together with the assignment $S \mapsto \mathbb{Z}[S]^{\blacksquare}$ forms an analytic ring~\cite[p.~46]{sch19}. Denote by $(-)^{\blacksquare}: \mathrm{Cond}(\mathbf{Ab}) \rightarrow \mathrm{Solid}(\underline{\mathbb{Z}})$ the associated solidification functor. If $\mathcal{R}$ is a condensed ring, then $(\mathcal{R}, S \mapsto (\mathcal{R}[\underline{S}])^{\blacksquare})$ is an analytic ring~\cite[Proposition~3.6]{tan24}.
\end{exl}

\begin{lem}\label{lem:solidandctscohom}
Let $G$ be a profinite group, $L$ be profinite ring and denote by $L\llbracket G \rrbracket$ the completed group ring. Endow the condensed ring $\underline{L\llbracket G \rrbracket}$ with an analytic ring structure as in Example~\ref{exl:analyticrings}. As the $\underline{L\llbracket G \rrbracket}$-module $\underline{L}$ is a solid by~\cite[pp.~14--15]{tan24}, one defines solid group cohomology $H_{\mathrm{Solid}(L)}^{\bullet}(\underline{G}, -)$ as the Ext-functor $\mathrm{Ext}_{\mathrm{Solid}(\underline{L\llbracket G \rrbracket})}^{\bullet}(\underline{L}, -)$. Then complete solid cohomology $\widehat{H}_{\mathrm{Solid}(L)}^{\bullet}(\underline{G}, -)$ is isomorphic as a cohomological functor to complete continuous cohomology $\widehat{H}_L^{\bullet}(G, -)$.
\end{lem}

\begin{proof}
Let $PMod(L\llbracket G \rrbracket)$ denote the category of profinite $L\llbracket G \rrbracket$-modules. Then the result follows from~\cite[Theorem~3.14]{tan24} because taking condensates constitutes an exact functor $PMod(L\llbracket G \rrbracket) \rightarrow \mathrm{Solid}(\underline{R\llbracket G \rrbracket})$ that preserves projectives.
\end{proof}

\section*{Acknowledgements}

A special thanks goes to Peter H. Kropholler for sharing his expertise with me. He has discussed examples of complete cohomology groups of discrete groups with me and pointed out his joint paper with Jonathan Cornick ``On Complete Resolutions''. In particular, he introduced me to condensed mathematics. I would also like to express my gratitude to Emma Brink for her great help in condensed mathematics. \\
%I am indebted to an anonymous referee whose invaluable feedback has substantially improved the paper. \\

I would like to acknowledge Alejandro Adem, Jon F.\! Carlson, Joel Friedman, Kalle Karu, Zinovy Reichstein and Ben Williams who provided invaluable suggestions as part of assessing my PhD thesis. Moreover, I am thankful to Andrew Fisher for pointing out the paper ``Complete Homology over Associative Rings'' by Olgur Celikbas et al. Lastly, I am grateful to Ged Corob Cook for insightful discussions.

%next line adds the Bibliography to the contents page
\addcontentsline{toc}{section}{References}
%change bibliography name to references
\renewcommand{\bibname}{References}
\bibliographystyle{plain}  %use the plain bibliography style
\bibliography{refs}        %use a bibtex bibliography file refs.bib

\end{document}
